\documentclass[11pt]{article}

\usepackage[margin=1in]{geometry}
\usepackage[T1]{fontenc}
\usepackage{lmodern}
\usepackage{amsmath,amssymb,amsthm}
\usepackage{mathtools}
\usepackage{enumitem}
\usepackage{booktabs}
\usepackage{microtype}
\usepackage[hidelinks]{hyperref}

\newcommand{\R}{\mathbb{R}}
\newcommand{\one}{\mathbf{1}}
\newcommand{\M}{\mathcal{M}}
\newcommand{\lexle}{\le_{\mathrm{lex}}}
\newcommand{\lexlt}{<_{\mathrm{lex}}}
\newcommand{\VB}{\mathrm{VB}}
\newcommand{\VBB}{\mathcal{V}_{\mathrm B}}
\newcommand{\ub}{\bar u}
\DeclareMathOperator{\conv}{conv}
\DeclareMathOperator{\cone}{cone}
\DeclareMathOperator{\relint}{relint}
\DeclareMathOperator{\rbd}{rbd}
\DeclareMathOperator{\aff}{aff}
\DeclareMathOperator{\supp}{supp}
\DeclareMathOperator{\ext}{ext}
\DeclareMathOperator{\cl}{cl}
\DeclareMathOperator{\inter}{int}
\DeclareMathOperator{\spn}{span}

\theoremstyle{plain}
\newtheorem{theorem}{Theorem}[section]
\newtheorem{proposition}[theorem]{Proposition}
\newtheorem{lemma}[theorem]{Lemma}
\newtheorem{corollary}[theorem]{Corollary}
\theoremstyle{definition}
\newtheorem{definition}[theorem]{Definition}
\newtheorem{example}[theorem]{Example}
\theoremstyle{remark}
\newtheorem{remark}[theorem]{Remark}

\setlist{itemsep=2pt,topsep=4pt}

\title{Covering coherent sets of desirable gambles}
\author{Jason Konek \and Catrin Campbell-Moore}
\date{Draft, 25 October 2026}

\begin{document}
\maketitle

\begin{abstract}
Fix finitely many worlds. A model is a set of gambles: the gambles it accepts. A ranking
assigns to each world a reflexive relation on models; the relations need not be transitive
or complete, and are required only to satisfy a weak acyclicity condition. A coherent
model is \emph{covered} by a ranking if the ranking is strictly guided by the truth and
compares models only where they differ, every undominated model is coherent, every
incoherent model is dominated by a coherent undominated one, ties with coherent
undominated models preserve that status, the target is undominated, and the domain leaves
the target free to vary near its boundary. We prove four covering theorems. First,
veto--Brier rankings, on open convex one-sided cones containing the positive orthant,
cover exactly the credal sets in the open simplex all of whose extreme points but at most
one are uniformly visible from the vertices of the simplex; this class contains many
standard imprecise-probability models. Second, single-level orthant-diagonal rankings
cover every orthant-planar credal set, on the full domain and on one-sided domains; these
include normalised boxes and density-ratio classes, and in three worlds they are exactly
the polygons each of whose edge lines passes through a vertex of the simplex. Third,
veto--box rankings, which follow the vetoes by a box chart instead of the Brier score,
allow one whole normalised box, possibly degenerate, to be exempt from visibility; in
three worlds this covers triangles that lie in neither of the first two classes. Fourth,
a composition of veto--Brier rankings with orthant-diagonal maps covers new targets in four
or more worlds and new targets with an exempt point on the boundary of the simplex; in
three worlds it yields no new target of composed form inside the open simplex. We also
show that strict guidance and global coherence alone force every domain for an interior
target, in three or more worlds, to be essentially one-sided or to contain every gamble
that is not a sure loss. Finally, we collect exclusions, each scoped to a class of
rankings: we prove those for veto--Brier rankings, for vetoes followed by a Brier-type
error, and for composed rankings, and we record others without proof. We list open
problems. Whether every coherent model is covered by some ranking remains open.
\end{abstract}

\tableofcontents

\section{Introduction}\label{sec:intro}

A \emph{model}, in this paper, is a set of gambles on finitely many worlds: the gambles it
accepts. A \emph{ranking} assigns to each world a relation on models, read as an ordering
by how well a model's acceptances fare if that world is actual. A dominance argument for
coherence needs a ranking under which every undominated model is coherent, every incoherent
model is dominated by a coherent undominated model, and a given coherent model, the
\emph{target}, is undominated. The ranking should also meet basic constraints. It should
reward correct acceptances and penalise wrong ones (strict guidance); it should compare
models only through the gambles on which they differ (independence); and the class of
models it evaluates should leave the target free to vary near its boundary (local
freedom). When all of this holds we say that the ranking \emph{covers} the target
(Definition~\ref{def:cover}).

We do not assume that the relations are transitive or complete. We assume only that they
are reflexive, that almost everywhere equal models are tied, and that they satisfy the
condition (TC) of Section~\ref{sec:framework}: if $A$ is at least as good as $B$ and $B$ at
least as good as $C$, then $C$ is not strictly better than $A$. All the rankings we
construct are generated by integrated scores, and are complete and transitive. They also
achieve a stronger dominance property than coverage requires: every model is weakly
dominated by a coherent undominated model.

The paper has four positive theorems.
\begin{enumerate}[label=(\arabic*)]
  \item \emph{Veto--Brier coverage} (Section~\ref{sec:vb}). At world $i$ a veto--Brier
  ranking first charges wrong acceptances that fall below zero under a prescribed set of
  probabilities, and then minimises a Brier-type error. It covers the coherent model $D_K$
  whenever all extreme points of $K$ except at most one are uniformly visible from the
  vertices of the simplex (Theorems~\ref{thm:cover}, \ref{thm:exempt} and~\ref{thm:vis}).
  The undominated family is identified exactly (Theorem~\ref{thm:family}). For targets in
  the open simplex this visibility condition is also necessary for this construction
  (Theorem~\ref{thm:reach}).
  \item \emph{Orthant-planar coverage} (Section~\ref{sec:op}). If, on every orthant
  whose gambles have coordinates of both signs, the coherent cone $D_K$ is cut out by a single
  hyperplane, then the single-level ranking with integrand $c(h)=d_\sigma\odot h$ on each
  orthant $\sigma$ covers $K$, both on the full domain and on one-sided domains
  (Theorem~\ref{thm:opcover}). This class contains every normalised box, and so every
  density-ratio class (Proposition~\ref{prop:box}); on the full domain it contains every
  precise model. In three worlds it is exactly the class of polygons each of whose edge
  lines passes through a vertex of the simplex (Theorem~\ref{thm:cevian}), and it contains
  triangles none of whose vertices is visible (Example~\ref{ex:cyclic}).
  \item \emph{Veto--box coverage} (Section~\ref{sec:vbox}). Replacing the Brier score by a
  box chart $c_j(h)=\min(l_jh_j,u_jh_j)$ after the vetoes lets one normalised box,
  possibly degenerate, be exempt from visibility (Theorem~\ref{thm:BX}). Since a segment on
  a line through a vertex of the simplex is a degenerate box, in three worlds this covers
  every interior triangle with exactly two non-visible vertices joined by such an edge, and among
  interior triangles the veto--box targets are exactly these and the veto--Brier targets
  (Corollary~\ref{cor:BX3}). An explicit example is the triangle with vertices
  $\frac15(1,2,2)$, $\frac15(2,1,2)$ and $\frac15(3,1,1)$ (Example~\ref{ex:T}), which is
  neither a veto--Brier target nor orthant-planar.
  \item \emph{Composition} (Section~\ref{sec:comp}). Running a veto--Brier ranking on the
  image of an orthant-diagonal map covers the composed target $L\star K_0=\conv\{\nu(q\odot
  m):q\in L,m\in K_0\}$ (Theorem~\ref{thm:comp}). In four or more worlds this covers
  targets that lie outside the veto--Brier and orthant-planar classes
  (Theorem~\ref{thm:F}, Example~\ref{ex:n4}); so does an exempt point on a face of the
  simplex, in any number of worlds (Theorem~\ref{thm:BE}, Example~\ref{ex:pentagon}). In
  three worlds, every target of the form $L\star K_0$ of Theorem~\ref{thm:comp} that lies
  inside the open simplex is already a veto--Brier target (Theorem~\ref{thm:n3}).
\end{enumerate}
Alongside these we prove that strict guidance alone forces a strong restriction on
domains: for a target meeting no edge of the simplex, if the models that are best at a
single world are coherent, then the domain is either essentially one-sided or contains
every gamble that is not a sure loss (Theorem~\ref{thm:dichotomy}). Every
full-dimensional credal set is a Hausdorff limit of covered polytopes
(Proposition~\ref{prop:dense}).

Section~\ref{sec:excl} collects exclusions. Each concerns a stated class of rankings, and
most concern the pari-mutuel models, of which the \emph{reversed triangle}
$\{m\in\Delta_3:m_i\le\frac25\}$ is the simplest. These are results about constructions.
None of them shows that a pari-mutuel model cannot be covered by a ranking meeting the
requirements of Section~\ref{sec:framework}; that question is open
(Section~\ref{sec:open}).

\section{Framework}\label{sec:framework}

\subsection{Gambles and credal sets}

Let $\Omega=\{1,\dots,n\}$ with $n\ge2$. A \emph{gamble} is a vector $h\in\R^n$; $h_i$ is
its payoff at world $i$. Write $e_1,\dots,e_n$ for the standard basis, $\one\coloneqq
(1,\dots,1)$, $\Delta\coloneqq\{m\in\R^n_+:\one\cdot m=1\}$ for the simplex, and
$\aff\Delta=\{m:\one\cdot m=1\}$. For $S\subseteq\Omega$ nonempty, $\Delta_S\coloneqq\{m\in
\Delta:\supp m\subseteq S\}$, where $\supp m\coloneqq\{i:m_i\neq0\}$; the \emph{edges} of
$\Delta$ are the sets $\Delta_{\{i,j\}}=[e_i,e_j]$. For $C\subseteq\aff\Delta$, $\relint C$
is the interior of $C$ relative to $\aff\Delta$ and $\rbd C\coloneqq C\setminus\relint C$.
A compact convex $K\subseteq\Delta$ is \emph{full-dimensional} if $\aff K=\aff\Delta$, that
is, if $\relint K\neq\emptyset$. A \emph{credal set} is a nonempty closed convex subset of
$\Delta$. For $S\subseteq\R^n$, $\cone S$ is the set of finite non-negative combinations of
points of $S$. For $w\in\R^n$ with $\one\cdot w>0$ put $\nu(w)\coloneqq w/(\one\cdot w)$.
We write $x\odot y\coloneqq(x_1y_1,\dots,x_ny_n)$, and $x\oslash y$ for the coordinatewise
quotient when $y$ has no zero coordinate.

For $S\subseteq\Delta$ define
\[
  D_S\coloneqq\{h\in\R^n: m\cdot h>0\text{ for all }m\in S\},\qquad
  C_S\coloneqq\{h\in\R^n: m\cdot h\ge0\text{ for all }m\in S\},
\]
with $D_\emptyset=C_\emptyset\coloneqq\R^n$. For $x\in S$ put $F^S_x\coloneqq\{h\in
C_S:x\cdot h=0\}$; when $S=K$ is fixed we write $F_x$.

\begin{lemma}\label{lem:D}
Let $S,S_1,S_2\subseteq\Delta$ be nonempty and compact.
\begin{enumerate}[label=(\alph*)]
  \item $D_S$ is an open convex cone containing every $h$ with all $h_i>0$, and
  $D_S\cap(-D_S)=\emptyset$. Moreover $D_S=D_{\conv S}$ and $D_{S_1}\cap D_{S_2}
  =D_{S_1\cup S_2}$.
  \item $\cl D_S=C_S$, and $\partial D_S=\{h:\min_{m\in S}m\cdot h=0\}$.
  \item If $S_1,S_2$ are convex and $D_{S_1}\subseteq D_{S_2}$, then $S_2\subseteq S_1$.
\end{enumerate}
\end{lemma}

\begin{proof}
(a) The function $h\mapsto\min_{m\in S}m\cdot h$ is continuous, since $|m\cdot h-m\cdot
h'|\le\max_j|h_j-h'_j|$ for $m\in\Delta$; so $D_S$ is open. Convexity, the cone property
and positivity on gambles with positive coordinates are immediate, and $m\cdot h>0$,
$m\cdot(-h)>0$ cannot both hold. A linear function has the same minimum over $S$ and
over $\conv S$, and this minimum is attained on $S$; hence $D_S=D_{\conv S}$. The last
identity is the definition.
(b) $D_S\subseteq C_S$ and $C_S$ is closed. If $h\in C_S$ then $h+t\one\in D_S$ for
$t>0$, because $m\cdot\one=1$. Hence $\cl D_S=C_S$, and $\partial D_S=C_S\setminus D_S$
since $D_S$ is open.
(c) Suppose $k\in S_2\setminus S_1$. Strict separation of $k$ from the compact convex set
$S_1$ gives $g\in\R^n$ and $c\in\R$ with $g\cdot m>c>g\cdot k$ for all $m\in S_1$. Put
$h\coloneqq g-c\one$. Since $\one\cdot m=1$ on $\Delta$, $m\cdot h>0$ on $S_1$ and
$k\cdot h<0$, so $h\in D_{S_1}\setminus D_{S_2}$.
\end{proof}

\begin{lemma}\label{lem:fulldim}
Let $K\subseteq\Delta$ be a full-dimensional compact convex set.
\begin{enumerate}[label=(\alph*)]
  \item Every affine function $\ell$ on $\aff\Delta$ has the form $\ell(m)=m\cdot h$ for
  exactly one $h\in\R^n$.
  \item $\relint K\subseteq\relint\Delta$, and $\ext K\subseteq\rbd K$.
  \item If $h\in C_K\setminus\{0\}$, then $m\cdot h>0$ for every $m\in\relint K$. In
  particular $C_K\cap(-C_K)=\{0\}$.
  \item If $x\in\rbd K$, then $F_x\neq\{0\}$.
\end{enumerate}
\end{lemma}

\begin{proof}
(a) If $\ell(m)=a\cdot m+c$ on $\aff\Delta$, then $\ell(m)=(a+c\one)\cdot m$ there. If
$m\cdot h=0$ for all $m\in\aff\Delta$, then $h_i=e_i\cdot h=0$ for every $i$.
(b) $\relint K$ is open in $\aff\Delta$ and contained in $\Delta$, hence in $\relint\Delta$.
A point $x\in\relint K$ is the midpoint of the segment $[x-td,x+td]\subseteq K$ for any
$d\neq0$ with $\one\cdot d=0$ and small $t>0$, so it is not extreme.
(c) If $m\mapsto m\cdot h$ is constant on $\aff\Delta$, then by (a) $h=c\one$, and $c>0$
because $h\in C_K\setminus\{0\}$. Otherwise suppose $x\in\relint K$ and $x\cdot h=0$.
Choose $d$ with $\one\cdot d=0$ and $d\cdot h\neq0$. For small $t>0$ both $x\pm td$ lie in
$K$, and one of $(x\pm td)\cdot h=\pm t\,d\cdot h$ is negative, contradicting $h\in C_K$.
If $h,-h\in C_K$, then $m\cdot h=0$ on $\relint K\ne\emptyset$, so $h=0$.
(d) By the supporting hyperplane theorem in $\aff\Delta$ \cite[Thm.~11.6]{Rockafellar}
there is a non-constant affine $\ell$ on $\aff\Delta$ with $\ell\ge\ell(x)$ on $K$. By (a)
$\ell-\ell(x)$ is $m\mapsto m\cdot h$ for some $h$, and $h\neq0$. Then $h\in F_x$.
\end{proof}

We use the following standard facts of convex geometry \cite{Rockafellar,Schrijver}: the
convex hull of a compact subset of $\R^n$ is compact; if $S$ is compact then $\ext\conv
S\subseteq S$; every compact convex set is the convex hull of its extreme points, and every
point of it is a convex combination of finitely many extreme points; if $x\in\relint K$ and
$k\in K$ then $(1-s)x+sk\in\relint K$ for $0\le s<1$; a full-dimensional polytope is the
intersection of the closed half-spaces defined by its facets; and a cone is polyhedral if
and only if it is finitely generated \cite[Cor.~7.1a]{Schrijver}. A \emph{face} of a convex
set $C$ is a convex $G\subseteq C$ such that whenever a point of $G$ is a convex
combination of points of $C$ with positive weights, all those points lie in $G$; extreme
points of a face are extreme in $C$. If a linear function is non-negative on $C$, its zero
set in $C$ is a face (an \emph{exposed} face).

For $q\in\Delta$ and $m\in\relint\Delta$ we have $q\cdot m>0$, so $\nu(q\odot m)$ is
defined. For fixed $m\in\relint\Delta$ the map $T_m(q)\coloneqq\nu(q\odot m)$ is a bijection
of $\Delta$ onto itself, with inverse $q\mapsto\nu(q\oslash m)$; it is continuous, maps
$\relint\Delta$ onto itself, fixes each vertex $e_j$, and maps segments onto segments,
because
\begin{equation}\label{eq:proj}
  \nu\bigl(((1-\lambda)q+\lambda q')\odot m\bigr)=(1-\mu)\,\nu(q\odot m)+\mu\,\nu(q'\odot m),
  \qquad \mu=\frac{\lambda\, q'\cdot m}{(1-\lambda)\,q\cdot m+\lambda\,q'\cdot m}.
\end{equation}
In particular $T_m$ is a homeomorphism of $\relint\Delta$, maps convex sets to convex sets,
maps each line through a vertex $e_j$ into a line through $e_j$, and maps extreme points of
a convex subset of $\Delta$ to extreme points of its image. The same holds with the roles
of $q$ and $m$ exchanged.

\subsection{Measure, models and domains}

\begin{itemize}
  \item[(M)] $\mu$ is a finite Borel measure on $\R^n$ with the same null sets as Lebesgue
  measure, and $\int|h|\,d\mu(h)<\infty$.
\end{itemize}
The standard Gaussian measure satisfies (M). The requirements of
Section~\ref{sec:framework} depend on $\mu$ only through its null sets; (M) is used only to
make the scores of the constructed rankings finite. So a ranking built with a measure
satisfying (M) is also a ranking in the sense of the framework for Lebesgue measure, or
for any measure with a positive density. Equalities and inclusions between sets of
gambles hold \emph{almost everywhere} (a.e.) if they hold up to a $\mu$-null set. Every
hyperplane is null and a nonempty open set is never null.

\begin{lemma}\label{lem:open}
Let $S\subseteq\Delta$ be nonempty and compact. If $O$ is open and $O\subseteq D_S$
a.e., then $O\subseteq D_S$. If $S_1,S_2\subseteq\Delta$ are nonempty, compact and
convex, and $D_{S_1}=D_{S_2}$ a.e., then $S_1=S_2$.
\end{lemma}

\begin{proof}
The open set $O\setminus\cl D_S$ is null, hence empty. So $O\subseteq\cl D_S$ and, being
open, $O\subseteq\inter\cl D_S$. For a convex set with nonempty interior the interior of
the closure is the interior \cite[Thm.~6.3]{Rockafellar}, so $O\subseteq D_S$. The second % citation number to be verified
claim follows from the first, applied to $O=D_{S_1}$ and $O=D_{S_2}$, and
Lemma~\ref{lem:D}(c).
\end{proof}

A \emph{domain} is a measurable set $U\subseteq\R^n$. The \emph{models} on $U$ are the
measurable sets $X\subseteq U$, identified when they agree a.e.; their class is $\M_U$. A
model $X$ \emph{accepts} $h$ if $h\in X$, and is otherwise silent about $h$. At world $i$,
accepting $h$ is \emph{correct} if $h_i>0$ and \emph{wrong} if $h_i<0$. The class $\M_U$ is
\emph{decomposable}: it is closed under splices $(X\cap S)\cup(Y\setminus S)$ for measurable
$S$, and under limits in measure. The \emph{full domain} is $U=\R^n$. Let $M$ be the set of
gambles with a positive and a negative coordinate. We call $U$ \emph{one-sided} if
$U\cap(-U)=\emptyset$, and \emph{essentially one-sided} if $U\cap(-U)\cap M$ is null.
Gambles outside $M$ are sure gains or sure losses (or zero).

A model $X$ is \emph{coherent} if $X=D_Q$ a.e.\ for some credal set $Q$. These are the
coherent sets of desirable gambles of the form $D_Q$, taken up to null sets. The
condition is global: a model on a domain $U$ is coherent only if it agrees a.e.\ with a
coherent cone in $\R^n$, not merely with the trace of one on $U$. By
Lemma~\ref{lem:open}, $Q$ is unique. If $U=D_{K'}$ for a nonempty compact convex $K'$, then
$X\subseteq U$ and Lemmas~\ref{lem:open} and~\ref{lem:D}(c) give $K'\subseteq Q$.

\subsection{Rankings and dominance}

A \emph{ranking} on $\M_U$ is a family $(\preceq_i)_{i\in\Omega}$ of binary relations on
$\M_U$. Lower is better: $X\preceq_iY$ says that $X$ does at least as well as $Y$ at world
$i$. We write $X\prec_iY$ if $X\preceq_iY$ and not $Y\preceq_iX$; $X$ and $Y$ are
\emph{tied} at $i$ if $X\preceq_iY$ and $Y\preceq_iX$, and \emph{incomparable} at $i$ if
neither holds. We assume throughout that each $\preceq_i$ is reflexive, that a.e.-equal
models are tied, and
\begin{itemize}
  \item[(TC)] if $A\preceq_iB$ and $B\preceq_iC$, then not $C\prec_iA$.
\end{itemize}
Every transitive relation satisfies (TC). We do \emph{not} assume transitivity or
completeness. For instance, (TC) permits $Y\preceq_iX$ and $X\prec_iM$ with $Y$ and $M$
incomparable at $i$.

\begin{definition}\label{def:dom}
Let $X,Y\in\M_U$. $Y$ \emph{weakly dominates} $X$ if $Y\preceq_iX$ for every $i$. $Y$
\emph{dominates} $X$ if $Y$ weakly dominates $X$ and $Y\prec_iX$ for some $i$. $Y$
\emph{strictly dominates} $X$ if $Y\prec_iX$ for every $i$. $X$ is \emph{undominated} if
no model dominates it. We write $\mathcal U$ for the undominated models, $\mathcal C$ for
the coherent models, and $\mathcal C^\ast\coloneqq\mathcal C\cap\mathcal U$.
\end{definition}

When a ranking is generated by real-valued losses $\ell(X,i)$, with $X\preceq_iY$ iff
$\ell(X,i)\le\ell(Y,i)$, one says that $Y$ \emph{uniformly dominates} $X$ if
$\ell(Y,i)+\varepsilon\le\ell(X,i)$ at every world for some $\varepsilon>0$. With finitely
many worlds this is strict dominance. The rankings below are lexicographic and uniform
dominance plays no role.

If $Y$ weakly dominates $X$ but does not dominate it, then for every $i$ we have
$Y\preceq_iX$ and not $Y\prec_iX$, hence $X\preceq_iY$: the two models are tied at every
world. This uses only the definition of $\prec_i$, not completeness.

\subsection{Requirements}

\begin{itemize}
  \item[(G)] \emph{Strict guidance.} If $X\setminus Y\subseteq\{h_i>0\}$ a.e.,
  $Y\setminus X\subseteq\{h_i<0\}$ a.e., and $X\triangle Y$ is not null, then
  $X\prec_iY$.
  \item[(IB)] \emph{Independence.} Let $S$ be measurable and $X,Y,X',Y'\in\M_U$ with
  $X\setminus S=Y\setminus S$, $X'\setminus S=Y'\setminus S$, $X\cap S=X'\cap S$ and
  $Y\cap S=Y'\cap S$, all a.e. Then $X\preceq_iY$ if and only if $X'\preceq_iY'$.
  \item[(GC)] \emph{Global coherence.} Every undominated model is coherent.
  \item[(C)] \emph{Direct closure under ties.} A model that is tied at every world with a
  member of $\mathcal C^\ast$ belongs to $\mathcal C^\ast$.
  \item[(ID)] \emph{Incoherent-model dominance.} Every incoherent model is dominated by a
  member of $\mathcal C^\ast$.
  \item[(LF)] \emph{Local freedom at $D_K$.} There is an open set $N\supseteq\partial
  D_K\setminus\{0\}$ such that every measurable $X$ with $X\triangle D_K\subseteq N$ a.e.\
  belongs to $\M_U$, that is, satisfies $X\subseteq U$ a.e.
\end{itemize}
By (IB), the comparison $X\preceq_iY$ depends only on the pair $(X\cap S,Y\cap S)$ for any
$S\supseteq X\triangle Y$.

\begin{definition}\label{def:cover}
Let $K$ be a credal set and $U$ a domain. A ranking on $\M_U$ \emph{covers} $K$ if:
\begin{enumerate}[label=(\roman*)]
  \item it satisfies (G) and (IB);
  \item it satisfies global coherence (GC);
  \item it satisfies direct closure under ties (C);
  \item $D_K\in\M_U$ and $D_K$ is undominated;
  \item it satisfies incoherent-model dominance (ID);
  \item (LF) holds at $D_K$.
\end{enumerate}
The domain $\M_U$ is decomposable by construction.
\end{definition}

Several conditions are \emph{not} required: transitivity, completeness, numerical
representation, continuity, homogeneity, uniqueness of optima, and the following stronger
dominance objective, which our constructions nevertheless achieve.
\begin{itemize}
  \item[(W)] Every model is weakly dominated by a member of $\mathcal C^\ast$.
\end{itemize}

\begin{lemma}\label{lem:dominance}
Consider a ranking on $\M_U$.
\begin{enumerate}[label=(\alph*)]
  \item Suppose every model is weakly dominated by an undominated model, and (GC) holds.
  Then (W) holds and $\mathcal U=\mathcal C^\ast$. If in addition (C) holds, then (ID)
  holds.
  \item If every $\preceq_i$ is transitive, then (GC) implies (C).
\end{enumerate}
\end{lemma}

\begin{proof}
(a) (W) and $\mathcal U=\mathcal C^\ast$ are immediate. Let $X$ be incoherent and let
$Y\in\mathcal C^\ast$ weakly dominate it. If $Y$ did not dominate $X$, then $X$ and $Y$
would be tied at every world (see the remark after Definition~\ref{def:dom}), and (C)
would make $X$ coherent. So $Y$ dominates $X$.

(b) Let $X$ be tied at every world with $Y\in\mathcal C^\ast$, and suppose $Z$ dominates
$X$. For every $i$, $Z\preceq_iX\preceq_iY$ gives $Z\preceq_iY$ by transitivity. For some
$i$, $Z\prec_iX$; if $Y\preceq_iZ$, then $X\preceq_iY\preceq_iZ$ would give $X\preceq_iZ$,
contradicting $Z\prec_iX$. So $Z\prec_iY$ and $Z$ dominates $Y$, which is impossible.
Hence $X\in\mathcal U$, and $X\in\mathcal C$ by (GC).
\end{proof}

\begin{remark}\label{rem:TC}
The derivation in (b) uses transitivity. For abstract relations satisfying (TC) the
implication fails. Consider three abstract objects $X,Y,Z$ and one world, with $X$ and $Y$
tied, $Z\prec X$, and $Z$ and $Y$ incomparable; the relation consists of these comparisons
and reflexivity. It satisfies (TC): the only strict comparison is $Z\prec X$, so a
violation would need a chain $X\preceq B\preceq Z$, but the only $B$ with $B\preceq Z$ is
$Z$, and $X\preceq Z$ fails. Here $Y$ is undominated, $X$ is tied with $Y$, and $Z$
dominates $X$. This configuration is not a ranking on a class $\M_U$ and need not satisfy
(G) or (IB). We do not know whether (G), (IB), (TC) and (GC) together imply (C); this is
why (C) is listed as a requirement.
\end{remark}

\subsection{What strict guidance alone forces}\label{sec:guidance}

The following results use only (G) at single worlds and the structure of $\M_U$. They hold
for every ranking in the sense of this section; no transitivity, completeness or (IB) is
used.

\begin{lemma}\label{lem:worldbest}
Let $U$ be a domain and $i\in\Omega$. Under (G) at world $i$, the model $B_i\coloneqq
U\cap\{h_i>0\}$ is undominated.
\end{lemma}

\begin{proof}
Let $Y\in\M_U$ with $Y\neq B_i$ a.e. Then $B_i\setminus Y\subseteq\{h_i>0\}$ and
$Y\setminus B_i\subseteq\{h_i\le0\}$, which is a.e.\ $\{h_i<0\}$. By (G), $B_i\prec_iY$,
so not $Y\preceq_iB_i$, and $Y$ does not weakly dominate $B_i$. A model a.e.\ equal to
$B_i$ is $B_i$.
\end{proof}

\begin{theorem}\label{thm:dichotomy}
Let $n\ge3$, let $K$ be a credal set that meets no edge of $\Delta$, and let $U$ be a domain
with $D_K\subseteq U$ a.e. Suppose some ranking on $\M_U$ satisfies (G), and that the $n$
models $B_i=U\cap\{h_i>0\}$ are coherent (for instance, because (GC) holds). Write
$B_i=D_{Q_i}$ a.e. Then $e_i\in Q_i\subseteq\conv(K\cup\{e_i\})$ and $Q_j\subseteq
\conv(Q_i\cup\{e_j\})$ for all $i,j$, and exactly one of the following holds.
\begin{enumerate}[label=(\alph*)]
  \item Every two of the sets $Q_i$ intersect, and $U$ is essentially one-sided.
  \item $Q_i=\{e_i\}$ for every $i$, and $U\supseteq\{h:h_l>0\text{ for some }l\}$ a.e.
\end{enumerate}
\end{theorem}

\begin{proof}
\emph{Step 1.} $D_{Q_i}$ is open and a.e.\ contained in $\{h_i>0\}=D_{\{e_i\}}$, so
Lemma~\ref{lem:open} gives $D_{Q_i}\subseteq D_{\{e_i\}}$, and Lemma~\ref{lem:D}(c) gives
$e_i\in Q_i$.

\emph{Step 2.} $D_{Q_i}\subseteq U$ a.e., so $D_{\conv(Q_i\cup\{e_j\})}=D_{Q_i}\cap\{h_j>0\}
\subseteq U\cap\{h_j>0\}=D_{Q_j}$ a.e. Lemma~\ref{lem:open} and Lemma~\ref{lem:D}(c) give
$Q_j\subseteq\conv(Q_i\cup\{e_j\})$. Likewise $D_{\conv(K\cup\{e_i\})}=D_K\cap\{h_i>0\}
\subseteq D_{Q_i}$ a.e.\ gives $Q_i\subseteq\conv(K\cup\{e_i\})$.

\emph{Step 3: disjoint pairs lie in an edge.} Let $i\ne j$ and $f(m)\coloneqq1-m_i-m_j$,
which is affine, non-negative on $\Delta$ and zero at $e_i,e_j$. Put $s_i\coloneqq
\max_{Q_i}f$ and $s_j\coloneqq\max_{Q_j}f$. Every point $(1-\lambda)e_i+\lambda q$ of
$\conv(Q_j\cup\{e_i\})$ ($q\in Q_j$) has $f=\lambda f(q)\le s_j$, so Step 2 gives $s_i\le
s_j$; by symmetry $s_i=s_j\eqqcolon s$. If $s>0$, take $a\in Q_i$ with $f(a)=s$ and write
$a=(1-\lambda)e_i+\lambda q$ with $q\in Q_j$; then $s=\lambda f(q)\le\lambda s$ forces
$\lambda=1$ and $a=q\in Q_i\cap Q_j$. So if $Q_i\cap Q_j=\emptyset$ then $s=0$, that is,
$Q_i,Q_j\subseteq[e_i,e_j]$.

\emph{Step 4: an edge forces precise vertices.} A point $(1-t)k+te_i$ with $k\in K$ and
$t<1$ lying in $[e_i,e_j]$ would have $k_l=0$ for all $l\notin\{i,j\}$, so $k\in[e_i,e_j]$,
which is excluded. Hence $\conv(K\cup\{e_i\})\cap[e_i,e_j]=\{e_i\}$, and if $Q_i\cap Q_j=
\emptyset$ then Step 2 gives $Q_i=\{e_i\}$, and likewise $Q_j=\{e_j\}$. For any other $l$,
Step 2 gives $Q_l\subseteq\conv(Q_i\cup\{e_l\})=[e_i,e_l]$ and $Q_l\subseteq\conv(K\cup
\{e_l\})$, and the same argument with the edge $[e_i,e_l]$ gives $Q_l=\{e_l\}$.

\emph{Step 5: sidedness.} Suppose $U\cap(-U)\cap M$ is not null. Since $M=\bigcup_{i\neq
j}\{h_i>0>h_j\}$, for some $i\ne j$ the set $E\coloneqq U\cap(-U)\cap\{h_i>0>h_j\}$ is not
null. Then $E\subseteq B_i=D_{Q_i}$ a.e.\ and $-E\subseteq B_j=D_{Q_j}$ a.e. If
$q\in Q_i\cap Q_j$, almost every $h\in E$ would satisfy $q\cdot h>0$ and $q\cdot(-h)>0$,
which is impossible. So some pair $Q_i,Q_j$ is disjoint, and Step 4 gives (b). Conversely,
in case (b), $U\cap\{h_l>0\}=D_{\{e_l\}}=\{h_l>0\}$ a.e.\ for every $l$, and two sets
$\{e_i\},\{e_j\}$ with $i\ne j$ are disjoint, so (a) fails. In summary: if all pairs
intersect, Step~5 shows that $U$ is essentially one-sided; if some pair is disjoint,
Step~4 gives (b); and (a) and (b) exclude each other.
\end{proof}

\begin{proposition}\label{prop:lowdim}
Let $K$ be a credal set that is not full-dimensional, and let $U$ be a domain with
$D_K\subseteq U$ a.e.\ and (LF) at $D_K$. Then $U$ is not essentially one-sided. If
moreover $n\ge3$, $K$ meets no edge of $\Delta$, and some ranking on $\M_U$ satisfies (G)
and (GC), then $U\supseteq\{h:h_l>0\text{ for some }l\}$ a.e.
\end{proposition}

\begin{proof}
The linear span of $K$ has dimension less than $n$, so some $h\neq0$ has $m\cdot h=0$ for
all $m\in K$. Then $\pm h\in C_K\setminus D_K=\partial D_K$. Let $N$ be as in (LF). The
model $D_K\cup N$ belongs to $\M_U$, so $N\subseteq U$ a.e. Choose $r>0$ with $B(\pm
h,r)\subseteq N$. Then $B(h,r)\subseteq U$ a.e.\ and $B(h,r)=-B(-h,r)\subseteq-U$ a.e.
Since $m\cdot h=0$ for some $m\in\Delta$ and $h\ne0$, the vector $h$ has a coordinate
$h_a\ge0$ and a coordinate $h_b\le0$ with $a\ne b$ (if all coordinates have one sign, some
coordinate is zero and another is not). Then $h+t(e_a-e_b)\in M$ for $t>0$, so every ball
around $h$ meets the open set $M$ in a nonempty open set. Hence $U\cap(-U)\cap M$ is not null. The second claim
then follows from Theorem~\ref{thm:dichotomy}.
\end{proof}

So targets that are not full-dimensional, including precise models $\{p\}$, require
two-sided domains, and under (GC) a precise model in $\relint\Delta$ with $n\ge3$ requires
essentially the full domain. Such targets are covered there (Corollary~\ref{cor:precise}).

\section{Two-level integral rankings}\label{sec:route}

All rankings constructed in this paper belong to the following class. Throughout this
section $U$ is a domain.

\begin{definition}\label{def:twolevel}
A \emph{two-level integral ranking} on $\M_U$ is given by $\mu$-integrable functions
$\varphi_i:U\to[0,\infty)$ and $\psi_i:U\to\R$, $i\in\Omega$, through
\[
  a_i(X)\coloneqq\int_X\varphi_i\,d\mu,\qquad b_i(X)\coloneqq\int_X\psi_i\,d\mu,\qquad
  X\preceq_iY\iff(a_i(X),b_i(X))\lexle(a_i(Y),b_i(Y)),
\]
where $(x_1,x_2)\lexle(y_1,y_2)$ means $x_1<y_1$, or $x_1=y_1$ and $x_2\le y_2$. It is
\emph{single-level} if every $\varphi_i=0$. It satisfies the \emph{sign condition} (SC) if,
for every $i$, $\varphi_i=0$ a.e.\ on $U\cap\{h_i>0\}$, $\psi_i<0$ a.e.\ on $U\cap
\{h_i>0\}$, and $\psi_i>0$ a.e.\ on $U\cap\{h_i<0\}$.
\end{definition}

\begin{lemma}\label{lem:axioms}
Every two-level integral ranking consists of total preorders under which a.e.-equal models
are tied; in particular it satisfies reflexivity and (TC). It satisfies (IB). Under (SC) it
satisfies (G).
\end{lemma}

\begin{proof}
$\lexle$ is a total order on $\R^2$, and the scores depend only on the a.e.-class of $X$.
(IB): with $X,Y,X',Y',S$ as in (IB), $a_i(X)-a_i(Y)=\int_S(1_X-1_Y)\varphi_i\,d\mu=
a_i(X')-a_i(Y')$, and likewise for $b_i$; and $p\lexle q$ iff $0\lexle q-p$. (G): let
$X,Y$ be as in (G). Since $\varphi_i=0$ a.e.\ on $X\setminus Y$, $a_i(X)-a_i(Y)=
-\int_{Y\setminus X}\varphi_i\,d\mu\le0$. Also
\[
  b_i(X)-b_i(Y)=\int_{X\setminus Y}\psi_i\,d\mu-\int_{Y\setminus X}\psi_i\,d\mu<0,
\]
because $\psi_i<0$ a.e.\ on $X\setminus Y$, $\psi_i>0$ a.e.\ on $Y\setminus X$, and
$X\triangle Y$ is not null. Hence $(a_i(X),b_i(X))\lexlt(a_i(Y),b_i(Y))$.
\end{proof}

Write $z(X)\coloneqq(a_1(X),\dots,a_n(X),b_1(X),\dots,b_n(X))\in\R^{2n}$ and
$Z\coloneqq\{z(X):X\in\M_U\}$. For $y,y'\in\R^{2n}$ write $y\trianglelefteq y'$ if
$(y_{a_i},y_{b_i})\lexle(y'_{a_i},y'_{b_i})$ for every $i$, where $y_{a_i}$ and $y_{b_i}$
are the coordinates in the positions of $a_i$ and $b_i$. This is a partial order, because
$\lexle$ is antisymmetric. By definition, $Y$ weakly dominates $X$ iff $z(Y)\trianglelefteq
z(X)$, and $Y$ dominates $X$ iff in addition $z(Y)\neq z(X)$.

\begin{lemma}\label{lem:lyap}
$Z$ is a compact convex subset of $\R^{2n}$, and $0=z(\emptyset)\in Z$.
\end{lemma}

\begin{proof}
$z(X)=\int_Xf\,d\mu$ with $f\coloneqq(\varphi_1,\dots,\varphi_n,\psi_1,\dots,\psi_n)$, which
is $\mu$-integrable on $U$. Each coordinate of the vector measure $X\mapsto\int_Xf\,d\mu$
on the measurable subsets of $U$ is absolutely continuous with respect to $\mu$, which has
no atoms. By Lyapunov's convexity theorem \cite{Lyapunov,Lindenstrauss} its range is
compact and convex. Since every measurable subset of $U$ is a model, this range is $Z$.
\end{proof}

\begin{lemma}\label{lem:zorn}
Every model is weakly dominated by an undominated model.
\end{lemma}

\begin{proof}
Fix $X_0$ and put $D\coloneqq\{y\in Z:y\trianglelefteq z(X_0)\}$. We show that every chain
$\mathcal K\subseteq D$ has a lower bound in $D$. If $\mathcal K=\emptyset$, take
$z(X_0)$. Otherwise, for $\gamma\in\mathcal K$ let $T_\gamma\coloneqq\{\gamma'\in
\mathcal K:\gamma'\trianglelefteq\gamma\}$. Since $\mathcal K$ is a chain, the intersection
of finitely many sets $T_\gamma$ contains the least of their indices, so the closed sets
$\cl T_\gamma$ have the finite intersection property. As $Z$ is compact, some $y^\infty\in
Z$ lies in every $\cl T_\gamma$.

Fix $\gamma\in\mathcal K$ and a world $i$. Every $\gamma'\in T_{\gamma''}$ has
$\gamma'_{a_i}\le\gamma''_{a_i}$, so $y^\infty_{a_i}\le\gamma''_{a_i}$ for every
$\gamma''\in\mathcal K$. If $y^\infty_{a_i}<\gamma_{a_i}$, then
$(y^\infty_{a_i},y^\infty_{b_i})\lexlt(\gamma_{a_i},\gamma_{b_i})$. Otherwise
$y^\infty_{a_i}=\gamma_{a_i}$, and for $\gamma'\in T_\gamma$ we get $\gamma_{a_i}=
y^\infty_{a_i}\le\gamma'_{a_i}\le\gamma_{a_i}$. So $a_i$ is constant on $T_\gamma$, and
the lexicographic order then gives $\gamma'_{b_i}\le\gamma_{b_i}$ on $T_\gamma$. Since
$y^\infty\in\cl T_\gamma$, $y^\infty_{b_i}\le\gamma_{b_i}$. In both cases
$(y^\infty_{a_i},y^\infty_{b_i})\lexle(\gamma_{a_i},\gamma_{b_i})$.

Hence $y^\infty\trianglelefteq\gamma\trianglelefteq z(X_0)$ for all $\gamma\in\mathcal K$,
so $y^\infty\in D$ is a lower bound. By Zorn's lemma $D$ has a minimal element $y^\ast$.
Let $X^\ast$ be a model with $z(X^\ast)=y^\ast$; it weakly dominates $X_0$. If $Y$
dominated $X^\ast$, then $z(Y)\trianglelefteq y^\ast$ and $z(Y)\neq y^\ast$, so $z(Y)\in
D$ would contradict minimality.
\end{proof}

For $W_0\subseteq\Omega$ put
\[
  N(W_0)\coloneqq U\cap\bigcap_{\omega\in W_0}\{\varphi_\omega=0\}.
\]
Since $\varphi_\omega\ge0$, a model $X$ has $a_\omega(X)=0$ iff $X\subseteq
\{\varphi_\omega=0\}$ a.e. A triple $(W_0,\beta,\alpha)$ is \emph{admissible} if
$W_0\subseteq\Omega$, $\beta\in\R^n_+\setminus\{0\}$ with $\supp\beta\subseteq W_0$, and
$\alpha\in\R_+^{\Omega\setminus W_0}$. For such a triple put
\[
  J(X)\coloneqq\sum_{\omega\in W_0}\beta_\omega b_\omega(X)
    +\sum_{\omega\notin W_0}\alpha_\omega a_\omega(X),\qquad
  g\coloneqq\sum_{\omega\notin W_0}\alpha_\omega\varphi_\omega+\sum_{\omega\in W_0}
  \beta_\omega\psi_\omega,
\]
and $X^\ast(W_0,\beta,\alpha)\coloneqq N(W_0)\cap\{g<0\}$. Then $J(X)=\int_Xg\,d\mu$
whenever $X\subseteq N(W_0)$ a.e. For a single-level ranking, $N(W_0)=U$ for every $W_0$,
and we may always take $W_0=\Omega$.

\begin{lemma}\label{lem:sep}
Let $X^\ast$ be undominated and $W_0\coloneqq\{\omega:a_\omega(X^\ast)=0\}$. Then $W_0\neq
\emptyset$, $X^\ast\subseteq N(W_0)$ a.e., and there are $\beta,\alpha$ with
$(W_0,\beta,\alpha)$ admissible such that $X^\ast$ minimises $J$ over the models $X\subseteq
N(W_0)$.
\end{lemma}

\begin{proof}
If $W_0=\emptyset$, every $a_\omega(X^\ast)>0=a_\omega(\emptyset)$, so $\emptyset$ strictly
dominates $X^\ast$. Hence $W_0\neq\emptyset$, and $X^\ast\subseteq N(W_0)$ a.e.\ by the
remark above.

Let $Z'\coloneqq\{y\in Z:y_{a_\omega}=0\text{ for }\omega\in W_0\}$, a convex set
containing $z(X)$ for every model $X\subseteq N(W_0)$, and $y^\ast\coloneqq z(X^\ast)\in
Z'$. Let $\mathrm{pr}$ be the projection of $\R^{2n}$ onto the coordinates $y_{b_\omega}$
($\omega\in W_0$) and $y_{a_\omega},y_{b_\omega}$ ($\omega\notin W_0$), and let
\[
  P\coloneqq\{d:d_{b_\omega}<0\text{ for }\omega\in W_0,\ d_{a_\omega}<0\text{ for }
  \omega\notin W_0\}
\]
in the target of $\mathrm{pr}$, a nonempty open convex cone. If some $y\in Z'$ had
$\mathrm{pr}(y-y^\ast)\in P$, the model realising $y$ would be strictly better than
$X^\ast$ at every world: at $\omega\in W_0$ its first level is equal (zero) and its second
level strictly smaller, and at $\omega\notin W_0$ its first level is strictly smaller. So
the convex set $E\coloneqq\mathrm{pr}(Z'-y^\ast)$ is disjoint from $P$.

By the separation theorem \cite[Thm.~11.3]{Rockafellar} there is a nonzero linear % citation number to be verified
functional $L$ with $\sup_PL\le\inf_EL$. Since $0\in E$ and $P$ is a cone with $0\in\cl P$,
$\sup_PL=0$; so $L\le0$ on $\cl P$ and $L\ge0$ on $E$. Write
\[
  L(d)=\sum_{\omega\in W_0}\beta_\omega d_{b_\omega}+\sum_{\omega\notin W_0}
  (\alpha_\omega d_{a_\omega}+\gamma_\omega d_{b_\omega}).
\]
Testing $L\le0$ on the vectors of $\cl P$ with one coordinate equal to $-1$ and the others
zero, and on $\pm$ the unit vectors in the coordinates $d_{b_\omega}$ ($\omega\notin W_0$),
which are unconstrained in $\cl P$, gives $\beta\ge0$, $\alpha\ge0$ and $\gamma=0$. Extend
$\beta$ by zero outside $W_0$. Then $L(\mathrm{pr}(y-y^\ast))=J(y)-J(y^\ast)\ge0$ for $y\in
Z'$, where $J$ is read as a function of the score vector. So $X^\ast$ minimises $J$ over the
models $X\subseteq N(W_0)$.

It remains to show $\beta\neq0$. Otherwise $\alpha\neq0$ and $J=\sum_{\omega\notin W_0}
\alpha_\omega a_\omega\ge0$ on $Z'$, with $J(0)=0$ and $0=z(\emptyset)\in Z'$. Hence
$J(y^\ast)=0$, so $\alpha_\omega a_\omega(X^\ast)=0$ for every $\omega\notin W_0$. But
$a_\omega(X^\ast)>0$ for $\omega\notin W_0$, so $\alpha=0$, a contradiction.
\end{proof}

We say that the ranking has \emph{null ties} if $N(W_0)\cap\{g=0\}$ is null for every
admissible triple.

\begin{lemma}\label{lem:unique}
Let $(W_0,\beta,\alpha)$ be admissible and suppose $N\coloneqq N(W_0)$ meets $\{g=0\}$ in a
null set. Then the models $X\subseteq N$ minimising $\int_Xg\,d\mu$ are exactly those equal
a.e.\ to $X^\ast(W_0,\beta,\alpha)$.
\end{lemma}

\begin{proof}
For $X\subseteq N$,
\[
  \int_Xg\,d\mu=\int_{X\cap\{g<0\}}g\,d\mu+\int_{X\cap\{g>0\}}g\,d\mu\ \ge\
  \int_{N\cap\{g<0\}}g\,d\mu,
\]
with equality iff $(N\cap\{g<0\})\setminus X$ and $X\cap\{g>0\}$ are null, that is, iff
$X=N\cap\{g<0\}$ a.e.
\end{proof}

\begin{theorem}\label{thm:route}
Let $(\varphi,\psi)$ be a two-level integral ranking on $\M_U$ satisfying (SC) and having
null ties.
\begin{enumerate}[label=(\alph*)]
  \item It consists of total preorders and satisfies (TC), (G) and (IB).
  \item Every model is weakly dominated by an undominated model.
  \item A model is undominated if and only if it equals $X^\ast(W_0,\beta,\alpha)$ a.e.\
  for some admissible triple.
  \item If a model is tied at every world with an undominated model $X$, it equals $X$
  a.e.
  \item If every model $X^\ast(W_0,\beta,\alpha)$ is coherent, then (GC), (C), (ID) and
  (W) hold, and $\mathcal U=\mathcal C^\ast$.
\end{enumerate}
\end{theorem}

\begin{proof}
(a) is Lemma~\ref{lem:axioms} and (b) is Lemma~\ref{lem:zorn}.

(c) Let $X^\ast$ be undominated. Lemma~\ref{lem:sep} gives an admissible triple such that
$X^\ast\subseteq N(W_0)$ minimises $\int_Xg\,d\mu$ over the models $X\subseteq N(W_0)$, and
Lemma~\ref{lem:unique} gives $X^\ast=X^\ast(W_0,\beta,\alpha)$ a.e. Conversely, let
$(W_0,\beta,\alpha)$ be admissible, $X^\ast\coloneqq X^\ast(W_0,\beta,\alpha)$, and let $Y$
weakly dominate $X^\ast$. For $\omega\in W_0$, $a_\omega(X^\ast)=0$, so $a_\omega(Y)\le0$;
hence $a_\omega(Y)=0$, and the lexicographic order then gives $b_\omega(Y)\le
b_\omega(X^\ast)$. For $\omega\notin W_0$, $a_\omega(Y)\le a_\omega(X^\ast)$. Thus
$Y\subseteq N(W_0)$ a.e.\ and, as the weights are non-negative, $J(Y)\le J(X^\ast)$. By
Lemma~\ref{lem:unique}, $Y=X^\ast$ a.e. So every weak dominator of $X^\ast$ equals it
a.e., and $X^\ast$ is undominated.

(d) Let $Y$ be tied with $X$ at every world. Then $z(Y)=z(X)$, because $\lexle$ is
antisymmetric. By (c), $X=X^\ast(W_0,\beta,\alpha)$ a.e.\ for an admissible triple. Then
$a_\omega(Y)=0$ for $\omega\in W_0$, so $Y\subseteq N(W_0)$ a.e., and $J(Y)=J(X)$. By
Lemma~\ref{lem:unique}, $Y=X$ a.e.

(e) (GC) follows from (c). (C) follows from (d): a model tied at every world with a member
of $\mathcal C^\ast$ is that member. Lemma~\ref{lem:dominance}(a) gives (W), (ID) and
$\mathcal U=\mathcal C^\ast$.
\end{proof}

So for a two-level integral ranking with (SC), null ties and coherent minimisers, coverage
of $K$ reduces to two further facts: $D_K$ is one of the models $X^\ast(W_0,\beta,\alpha)$
(or is otherwise shown to be undominated), and (LF) holds at $D_K$. Every ranking
constructed in Sections~\ref{sec:vb}--\ref{sec:comp} is a two-level integral ranking with
(SC), null ties and coherent minimisers. So each of them is generated by integrated
scores, is complete and transitive, and satisfies the stronger objective (W) as well as
the requirements of Definition~\ref{def:cover}.

\begin{lemma}\label{lem:fubini}
Let $F:\R^n\to\R$ be continuous, and suppose $t\mapsto F(h+t\one)$ is strictly monotone
for every $h$. Then $\{F=0\}$ is null.
\end{lemma}

\begin{proof}
$\{F=0\}$ is closed. Every $h$ is uniquely $y+t\one$ with $y\in\one^\perp$ and $t\in\R$;
this linear change of coordinates multiplies Lebesgue measure by a constant. For each $y$
the section $\{t:F(y+t\one)=0\}$ has at most one point. By Fubini's theorem $\{F=0\}$ is
Lebesgue-null, hence null by (M).
\end{proof}

\section{Veto--Brier rankings and the visibility class}\label{sec:vb}

\subsection{Veto--Brier rankings}

Throughout this section $K'\subseteq\Delta$ is nonempty, compact and convex, and
$U=D_{K'}$. Everything in this section holds for every measure satisfying (M).

\begin{definition}\label{def:vb}
A \emph{veto profile} over $K'$ is a tuple $B=(B_1,\dots,B_n)$ of compact convex sets
with $0\in B_i\subseteq\cone(K'\cup\{e_i\})$. Put
\[
  \varphi_i(h)\coloneqq\max_{v\in B_i}(-v\cdot h),\qquad
  a_i(X)\coloneqq\int_X\varphi_i\,d\mu,\qquad
  b_i(X)\coloneqq-\int_Xh_i\,d\mu .
\]
The \emph{veto--Brier ranking} $\VB(K',B)$ on $\M_U$ is the two-level integral ranking
with integrands $\varphi_i$ and $\psi_i(h)\coloneqq-h_i$. The \emph{normalised veto set} of
world $i$ is $\hat A_i\coloneqq\cl\{\nu(v):v\in B_i\setminus\{0\}\}$ (empty if
$B_i=\{0\}$).
\end{definition}

The error of $X$ at world $i$, weighting each wrong acceptance and each silence on a
correct gamble by $|h_i|$, is $\int_{X\cap\{h_i<0\}}|h_i|\,d\mu+\int_{(U\setminus X)
\cap\{h_i>0\}}|h_i|\,d\mu$, which equals $b_i(X)$ plus a constant. So the second level is a
Brier-type error.

The main case is the following. Let $A_1,\dots,A_n$ be compact (possibly empty) sets with
$A_i\subseteq\conv(K'\cup\{e_i\})$, and put $B_i\coloneqq\conv(\{0\}\cup A_i)$. Then
\[
  \varphi_i(h)=\max\Bigl(0,\max_{v\in A_i}(-v\cdot h)\Bigr),
\]
so the veto charges an acceptance $h$ by the amount by which its expectation falls below
zero under the least favourable member of $A_i$. We write $\VB(K',A)$ for this ranking.
Here $B_i\setminus\{0\}=\{ta:t\in(0,1],\,a\in\conv A_i\}$ and $\one\cdot a=1$, so
$\hat A_i=\conv A_i$.

\begin{lemma}\label{lem:veto}
For every $i$:
\begin{enumerate}[label=(\alph*)]
  \item $\varphi_i$ is non-negative, convex, positively homogeneous and continuous, and
  $\varphi_i(h)\le c_i\max_j|h_j|$ with $c_i\coloneqq\max_{v\in B_i}\one\cdot v$;
  \item $\varphi_i(h)=0$ whenever $h\in U$ and $h_i\ge0$;
  \item $\varphi_i(h+t\one)\le\varphi_i(h)$ for $t\ge0$;
  \item $\hat A_i$ is a compact subset of $\conv(K'\cup\{e_i\})$, and
  $\{\varphi_i=0\}=C_{\hat A_i}$.
\end{enumerate}
\end{lemma}

\begin{proof}
Since $K'$ is convex, $\conv(K'\cup\{e_i\})=\{(1-\lambda)k+\lambda e_i:k\in K',\,
\lambda\in[0,1]\}$, and $\cone(K'\cup\{e_i\})=\{tx:t\ge0,\,x\in\conv(K'\cup\{e_i\})\}
\subseteq\R^n_+$.
(a) $\varphi_i$ is a maximum of linear functions over a compact set, and $0\in B_i$ gives
$\varphi_i\ge0$. For $v\ge0$, $|v\cdot h|\le(\one\cdot v)\max_j|h_j|$.
(b) Write $v\in B_i$ as $v=t((1-\lambda)k+\lambda e_i)$. For $h\in U$ with $h_i\ge0$,
$v\cdot h=t((1-\lambda)k\cdot h+\lambda h_i)\ge0$, so $\varphi_i(h)\le0$.
(c) $v\cdot\one\ge0$ for $v\in B_i$.
(d) Each $v\in B_i\setminus\{0\}$ is $tx$ with $t>0$ and $x\in\conv(K'\cup\{e_i\})$, and
$\nu(v)=x$. That set is closed, so it contains $\hat A_i$, which is therefore compact.
Finally $\varphi_i(h)=0$ iff $v\cdot h\ge0$ for all $v\in B_i$, iff $a\cdot h\ge0$ for all
$a\in\nu(B_i\setminus\{0\})$, iff (by continuity) $a\cdot h\ge0$ for all $a\in\hat A_i$.
\end{proof}

\begin{lemma}\label{lem:coh}
$\VB(K',B)$ satisfies (SC) and has null ties. For an admissible triple $(W_0,\beta,\alpha)$
put
\[
  W\coloneqq\Bigl\{\beta+\sum_{\omega\notin W_0}\alpha_\omega v_\omega:v_\omega\in
  B_\omega\Bigr\},\qquad\hat W\coloneqq\nu(W).
\]
Then $\hat W\subseteq\Delta$ is compact, and $X^\ast(W_0,\beta,\alpha)=D_{K(W_0,\beta,
\alpha)}$ a.e., where
\[
  K(W_0,\beta,\alpha)\coloneqq\conv\Bigl(K'\cup\bigcup_{\omega\in W_0}\hat A_\omega\cup
  \hat W\Bigr)
\]
is a credal set containing $K'$. In particular $X^\ast(W_0,\beta,\alpha)$ is coherent.
\end{lemma}

\begin{proof}
By Lemma~\ref{lem:veto}(a) and (M), $\varphi_i$ and $h_i$ are integrable. (SC) follows from
Lemma~\ref{lem:veto}(b). For an admissible triple, $g=\sum_{\omega\notin W_0}
\alpha_\omega\varphi_\omega-\beta\cdot h$ is continuous, and by Lemma~\ref{lem:veto}(c),
$g(h+t\one)\le g(h)-t\,\one\cdot\beta<g(h)$ for $t>0$, since $\one\cdot\beta>0$.
Lemma~\ref{lem:fubini} makes $\{g=0\}$ null; so the ranking has null ties.

Since $\alpha\ge0$, a sum of maxima over independent choices is the maximum of the sums:
\[
  g(h)=\max_{w\in W}(-w\cdot h).
\]
$W$ is compact, as a continuous image of a product of compact sets. Each $w\in W$ is
non-negative with $\one\cdot w\ge\one\cdot\beta>0$, so $\nu$ is continuous on $W$ and
$\hat W\subseteq\Delta$ is compact. As the maximum is attained, $\{g<0\}=\{h:w\cdot h>0
\text{ for all }w\in W\}=D_{\hat W}$.

Let $\omega\in W_0$. By Lemma~\ref{lem:veto}(d), $\{\varphi_\omega=0\}=C_{\hat
A_\omega}$. If $\hat A_\omega=\emptyset$, this is $\R^n=D_{\hat A_\omega}$. Otherwise
$\psi_\omega'(h)\coloneqq\min_{a\in\hat A_\omega}a\cdot h$ is continuous with
$\psi_\omega'(h+t\one)=\psi_\omega'(h)+t$, so $\{\psi_\omega'=0\}$ is null by
Lemma~\ref{lem:fubini}, and $C_{\hat A_\omega}=D_{\hat A_\omega}\cup\{\psi_\omega'=0\}$.
Either way $\{\varphi_\omega=0\}=D_{\hat A_\omega}$ a.e. By Lemma~\ref{lem:D}(a),
\[
  X^\ast(W_0,\beta,\alpha)=D_{K'}\cap\bigcap_{\omega\in W_0}D_{\hat A_\omega}\cap
  D_{\hat W}=D_{K'\cup\bigcup_{\omega\in W_0}\hat A_\omega\cup\hat W}=
  D_{K(W_0,\beta,\alpha)}\quad\text{a.e.}
\]
The union is a finite union of compact subsets of $\Delta$, so its convex hull is a
credal set.
\end{proof}

\begin{theorem}\label{thm:family}
Let $B$ be a veto profile over $K'$. The ranking $\VB(K',B)$ on $\M_U$ is complete and
transitive and satisfies (G), (IB), (GC), (C), (ID) and (W). Every model is weakly dominated
by an undominated model. A model is undominated if and only if it equals
$D_{K(W_0,\beta,\alpha)}$ a.e.\ for some admissible triple, and a model tied at every world
with an undominated model equals it a.e.
\end{theorem}

\begin{proof}
Lemma~\ref{lem:coh} and Theorem~\ref{thm:route}.
\end{proof}

\subsection{The covering theorem}\label{sec:cover}

\begin{lemma}\label{lem:LF}
Let $K\subseteq\Delta$ be full-dimensional, compact and convex, and let $K'\subseteq\Delta$
be nonempty, compact and convex, with $U=D_{K'}$.
\begin{enumerate}[label=(\alph*)]
  \item If $K'\subseteq\relint K$, then $D_K\subseteq U$ and $\partial D_K\setminus\{0\}
  \subseteq U$. Hence $D_K\in\M_U$ and (LF) holds at $D_K$ with $N=U$.
  \item Conversely, if $D_K\subseteq U$ a.e.\ and (LF) holds at $D_K$, then $K'\subseteq
  \relint K$.
\end{enumerate}
\end{lemma}

\begin{proof}
(a) $K'\subseteq K$ gives $D_K\subseteq D_{K'}$. Let $h\in\partial D_K\setminus\{0\}$. By
Lemma~\ref{lem:D}(b), $h\in C_K\setminus\{0\}$, so Lemma~\ref{lem:fulldim}(c) gives
$k'\cdot h>0$ for every $k'\in K'$, that is, $h\in U$. $U$ is open, and every measurable
$X$ with $X\triangle D_K\subseteq U$ a.e.\ satisfies $X\subseteq U$ a.e.

(b) By Lemma~\ref{lem:open}, $D_K\subseteq U$, and Lemma~\ref{lem:D}(c) gives
$K'\subseteq K$. Let $N$ be as in (LF). By (LF) the set $D_K\cup N$ belongs to $\M_U$, so
$N\subseteq U$ a.e., and $N\subseteq U$ by Lemma~\ref{lem:open}. Suppose $k'\in K'\cap\rbd
K$. By Lemma~\ref{lem:fulldim}(d) there is $h\in F_{k'}\setminus\{0\}$. Then $h\in
C_K\setminus D_K=\partial D_K$, so $h\in N\subseteq U$; but $k'\cdot h=0$, so $h\notin U$.
\end{proof}

\begin{theorem}\label{thm:cover}
Let $K\subseteq\Delta$ be a full-dimensional credal set. Suppose that $K'\subseteq\relint
K$ is nonempty, compact and convex, that $A_1,\dots,A_n$ are compact (possibly empty)
with $A_i\subseteq\conv(K'\cup\{e_i\})$, and that
\[
  K=\conv(K'\cup A_1\cup\dots\cup A_n).
\]
Then $\VB(K',A)$ covers $K$ on $\M_{D_{K'}}$, and satisfies (W). Its undominated models
are exactly the coherent models $D_{K(W_0,\beta,\alpha)}$ of Theorem~\ref{thm:family}.
\end{theorem}

\begin{proof}
Theorem~\ref{thm:family} gives every requirement except (iv) and (vi) of
Definition~\ref{def:cover}. Lemma~\ref{lem:LF}(a) gives $D_K\in\M_U$ and (LF). Fix $p\in
K'$. The triple $(\Omega,p,\emptyset)$ is admissible (there are no $\alpha$-weights), and
$W=\{p\}$, $\hat W=\{p\}$. Since $\hat A_\omega=\conv A_\omega$,
\[
  K(\Omega,p,\emptyset)=\conv\Bigl(K'\cup\bigcup_\omega\conv A_\omega\cup\{p\}\Bigr)
  =\conv(K'\cup A_1\cup\dots\cup A_n)=K .
\]
By Theorem~\ref{thm:family}, $D_K$ is undominated.
\end{proof}

\begin{remark}\label{rem:direct}
The target step can also be seen directly. If $Y$ weakly dominates $D_K$, then
$a_i(Y)\le a_i(D_K)=0$ for every $i$, so $Y\subseteq N(\Omega)=D_K$ a.e. With the vetoes
tied, $b_i(Y)\le b_i(D_K)$ for every $i$. For $p\in K$, $\sum_ip_i(b_i(Y)-b_i(D_K))=
\int_{D_K\setminus Y}p\cdot h\,d\mu$, which is positive unless $D_K\setminus Y$ is null.
Hence $Y=D_K$ a.e.
\end{remark}

\begin{corollary}\label{cor:contam}
Let $p\in\relint\Delta$, $0<\varepsilon\le1$ and $K\coloneqq(1-\varepsilon)p+
\varepsilon\Delta$, with vertices $v_i\coloneqq(1-\varepsilon)p+\varepsilon e_i$. For
$0\le s<1$ let $K'\coloneqq p+s(K-p)$ and $A_i\coloneqq\{v_i\}$. Then $\VB(K',A)$ covers
$K$ on $\M_{D_{K'}}$, in every number $n\ge2$ of worlds. At world $i$ the ranking compares
\[
  \Bigl(\int_X\bigl[-(1-\varepsilon)\,p\cdot h-\varepsilon h_i\bigr]_+d\mu(h),\
  -\int_Xh_i\,d\mu(h)\Bigr)
\]
lexicographically. The case $\varepsilon=1$ is the vacuous model $K=\Delta$.
\end{corollary}

\begin{proof}
The map $m\mapsto(1-\varepsilon)p+\varepsilon m$ is an affine bijection of $\aff\Delta$
mapping $\Delta$ onto $K$ and fixing $p$, so $K$ is full-dimensional and $p\in\relint K$.
Hence $K'\subseteq\relint K$. Also $v_i\in\conv(\{p\}\cup\{e_i\})\subseteq\conv(K'\cup
\{e_i\})$, and $K=\conv\{v_1,\dots,v_n\}=\conv(K'\cup A_1\cup\dots\cup A_n)$ because
$K'\subseteq K$. Apply Theorem~\ref{thm:cover}; the veto is $\max(0,-v_i\cdot h)$.
\end{proof}

\begin{theorem}\label{thm:exempt}
Let $K'\subseteq\Delta$ be nonempty, compact and convex, let $A_1,\dots,A_n$ be compact
with $A_i\subseteq\conv(K'\cup\{e_i\})$, and let $p\in\Delta$ be arbitrary. Put
\[
  K\coloneqq\conv(K'\cup A_1\cup\dots\cup A_n\cup\{p\}).
\]
If $K$ is full-dimensional and $K'\subseteq\relint K$, then $\VB(K',A)$ covers $K$ on
$\M_{D_{K'}}$, and satisfies (W).
\end{theorem}

\begin{proof}
The triple $(\Omega,p,\emptyset)$ is admissible, since $p\in\R^n_+\setminus\{0\}$, and
$\hat W=\{p\}$. As in the proof of Theorem~\ref{thm:cover}, $K(\Omega,p,\emptyset)=K$, so
$D_K$ is undominated by Theorem~\ref{thm:family}. Lemma~\ref{lem:LF}(a) gives $D_K\in
\M_U$ and (LF), and Theorem~\ref{thm:family} gives the remaining requirements.
\end{proof}

\begin{remark}\label{rem:exempt}
The exempt point $p$ may lie on the boundary of $\Delta$. The set
$K_0\coloneqq\conv(K'\cup A_1\cup\dots\cup A_n)$ need not be full-dimensional, and $K'$
need not lie in $\relint K_0$. The ranking does not depend on $p$: a single ranking
$\VB(K',A)$ covers every $\conv(K_0\cup\{p\})$, $p\in\Delta$, that is full-dimensional and
contains $K'$ in its relative interior.
\end{remark}

\subsection{Visibility}\label{sec:vis}

Throughout this subsection $K\subseteq\Delta$ is a full-dimensional credal set.

\begin{definition}\label{def:vis}
Let $x\in K$ and $j\in\Omega$.
\begin{enumerate}[label=(\alph*)]
  \item $x$ is \emph{visible from $j$} if $x=e_j$, or the open ray $\{x+s(x-e_j):s>0\}$
  meets $\relint K$. $x$ is \emph{visible} if it is visible from some world.
  \item $x$ is \emph{weakly visible from $j$} if $h_j\le0$ for every $h\in F_x$, and
  \emph{tangential from $j$} if it is weakly visible but not visible from $j$. It is
  \emph{anti-visible from $j$} if $h_j\ge0$ for every $h\in F_x$, and $j$ is
  \emph{crossing} at $x$ if $F_x$ contains gambles with $h_j>0$ and gambles with $h_j<0$.
  \item For $\delta>0$ let $K_\delta\coloneqq\{k\in K:\bar B(k,\delta)\cap\aff\Delta
  \subseteq K\}$, where $\bar B(k,\delta)$ is the closed Euclidean ball. A set $S\subseteq
  K$ is \emph{uniformly visible} if, for some $\delta>0$,
  \[
    S\subseteq\bigcup_{j\in\Omega}\conv(K_\delta\cup\{e_j\}).
  \]
\end{enumerate}
\end{definition}

\begin{lemma}\label{lem:visbasic}
\begin{enumerate}[label=(\alph*)]
  \item Let $x\in K$ and $x\neq e_j$. Then $x$ is visible from $j$ if and only if
  $x=(1-t)k+te_j$ for some $k\in\relint K$ and $t\in[0,1)$.
  \item Each $K_\delta$ is compact, convex and contained in $\relint K$; $K_\delta
  \supseteq K_{\delta'}$ for $\delta\le\delta'$; $K_\delta\neq\emptyset$ for small
  $\delta$; and every compact subset of $\relint K$ lies in some $K_\delta$.
  \item Let $x=(1-t)k+te_j$ with $k\in K_\delta$ and $t\in[0,1)$. Then every $y\in K$ with
  $|y-x|<(1-t)\delta/2$ lies in $\conv(K_{\delta/2}\cup\{e_j\})$.
\end{enumerate}
\end{lemma}

\begin{proof}
(a) If $k\coloneqq x+s(x-e_j)\in\relint K$ with $s>0$, then $x=\frac1{1+s}k+\frac
s{1+s}e_j$. Conversely let $x=(1-t)k+te_j$ with $k\in\relint K$. If $t>0$, then
$k=x+\frac t{1-t}(x-e_j)$ lies on the open ray. If $t=0$, then $x\in\relint K$, which is
open in $\aff\Delta$ and contains the points of the ray near $x$.

(b) Let $d(k)\coloneqq\inf\{|k-y|:y\in\aff\Delta\setminus K\}$, a $1$-Lipschitz function.
Since $K$ is closed, $K_\delta=\{k\in K:d(k)\ge\delta\}$, which is compact. If $k,k'\in
K_\delta$, $\lambda\in[0,1]$ and $y=\lambda k+(1-\lambda)k'+u$ with $|u|\le\delta$ and
$\one\cdot u=0$, then $y=\lambda(k+u)+(1-\lambda)(k'+u)\in K$; so $K_\delta$ is convex.
Clearly $K_\delta\subseteq\relint K$ and the sets increase as $\delta$ decreases. $d$ is
positive and continuous on the nonempty set $\relint K$, so it has a positive minimum on
any compact subset of $\relint K$.

(c) Put $k_y\coloneqq(y-te_j)/(1-t)\in\aff\Delta$. Then $|k_y-k|=|y-x|/(1-t)<\delta/2$, so
$\bar B(k_y,\delta/2)\subseteq\bar B(k,\delta)$ and $k_y\in K_{\delta/2}$. Moreover
$y=(1-t)k_y+te_j$.
\end{proof}

\begin{lemma}\label{lem:V1}
Let $x\in K$ and $x\neq e_j$. Then $x$ is visible from $j$ if and only if $h_j<0$ for
every nonzero $h\in F_x$. In particular, every point of $\relint K$ is visible from every
world, and visible points are weakly visible.
\end{lemma}

\begin{proof}
($\Rightarrow$) Let $k=x+s(x-e_j)\in\relint K$ with $s>0$, and $h\in F_x\setminus\{0\}$.
Then $k\cdot h=(1+s)\,x\cdot h-s\,h_j=-s\,h_j$, and $k\cdot h>0$ by
Lemma~\ref{lem:fulldim}(c). So $h_j<0$.

($\Leftarrow$) Suppose the open ray $R$ misses $\relint K$. Both are convex subsets of
$\aff\Delta$ and $\relint K$ is relatively open and nonempty. By separation in $\aff\Delta$
\cite[Thm.~11.3]{Rockafellar} there are a non-constant affine $\ell$ on $\aff\Delta$ and % citation number to be verified
$c\in\R$ with $\ell\ge c$ on $\relint K$, hence on $K$, and $\ell\le c$ on $R$. Letting
$s\to0$ along $R$ gives $\ell(x)\le c$, and $x\in K$ gives $\ell(x)\ge c$; so $\ell(x)=c$.
By Lemma~\ref{lem:fulldim}(a), $\ell-c$ is $m\mapsto m\cdot h$ for some $h\neq0$. Then
$h\in F_x$, and on $R$, $(x+s(x-e_j))\cdot h=-s\,h_j\le0$, so $h_j\ge0$.

For $x\in\relint K$, $F_x=\{0\}$ by Lemma~\ref{lem:fulldim}(c).
\end{proof}

\begin{theorem}\label{thm:vis}
The following are equivalent.
\begin{enumerate}[label=(\roman*)]
  \item There are $K'$ and $A_1,\dots,A_n$ satisfying the hypotheses of
  Theorem~\ref{thm:cover}.
  \item $\ext K$ is uniformly visible.
  \item Both of the following hold:
  \begin{enumerate}[label=(\alph*)]
    \item every point of $\cl(\ext K)$ that is not a vertex of $\Delta$ is visible;
    \item no vertex $e_i$ of $\Delta$ that lies in $K$ is an accumulation point of
    $\ext K$.
  \end{enumerate}
\end{enumerate}
Under these conditions $K$ is covered by $\VB(K_\delta,A)$ with
$A_j\coloneqq K\cap\conv(K_\delta\cup\{e_j\})$, for $\delta$ as in (ii) and small enough
that $K_\delta\neq\emptyset$.
\end{theorem}

\begin{proof}
(i)$\Rightarrow$(ii). $\ext K\subseteq K'\cup\bigcup_iA_i$, since $K$ is the convex hull of
this compact set, and $K'\cap\ext K=\emptyset$ by Lemma~\ref{lem:fulldim}(b). By
Lemma~\ref{lem:visbasic}(b), $K'\subseteq K_\delta$ for some $\delta>0$, so $\ext
K\subseteq\bigcup_i\conv(K_\delta\cup\{e_i\})$.

(ii)$\Rightarrow$(i). Decreasing $\delta$ preserves (ii), so we may assume $K_\delta\neq
\emptyset$. Put $K'\coloneqq K_\delta\subseteq\relint K$ and $A_j\coloneqq K\cap
\conv(K_\delta\cup\{e_j\})$, which is compact. Then $\ext K\subseteq\bigcup_jA_j$, and
$K=\conv\ext K\subseteq\conv(K'\cup\bigcup_jA_j)\subseteq K$. This also proves the last
sentence of the theorem, by Theorem~\ref{thm:cover}.

(ii)$\Rightarrow$(iii). Let $x\in\cl(\ext K)$, $x=\lim v_m$ with $v_m\in\ext K$, and write
$v_m=(1-t_m)k_m+t_me_{j_m}$ with $k_m\in K_\delta$ and $t_m\in[0,1]$. Passing to a
subsequence, $j_m=j$, $t_m\to t$ and $k_m\to k\in K_\delta$, so $x=(1-t)k+te_j$. For (a):
if $x$ is not a vertex of $\Delta$, then $t<1$, and $x$ is visible from $j$ by
Lemma~\ref{lem:visbasic}(a). For (b): suppose $e_i\in K$ and $v_m\to e_i$ with $v_m\neq
e_i$. Then $e_i=(1-t)k+te_j$ with $k\in\relint K\subseteq\relint\Delta$
(Lemma~\ref{lem:fulldim}(b)); since $k$ has full support and $n\ge2$, this forces $t=1$
and $j=i$. So $v_m=(1-t_m)k_m+t_me_i$ with $t_m<1$ (as $v_m\ne e_i$) and $t_m>0$ (as
$v_m\notin\relint K$). Then $v_m$ lies in the open segment between $k_m\in K$ and $e_i\in
K$, so it is not extreme, a contradiction.

(iii)$\Rightarrow$(ii). For $x\in\cl(\ext K)$ not a vertex of $\Delta$, (a) and
Lemma~\ref{lem:visbasic}(a) give $x=(1-t)k+te_j$ with $k\in\relint K$, $t\in[0,1)$; let
$\delta(x)\coloneqq d(k)>0$, so $k\in K_{\delta(x)}$. By Lemma~\ref{lem:visbasic}(c) the open
ball $O_x$ of radius $(1-t)\delta(x)/2$ about $x$ satisfies $O_x\cap K\subseteq
\conv(K_{\delta(x)/2}\cup\{e_j\})$. For a vertex $e_i\in\cl(\ext K)$, (b) gives an open
ball $O_{e_i}$ about $e_i$ meeting $\ext K$ only in $e_i$, and $e_i\in\conv(K_\delta\cup
\{e_i\})$ for every $\delta$ with $K_\delta\ne\emptyset$. The compact set $\cl(\ext K)$ is
covered by finitely many of these balls. Let $\delta$ be the least of the corresponding
numbers $\delta(x)/2$, decreased if necessary so that $K_\delta\neq\emptyset$. Since
$K_{\delta'}\subseteq K_\delta$ for $\delta'\ge\delta$, every extreme point lies in some
$\conv(K_\delta\cup\{e_j\})$.
\end{proof}

\begin{corollary}\label{cor:exemptvis}
If, for some $p\in K$, the set $\ext K\setminus\{p\}$ is uniformly visible, then $K$ is
covered by $\VB(K_\delta,A)$ with $A_j\coloneqq K\cap\conv(K_\delta\cup\{e_j\})$, for
suitable $\delta>0$.
\end{corollary}

\begin{proof}
Choose $\delta$ for $\ext K\setminus\{p\}$ with $K_\delta\neq\emptyset$, and put
$K'\coloneqq K_\delta$. Then $\ext K\subseteq\bigcup_jA_j\cup\{p\}$, so
$K=\conv(K'\cup\bigcup_jA_j\cup\{p\})$, and $K'\subseteq\relint K$. Apply
Theorem~\ref{thm:exempt}.
\end{proof}

For polytopes, visibility can be read off the facet normals.

\begin{lemma}\label{lem:poly}
Suppose $K=\{m\in\aff\Delta:g_f\cdot m\ge0\text{ for }f\in I\}$ with $I$ finite. Then
$C_K=\cone\{g_f:f\in I\}+\R_+\one$, and for $x\in K$,
\[
  F_x=\cone\{g_f:f\in I,\ g_f\cdot x=0\}.
\]
Hence $x\neq e_j$ is visible from $j$ if and only if $(g_f)_j<0$ for every $f$ with
$g_f\neq0$ and $g_f\cdot x=0$.
\end{lemma}

\begin{proof}
By the affine Farkas lemma \cite[Cor.~7.1h]{Schrijver}, applied to the system $g_f\cdot
m\ge0$, $\one\cdot m\ge1$, $-\one\cdot m\ge-1$, the inequality $h\cdot m\ge0$ holds on
$K$ iff $h=\sum_f\lambda_fg_f+(\eta_+-\eta_-)\one$ with $\lambda,\eta_\pm\ge0$ and
$\eta_+-\eta_-\ge0$. This is the formula for $C_K$. If such an $h$ has $x\cdot h=0$, then
$\sum_f\lambda_f\,g_f\cdot x+(\eta_+-\eta_-)=0$ with non-negative terms, so $\lambda_f=0$
unless $g_f\cdot x=0$, and the $\one$-term vanishes. The visibility criterion follows from
Lemma~\ref{lem:V1}: a nonzero element of $F_x$ has a positive coefficient on some nonzero
active $g_f$, and each nonzero active $g_f$ lies in $F_x$.
\end{proof}

Every polytope $K\subseteq\Delta$ has such a description: an affine inequality $a\cdot
m\ge c$ on $\aff\Delta$ is $(a-c\one)\cdot m\ge0$.

\begin{corollary}\label{cor:poly}
Let $K$ be a full-dimensional polytope. If every vertex of $K$, with at most one
exception, is a vertex of $\Delta$ or visible, then $K$ is covered by a veto--Brier
ranking. If there is no exception, the hypotheses of Theorem~\ref{thm:cover} hold.
\end{corollary}

\begin{proof}
A visible vertex $v\neq e_j$ has the form $(1-t)k+te_j$ with $k\in\relint K$, so
$v\in\conv(K_{d(k)}\cup\{e_j\})$; and $e_j\in\conv(K_\delta\cup\{e_j\})$. Taking the least
of finitely many depths, $\ext K$ minus the exception is uniformly visible. Apply
Corollary~\ref{cor:exemptvis} or Theorem~\ref{thm:vis}.
\end{proof}

\begin{corollary}\label{cor:smooth}
Let $K\subseteq\relint\Delta$ be full-dimensional, and suppose $K$ has exactly one
supporting hyperplane (in $\aff\Delta$) at each point of $\rbd K$. Then every point of $K$
is visible, the hypotheses of Theorem~\ref{thm:cover} hold, and $K$ is covered.
\end{corollary}

\begin{proof}
Let $x\in\rbd K$. By Lemma~\ref{lem:fulldim}(d) there is $h_x\in F_x\setminus\{0\}$. If
$h\in F_x\setminus\{0\}$, the affine functions $m\mapsto m\cdot h$ and $m\mapsto m\cdot
h_x$ on $\aff\Delta$ vanish on the same supporting hyperplane, so they are proportional;
both are non-negative on $K$, which has points off that hyperplane, so the factor is
positive, and by Lemma~\ref{lem:fulldim}(a) $h$ is a positive multiple of $h_x$. Thus
$F_x=\R_+h_x$. Since $x$ has full support, $x\cdot h_x=0$ and $h_x\neq0$, some coordinate
$(h_x)_j$ is negative, and $x$ is visible from $j$ by Lemma~\ref{lem:V1}. Points of
$\relint K$ are visible by Lemma~\ref{lem:V1}. No vertex of $\Delta$ lies in $K$, so
condition (iii) of Theorem~\ref{thm:vis} holds.
\end{proof}

Examples of such $K$, whenever they lie in $\relint\Delta$, are Euclidean balls
$\{m\in\aff\Delta:|m-c|\le r\}$; outer parallel bodies $\{m\in\aff\Delta:\operatorname{dist}
(m,L)\le r\}$ of a compact convex $L\subseteq\aff\Delta$ with $r>0$ (each boundary point
lies on the boundary of a ball of radius $r$ contained in the body, and any supporting
hyperplane there supports that ball); and compact sublevel sets $\{m\in\aff\Delta:G(m)\le
r\}$, $r>\inf G$, of a function $G$ that is convex and differentiable on an open convex
subset $V$ of $\aff\Delta$ containing the set. In the last
case, at a boundary point $x$ we have $G(x)=r$ and $\nabla G(x)\ne0$ (a critical point of
a convex function is a minimiser); if $y$ is an outward normal at $x$, then $y\cdot d\le0$
for every direction $d$ with $\nabla G(x)\cdot d<0$, since $x+td$ lies in the set for
small $t>0$, so $y$ is a non-negative multiple of $\nabla G(x)$. The last class includes
ellipsoids and balls $\{m:\mathrm{KL}(m\|q)\le r\}$ for the Kullback--Leibler divergence,
$q\in\relint\Delta$, with $V=\relint\Delta$.

\begin{corollary}\label{cor:n2}
If $n=2$, every full-dimensional credal set $K\subseteq\Delta$ is covered by a
veto--Brier ranking satisfying the hypotheses of Theorem~\ref{thm:cover}.
\end{corollary}

\begin{proof}
$K$ is a segment with endpoints $x,y$, $x_1<y_1$. If $y\neq e_1$, the points
$y+s(y-e_1)$ have first coordinate $y_1-s(1-y_1)\in(x_1,y_1)$ for small $s>0$, so they lie in
$\relint K$ and $y$ is visible from $1$. Symmetrically $x$ is $e_2$ or visible from $2$.
Apply Corollary~\ref{cor:poly}.
\end{proof}

\begin{remark}\label{rem:qual}
Three qualifications in these statements matter.
\begin{enumerate}[label=(\arabic*)]
  \item \emph{Strict, not weak, visibility.} Visible points are weakly visible, but not
  conversely. Let $K=\conv\{v,w,x\}$ with $v=(\frac25,\frac3{10},\frac3{10})$,
  $w=(\frac15,\frac25,\frac25)$ and $x=(\frac9{20},\frac9{20},\frac1{10})$. Its facet
  normals at $v$ are $(0,1,-1)$ (through $v,w$) and $(-21,19,9)$ (through $v,x$). Both have
  non-positive first coordinate, so $v$ is weakly visible from world $1$; but
  $w=v+\frac13(v-e_1)$, so the ray from $e_1$ through $v$ runs along the edge $[v,w]$ and
  misses $\relint K$. By Lemma~\ref{lem:poly}, $v$ is not visible from worlds $2$ or $3$
  either. So $v$ is tangential from world $1$ and is not visible; by
  Lemma~\ref{lem:visbasic}(a) it lies in no set $\conv(K''\cup\{e_j\})$ with $K''$ a
  compact subset of $\relint K$. Tangential points count as not visible in
  Theorems~\ref{thm:cover} and~\ref{thm:vis}, and also in Theorem~\ref{thm:reach} below.
  \item \emph{Uniformity.} Theorem~\ref{thm:vis} requires visibility on $\cl(\ext K)$, not
  only on $\ext K$, together with condition (iii)(b) at vertices of $\Delta$. For
  polytopes both reduce to visibility of the vertices.
  \item \emph{Dimension and interiority.} Full-dimensionality of $K$ is assumed throughout
  this section; it is forced by (LF) on one-sided domains
  (Proposition~\ref{prop:lowdim}). The assumption $K\subseteq\relint\Delta$ is used only in
  Corollary~\ref{cor:smooth} and in the exclusions of Section~\ref{sec:excl}.
\end{enumerate}
\end{remark}

\subsection{Examples}\label{sec:examples}

Interior $\varepsilon$-contaminations, including the vacuous model, are covered by
Corollary~\ref{cor:contam}, in every number of worlds. Smooth credal sets in
$\relint\Delta$ and all full-dimensional credal sets on two worlds are covered by
Corollaries~\ref{cor:smooth} and~\ref{cor:n2}.

\begin{proposition}\label{prop:tv}
Let $p\in\relint\Delta$, $0<\delta<\min_jp_j$, and $K\coloneqq\{m\in\Delta:\frac12\sum_i
|m_i-p_i|\le\delta\}$. Then $K$ is a full-dimensional polytope in $\relint\Delta$ whose
vertices are all visible, and $K$ is covered, for every $n\ge2$.
\end{proposition}

\begin{proof}
Every $m\in\aff\Delta$ with $\frac12\sum|m_i-p_i|\le\delta$ has $m_j\ge p_j-\delta>0$, so
$K=p+\delta Q$ with $Q\coloneqq\{d:\one\cdot d=0,\ \sum|d_i|\le2\}$, and $K\subseteq
\relint\Delta$ is full-dimensional. If $d\in Q$, $d\ne0$, write $d=d^+-d^-$ with
$\one\cdot d^+=\one\cdot d^-\eqqcolon\sigma\le1$; then $d=\sum_{i\neq j}\frac{d^+_id^-_j}\sigma
(e_i-e_j)$, with weights summing to $\sigma$, and $0=\frac12(e_1-e_2)+\frac12(e_2-e_1)$.
So $Q=\conv\{e_i-e_j:i\neq j\}$, and every vertex of $K$ has the form $v=p+\delta(e_i-e_j)$.
For small $s>0$ the point $y\coloneqq v+s(v-e_i)$ satisfies
\[
  y-p=(1+s)\delta(e_i-e_j)+s(p-e_i),\qquad
  \tfrac12\textstyle\sum_k|y_k-p_k|=\delta+s(\delta-p_j)<\delta ,
\]
by a direct computation of the coordinates $i$, $j$ and $k\notin\{i,j\}$. So $y\in\relint K$
and $v$ is visible from $i$. Apply Corollary~\ref{cor:poly}.
\end{proof}

Let $\underline P:2^\Omega\to[0,1]$ with $\underline P(\emptyset)=0$ and $\underline
P(\Omega)=1$, and put $K\coloneqq\{m\in\aff\Delta:m(A)\ge\underline P(A)\text{ for all
}A\subseteq\Omega\}$, where $m(A)\coloneqq\sum_{k\in A}m_k$. The constraints on singletons
give $K\subseteq\Delta$. An event $A$ is \emph{active} at $x\in K$ if $x(A)=\underline
P(A)$.

\begin{lemma}\label{lem:events}
Suppose $K$ is full-dimensional. A point $x\in K$ with $x\neq e_j$ is visible from $j$ if
and only if every active event $A$ with $\emptyset\neq A\neq\Omega$ satisfies $j\notin A$
and $\underline P(A)>0$.
\end{lemma}

\begin{proof}
On $\aff\Delta$, $m(A)-\underline P(A)=m\cdot g_A$ with $g_A\coloneqq1_A-\underline
P(A)\one$, so Lemma~\ref{lem:poly} applies. $g_\emptyset=g_\Omega=0$, and $g_A\neq0$ for
proper nonempty $A$. For such $A$, $\underline P(A)<1$: otherwise $m(\Omega\setminus A)=0$
on $K$, and $K$ would not be full-dimensional. Hence $(g_A)_j=1-\underline P(A)>0$ if
$j\in A$, and $(g_A)_j=-\underline P(A)$ if $j\notin A$.
\end{proof}

For a permutation $\sigma$ of $\Omega$ let $S_k\coloneqq\{\sigma(1),\dots,\sigma(k)\}$ and
define the marginal vector $v^\sigma$ by $v^\sigma_{\sigma(k)}\coloneqq\underline P(S_k)
-\underline P(S_{k-1})$. If $\underline P$ is supermodular, that is, $\underline P(A\cup
B)+\underline P(A\cap B)\ge\underline P(A)+\underline P(B)$, then $K$ is the core of
$\underline P$ and its extreme points are exactly the vectors $v^\sigma$
\cite{Shapley,Ichiishi}.

\begin{proposition}\label{prop:bel}
Let $\rho$ be a mass function on the nonempty subsets of $\Omega$, $\mathrm{Bel}(A)
\coloneqq\sum_{\emptyset\neq F\subseteq A}\rho(F)$, and suppose $\rho(\Omega)>0$ and
$\rho(\{k\})>0$ for every $k$. Then the credal set $K$ of $\mathrm{Bel}$ is a
full-dimensional polytope in $\relint\Delta$, each vertex $v^\sigma$ is visible from
$\sigma(n)$, and $K$ is covered.
\end{proposition}

\begin{proof}
Belief functions are supermodular \cite{Shafer}, so $\ext K=\{v^\sigma\}$. Since
$\mathrm{Bel}(S_k)-\mathrm{Bel}(S_{k-1})$ sums $\rho(F)$ over the $F$ whose
$\sigma$-last element is $\sigma(k)$,
\[
  v^\sigma(A)=\sum_{F:\,\mathrm{last}_\sigma(F)\in A}\rho(F)\ \ge\
  \sum_{F\subseteq A}\rho(F)=\mathrm{Bel}(A).
\]
If $\sigma(n)\in A\neq\Omega$, the set $F=\Omega$ contributes to the left side but not
the right, so $v^\sigma(A)\ge\mathrm{Bel}(A)+\rho(\Omega)>\mathrm{Bel}(A)$. Hence every
active proper event at $v^\sigma$ omits $\sigma(n)$, and it has $\mathrm{Bel}(A)\ge
\sum_{k\in A}\rho(\{k\})>0$. For each proper nonempty $A$ some $\sigma$ has $\sigma(n)
\in A$, so the average of all $v^\sigma$ satisfies every constraint strictly; it lies in
$\relint K$, and $K$ is full-dimensional. Also $m_k\ge\rho(\{k\})>0$ on $K$. By
Lemma~\ref{lem:events} each $v^\sigma$ is visible from $\sigma(n)$; apply
Corollary~\ref{cor:poly}.
\end{proof}

\begin{proposition}\label{prop:2mon}
Suppose $\underline P$ is \emph{strictly supermodular}: $\underline P(A\cup B)+
\underline P(A\cap B)>\underline P(A)+\underline P(B)$ whenever neither of $A,B$ contains
the other. Suppose also $\underline P(\{k\})>0$ for every $k$. Then $K$ is a
full-dimensional polytope in $\relint\Delta$, each vertex $v^\sigma$ is visible from
$\sigma(n)$, and $K$ is covered.
\end{proposition}

\begin{proof}
$\ext K=\{v^\sigma\}$ by supermodularity. Since $\underline P\ge0$ is supermodular, it is
monotone: for $A\subseteq B$, $\underline P(B)\ge\underline P(A)+\underline P(B\setminus
A)\ge\underline P(A)$. Let $A$ be nonempty, with elements $\sigma(k_1),\dots,\sigma(k_r)$,
$k_1<\dots<k_r$, and put $C_t\coloneqq A\cap S_{k_t}$. Applying supermodularity to
$S_{k_t-1}$ and $C_t$, whose union is $S_{k_t}$ and whose intersection is $C_{t-1}$,
\[
  \underline P(S_{k_t})-\underline P(S_{k_t-1})\ \ge\ \underline P(C_t)-\underline
  P(C_{t-1}),
\]
strictly unless $S_{k_t-1}\subseteq C_t$, that is, unless $S_{k_t-1}\subseteq A$ (note
$C_t\not\subseteq S_{k_t-1}$). Summing over $t$ gives $v^\sigma(A)\ge\underline P(A)$,
with equality only if $S_{k_r-1}\subseteq A\subseteq S_{k_r}$, so only if $A=S_{k_r}$ is an
initial segment of $\sigma$. The proper initial segments omit $\sigma(n)$ and have
$\underline P(S_k)\ge\underline P(\{\sigma(1)\})>0$. Each proper nonempty $A$ is not an
initial segment for some $\sigma$ (one with $\sigma(n)\in A$), so averaging the
$v^\sigma$ gives a point of $\relint K$, and $K$ is full-dimensional; $K\subseteq
\relint\Delta$ as before. Apply Lemma~\ref{lem:events} and Corollary~\ref{cor:poly}.
\end{proof}

\begin{corollary}\label{cor:distortion}
Let $p\in\relint\Delta$ and let $f:[0,1]\to[0,1]$ be strictly convex with $f(0)=0$,
$f(1)=1$ and $f>0$ on $(0,1]$. Then the credal set of $\underline P\coloneqq f\circ p$,
$\underline P(A)=f(p(A))$, is covered. This includes $f(x)=(1-\tau)x/(1-\tau x)$,
$0<\tau<1$, the event lower probability of the constant-odds-ratio model with reference
$p$ \cite{Walley}.
\end{corollary}

\begin{proof}
Let $A,B$ be non-nested, $x\coloneqq p(A\cap B)$, $a\coloneqq p(A\setminus B)>0$ and
$b\coloneqq p(B\setminus A)>0$. With $\lambda\coloneqq b/(a+b)\in(0,1)$, $x+a=\lambda x+
(1-\lambda)(x+a+b)$ and $x+b=(1-\lambda)x+\lambda(x+a+b)$. Strict convexity gives
$f(x+a)+f(x+b)<f(x)+f(x+a+b)$, which is strict supermodularity. Also $f(p_k)>0$. Apply
Proposition~\ref{prop:2mon}. For the example, $f''(x)=2\tau(1-\tau)(1-\tau x)^{-3}>0$.
\end{proof}

The covered set in the constant-odds-ratio case is the core of the event lower
probability. It differs from the density-ratio class $\{m:m_i/p_i\le km_j/p_j\}$,
$k=1/(1-\tau)$: for $n=3$, $p$ uniform and $\tau=\frac12$, the core contains
$(\frac15,\frac3{10},\frac12)$, which violates $m_3\le2m_1$. Density-ratio classes are not
reached by veto--Brier rankings (Corollary~\ref{cor:out}); they are covered by the
orthant-diagonal rankings of Section~\ref{sec:op} (Proposition~\ref{prop:box}).

\begin{proposition}\label{prop:pri}
Let $K\coloneqq\{m\in\Delta:l_i\le m_i\le u_i\}$ be full-dimensional with $K\subseteq
\relint\Delta$. A vertex of $K$ is visible if and only if at most one of its coordinates is
at its upper bound. Hence $K$ is covered if every vertex, with at most one exception, has
at most one coordinate at its upper bound.
\end{proposition}

\begin{proof}
On $\aff\Delta$ the constraints are $m\cdot(e_i-l_i\one)\ge0$ and $m\cdot(u_i\one-e_i)
\ge0$; at a vertex $v\in\relint\Delta$ the constraints $m_k\ge0$ are inactive. Let $L$ and
$H$ be the sets of coordinates of $v$ at their lower and upper bounds. For $i\in H$,
$u_i=v_i<1$, so $u_i\one-e_i$ is negative exactly at coordinate $i$. For $i\in L$,
$l_i=v_i>0$, so $e_i-l_i\one$ is negative exactly off $i$. By Lemma~\ref{lem:poly}, $v$ is
visible from $j$ iff $H\subseteq\{j\}$ and $j\notin L$. If $|H|\ge2$ no $j$ works. If
$H=\{i\}$, then $j=i\notin L$ works, since $l_i<u_i$ by full-dimensionality. If
$H=\emptyset$, then $L\neq\Omega$ (otherwise $\sum_il_i=1$ and $K=\{l\}$), and any $j\notin
L$ works. Apply Corollary~\ref{cor:poly}.
\end{proof}

\begin{corollary}\label{cor:hexagon}
Let $n=3$. Every probability-interval model $K\subseteq\relint\Delta$ that is a hexagon is
covered by a ranking satisfying the hypotheses of Theorem~\ref{thm:cover}.
\end{corollary}

\begin{proof}
Each edge of $K$ lies on one of the six lines $m_i=l_i$, $m_i=u_i$, and the six edges have
distinct outward normals, so each line carries exactly one edge. Within $\aff\Delta$ the
outward normal of the edge on $m_i=l_i$ is $-\pi(e_i)$ and that on $m_i=u_i$ is $\pi(e_i)$,
where $\pi$ is the orthogonal projection onto $\one^\perp$. The vectors $\pi(e_i)$ have
equal length and pairwise angles of $120^\circ$, so in angular order the six normals are
$\pi(e_1),-\pi(e_3),\pi(e_2),-\pi(e_1),\pi(e_3),-\pi(e_2)$, alternating in sign. Consecutive
edges meet at a vertex, and no vertex lies on a third line (a supporting line through a
vertex that carries an edge contains that vertex in the edge, and a vertex lies on exactly
two edges). So every vertex has exactly one coordinate at its upper bound. Apply
Proposition~\ref{prop:pri}.
\end{proof}

The \emph{reversed triangle} is $R\coloneqq\{m\in\Delta_3:m_i\le\frac25\}$, with vertices
$v_1=(\frac15,\frac25,\frac25)$, $v_2=(\frac25,\frac15,\frac25)$ and
$v_3=(\frac25,\frac25,\frac15)$. It is the pari-mutuel model with uniform reference and
$\delta=\frac15$, and lies in $\relint\Delta$. No vertex of $R$ is visible
(Corollary~\ref{cor:out}). Nearby targets are covered.

\begin{proposition}\label{prop:trunc}
For $0<\eta<\frac1{10}$ let $R_\eta\coloneqq\{m\in\Delta:\frac15+\eta\le m_i\le\frac25\}$.
\begin{enumerate}[label=(\alph*)]
  \item $R_\eta$ is covered by a ranking satisfying the hypotheses of
  Theorem~\ref{thm:cover}, and $d_H(R_\eta,R)=\eta\sqrt{3/2}$.
  \item $P_\eta\coloneqq\conv(R_\eta\cup\{v_1\})$ is covered. It coincides with $R$ on
  $\{m_1\le\frac15+\eta\}$, and $d_H(P_\eta,R)=\eta\sqrt{3/2}$.
\end{enumerate}
For $\eta=\frac1{10}$, $R_\eta=\{m:m_i\ge\frac3{10}\}=\frac9{10}u+\frac1{10}\Delta$, $u$
uniform, is covered by Corollary~\ref{cor:contam}.
\end{proposition}

\begin{proof}
(a) $R_\eta\subseteq\relint\Delta$ contains $u$ in its relative interior. A vertex has two
active constraints; if two coordinates were at the upper bound $\frac25$, the third would
be $\frac15<\frac15+\eta$. So every vertex has at most one coordinate at its upper bound,
and Proposition~\ref{prop:pri} applies (with no exception). The distance from a point of
$R$ to $R_\eta$ is a convex function, so $d_H(R_\eta,R)$ is attained at a vertex of $R$,
by symmetry at $v_1$. Since $R_\eta\subseteq\{m_1\ge\frac15+\eta\}$ and $|\pi(e_1)|=
\sqrt{2/3}$, $\operatorname{dist}(v_1,R_\eta)\ge\eta\sqrt{3/2}$, with equality at
$(\frac15+\eta,\frac25-\frac\eta2,\frac25-\frac\eta2)\in R_\eta$.

(b) By (a) and Theorem~\ref{thm:vis} there are $K'\subseteq\relint R_\eta$ and
$A_i\subseteq\conv(K'\cup\{e_i\})$ with $R_\eta=\conv(K'\cup\bigcup A_i)$. Then
$P_\eta=\conv(K'\cup\bigcup A_i\cup\{v_1\})$ is full-dimensional and $K'\subseteq
\relint R_\eta\subseteq\relint P_\eta$, so Theorem~\ref{thm:exempt} applies.
$R\cap\{m_1\le\frac15+\eta\}$ is the triangle with vertices $v_1$,
$(\frac15+\eta,\frac25,\frac25-\eta)$ and $(\frac15+\eta,\frac25-\eta,\frac25)$, whose last
two vertices lie in $R_\eta$; so it is contained in $P_\eta\subseteq R$. The distance
claim follows as in (a), at $v_2$.
\end{proof}

\subsection{Dense approximation}\label{sec:dense}

\begin{proposition}\label{prop:dense}
Let $K\subseteq\Delta$ be a full-dimensional credal set, $u_0\in\relint K$ and
$s\in(0,1)$. Put $K'_s\coloneqq u_0+s(K-u_0)$, $A_i\coloneqq K\cap\conv(K'_s\cup\{e_i\})$
and $K_s\coloneqq\conv(K'_s\cup A_1\cup\dots\cup A_n)$. Then $K'_s\subseteq\relint K_s$,
$K_s\subseteq K$,
\[
  d_H(K_s,K)\le(1-s)\max_{k\in K}|k-u_0|,
\]
and $\VB(K'_s,A)$ covers $K_s$ on $\M_{D_{K'_s}}$. In particular every full-dimensional
credal set is a Hausdorff limit of covered polytopes.
\end{proposition}

\begin{proof}
$K'_s$ is compact, convex, full-dimensional and contained in $\relint K$, since
$(1-s)u_0+sk\in\relint K$ for $k\in K$. Each $A_i$ is compact and contained in
$\conv(K'_s\cup\{e_i\})$, and $K'_s\subseteq K_s\subseteq K$. For $k\in K$, the point
$u_0+s(k-u_0)\in K_s$ lies at distance $(1-s)|k-u_0|$ from $k$, which gives the bound.

Suppose $x\in K'_s\cap\rbd K_s$. There is a supporting hyperplane of $K_s$ at $x$, given by
Lemma~\ref{lem:fulldim}(a) as $c\cdot m\le c\cdot x$ on $K_s$ with $m\mapsto c\cdot m$
non-constant on $\aff\Delta$, that is, $c$ not a multiple of $\one$. Since $x\in\relint
K\subseteq\relint\Delta$ has full support, $c\cdot x<\max_ic_i\eqqcolon c_{i^\ast}$. For
small $t>0$, $y_t\coloneqq(1-t)x+te_{i^\ast}\in K$ because $x\in\relint K$, so $y_t\in
A_{i^\ast}\subseteq K_s$; but $c\cdot y_t=c\cdot x+t(c_{i^\ast}-c\cdot x)>c\cdot x$, a
contradiction. Hence $K'_s\subseteq\relint K_s$, $K_s$ is full-dimensional, and
Theorem~\ref{thm:cover} applies. If $K$ is a polytope, so is each $A_i$ (an intersection
of polytopes), and hence so is $K_s$. Applying this to full-dimensional polytopes
approximating a general $K$ gives the last claim.
\end{proof}

The domain $D_{K'_s}$ depends on $s$: the approximation is by targets covered on their own
domains. For the reversed triangle, Proposition~\ref{prop:trunc} gives an explicit family
with exact distances.

\section{Orthant-planar targets}\label{sec:op}

\subsection{Orthant-planar credal sets}

For a sign vector $s\in\{-1,1\}^n$ let $\sigma_s\coloneqq\{h:s_ih_i>0\text{ for all }i\}$,
an open orthant. It is \emph{pure} if $s$ is constant and \emph{mixed} otherwise. For a
mixed orthant $\sigma$ we write $P(\sigma)\coloneqq\{i:s_i=1\}$ and
$N(\sigma)\coloneqq\{i:s_i=-1\}$, both nonempty. Let $O\coloneqq\bigcup_s\sigma_s$, the
gambles with no zero coordinate; its complement $Z\coloneqq\R^n\setminus O$ is null.

\begin{definition}\label{def:op}
A credal set $K\subseteq\relint\Delta$ is \emph{orthant-planar} if for every mixed orthant
$\sigma$ there is $d_\sigma\in\R^n$ with positive coordinates such that
\[
  D_K\cap\sigma=\{h:d_\sigma\cdot h>0\}\cap\sigma .
\]
An \emph{orthant-diagonal chart} for $K$ is the map $c:\R^n\to\R^n$ with
$c(h)\coloneqq d_\sigma\odot h$ for $h\in\sigma$, where $d_\sigma$ is as above for mixed
$\sigma$ and is an arbitrary vector with positive coordinates for the two pure orthants, and
$c(h)\coloneqq h$ for $h\in Z$.
\end{definition}

On the pure orthants nothing is required, since $D_K$ contains the positive orthant and
misses the negative one. On each open orthant $\sigma$, $c$ is the restriction of the
invertible linear map $\operatorname{diag}(d_\sigma)$, which maps $\sigma$ onto $\sigma$;
so $c$ maps $O$ bijectively onto $O$, $c_i(h)$ has the sign of $h_i$ for every $h$, and
$|c(h)|\le\max_\sigma\max_jd_{\sigma,j}\,|h|$.

\begin{lemma}\label{lem:closedorthant}
Let $K\subseteq\relint\Delta$ be orthant-planar and $\sigma$ a mixed orthant. Then
\[
  D_K\cap\cl\sigma=\{x:d_\sigma\cdot x>0\}\cap\cl\sigma .
\]
\end{lemma}

\begin{proof}
Note that $x+\varepsilon s\in\sigma$ whenever $x\in\cl\sigma$, $s\in\sigma$ and
$\varepsilon>0$.
($\subseteq$) Let $x\in D_K\cap\cl\sigma$ and suppose $d_\sigma\cdot x\le0$. Since
$\sigma$ is mixed and $d_\sigma$ has positive coordinates, some $s\in\sigma$ has
$d_\sigma\cdot s<0$. For small $\varepsilon>0$, $y\coloneqq x+\varepsilon s$ lies in
$\sigma\cap D_K$, since $D_K$ is open, but $d_\sigma\cdot y<0$, a contradiction.
($\supseteq$) Let $x\in\cl\sigma$ with $d_\sigma\cdot x>0$, and suppose $x\notin D_K$. Pick
$m\in K$ with $m\cdot x\le0$; since $m$ has positive coordinates, some $s\in\sigma$ has
$m\cdot s<0$. For small $\varepsilon>0$, $y\coloneqq x+\varepsilon s\in\sigma$ has
$d_\sigma\cdot y>0$ and $m\cdot y<0$, so $y\in(\{d_\sigma\cdot h>0\}\cap\sigma)\setminus
D_K$, a contradiction.
\end{proof}

For a compact $Q\subseteq\Delta$ and a compact $K\subseteq\relint\Delta$ put
\[
  Q\star K\coloneqq\conv\{\nu(q\odot m):q\in Q,\ m\in K\}.
\]
The generating set is compact, as the image of $Q\times K$ under the continuous map
$(q,m)\mapsto\nu(q\odot m)$ (note $q\cdot m\ge\min_{m\in K}\min_im_i>0$). So $Q\star K$
is a credal set. For $q\in\Delta$, $\{q\}\star K=\nu(q\odot K)$, the image of $K$ under
$m\mapsto\nu(q\odot m)$, which is convex by~\eqref{eq:proj}.

\begin{lemma}\label{lem:pullback}
Let $K\subseteq\relint\Delta$ be orthant-planar with chart $c$, and let $Q\subseteq\Delta$
be nonempty and compact. For every $h\in O$,
\[
  c(h)\in D_Q\iff h\in D_{Q\star K}.
\]
In particular, for $\pi\in\R^n_+\setminus\{0\}$, $\{h\in O:\pi\cdot c(h)>0\}=
D_{\nu(\pi\odot K)}\cap O$.
\end{lemma}

\begin{proof}
Fix $q\in Q$ and $h\in\sigma$. If $\sigma$ is the positive orthant, $q\cdot c(h)>0$ and
$(q\odot m)\cdot h>0$ for all $m\in K$; if it is the negative orthant, both are negative.
Let $\sigma$ be mixed. Then $q\cdot c(h)=d_\sigma\cdot x$ with $x\coloneqq q\odot h$, and
$x\in\cl\sigma\setminus\{0\}$ because $h$ has no zero coordinate and $q\neq0$. By
Lemma~\ref{lem:closedorthant}, $q\cdot c(h)>0$ iff $x\in D_K$, iff $m\cdot(q\odot h)=
(q\odot m)\cdot h>0$ for every $m\in K$, iff $\nu(q\odot m)\cdot h>0$ for every $m\in K$.
Taking all $q\in Q$, $c(h)\in D_Q$ iff $h\in D_S$ with $S\coloneqq\{\nu(q\odot m):q\in
Q,m\in K\}$, and $D_S=D_{Q\star K}$ by Lemma~\ref{lem:D}(a). The last claim is the case
$Q=\{\nu(\pi)\}$.
\end{proof}

\begin{lemma}\label{lem:ratio}
Let $K\subseteq\relint\Delta$ be a credal set and $\sigma$ a mixed orthant, with
$P=P(\sigma)$ and $N=N(\sigma)$.
\begin{enumerate}[label=(\alph*)]
  \item If some $d\in K$ satisfies $m_bd_a\le d_bm_a$ for all $m\in K$, $a\in P$ and $b\in
  N$ (that is, $d$ maximises every ratio $m_b/m_a$ with $a\in P$, $b\in N$), then
  $D_K\cap\sigma=\{d\cdot h>0\}\cap\sigma$.
  \item Conversely, if $D_K\cap\sigma=\{d\cdot h>0\}\cap\sigma$ for some $d$ with positive
  coordinates, then $\nu(d)\in K$, and $\nu(d)$ maximises every ratio $m_b/m_a$ with $a\in
  P$, $b\in N$.
\end{enumerate}
Hence $K$ is orthant-planar if and only if, for every mixed orthant, some point of $K$
maximises all the corresponding ratios; and we may then take $d_\sigma\in K$.
\end{lemma}

\begin{proof}
(a) $D_K\cap\sigma\subseteq\{d\cdot h>0\}$ because $d\in K$. Conversely let $h\in\sigma$
with $d\cdot h>0$ and $m\in K$. Put $\lambda\coloneqq\min_{a\in P}m_a/d_a>0$. Then
$m_a\ge\lambda d_a$ for $a\in P$, and for $b\in N$ and the $a$ attaining the minimum,
$m_b\le(m_a/d_a)d_b=\lambda d_b$. As $h_a>0$ on $P$ and $h_b<0$ on $N$,
$m\cdot h\ge\lambda\,d\cdot h>0$.

(b) First, $m_b/m_a\le d_b/d_a$ for $m\in K$, $a\in P$, $b\in N$. Fix $w\in\sigma$ and
$x>d_b/d_a$, and put $h_\delta\coloneqq xe_a-e_b+\delta w$. For small $\delta>0$,
$h_\delta\in\sigma$ and $d\cdot h_\delta=xd_a-d_b+\delta\,d\cdot w>0$, so $h_\delta\in D_K$
and $xm_a-m_b+\delta\,m\cdot w>0$. Letting $\delta\to0$ gives $xm_a\ge m_b$; letting
$x\downarrow d_b/d_a$ gives the claim. Next, choose $h\in\sigma$ with $d\cdot h=0$ (possible
as $\sigma$ is mixed), and $s\in\sigma$ with $d\cdot s>0$. Then $h+ts\in\sigma\cap
\{d\cdot y>0\}\subseteq D_K$ for $t>0$, so $h\in C_K$, while $h\notin D_K$. So $m\cdot h=0$
for some $m\in K$. With $\lambda\coloneqq\min_{a\in P}m_a/d_a$ the argument of (a) gives
$m_a\ge\lambda d_a$ on $P$ and $m_b\le\lambda d_b$ on $N$, so $0=m\cdot h\ge\lambda\,d\cdot
h=0$, with equality only if $m_a=\lambda d_a$ and $m_b=\lambda d_b$ for all $a\in P$,
$b\in N$, because $h$ has no zero coordinate. So $m=\lambda d$, $\lambda>0$, and
$\nu(d)=m\in K$. Its ratios are those of $d$.
\end{proof}

\begin{lemma}\label{lem:planarobs}
Let $K\subseteq\relint\Delta$ be orthant-planar.
\begin{enumerate}[label=(\alph*)]
  \item $K$ is a polytope.
  \item If $g\in\R^n$ has no zero coordinate and $\min_{m\in K}m\cdot g=0$, then this
  minimum is attained at only one point of $K$.
  \item Every edge of $K$ lies on a line through a vertex of $\Delta$.
  \item No facet normal of a full-dimensional orthant-planar $K$ has all coordinates
  nonzero, when $n\ge3$.
\end{enumerate}
\end{lemma}

\begin{proof}
(a) For a mixed orthant $\sigma$, $D_K\cap\sigma$ is nonempty and open, and for $x\in C_K
\cap\cl\sigma$ and $y\in D_K\cap\sigma$ the points $(1-t)x+ty$, $t\in(0,1]$, lie in
$D_K\cap\sigma$; so $C_K\cap\cl\sigma=\cl(D_K\cap\sigma)=\{d_\sigma\cdot h\ge0\}\cap
\cl\sigma$, a polyhedral cone. Also $C_K\cap\cl\sigma$ is $\cl\sigma$ for the positive
orthant and $\{0\}$ for the negative one. So the convex cone $C_K$ is a finite union of
polyhedral cones, hence generated by finitely many vectors $g_1,\dots,g_r$
\cite[Cor.~7.1a]{Schrijver}. Since $K$ is closed and convex, Lemma~\ref{lem:D}(c) applied
to points outside $K$ shows $K=\{m\in\Delta:m\cdot h\ge0\text{ for all }h\in C_K\}=
\{m\in\Delta:m\cdot g_t\ge0,\ t=1,\dots,r\}$, a polytope.

(b) Let $V$ be the vertex set of $K$, so $D_K=D_V$, and suppose $\min_Km\cdot g=0$ is
attained at two distinct points; then it is attained at two distinct vertices $a,b$. The
vector $g$ lies in an open orthant $\sigma$, which is mixed, because $m\cdot g$ would
otherwise have constant nonzero sign on $\Delta$. Let $G\coloneqq\{w\in V:w\cdot g=0\}$.
There is $r>0$ with $w\cdot y>0$ for all $w\in V\setminus G$ and $y\in B(g,r)\subseteq
\sigma$. For $z\in\R^n$ and small $t>0$, $g+tz\in B(g,r)$, and $g+tz\in D_K$ iff $w\cdot
z>0$ for every $w\in G$, while $g+tz\in\{d_\sigma\cdot y>0\}$ iff $d_\sigma\cdot z>0$ (note
$d_\sigma\cdot g=0$: $g\in\cl(D_K\cap\sigma)\setminus D_K$). By orthant-planarity,
$\{z:d_\sigma\cdot z>0\}=\{z:w\cdot z>0\ \forall w\in G\}\subseteq\{z:a\cdot z>0\}$. An
open half-space contained in another open half-space through $0$ has a positively
proportional normal, so $a$ and likewise $b$ are positive multiples of $d_\sigma$. Two
distinct points of $\Delta$ are not proportional, a contradiction.

(c) Let $[a,b]$ be an edge, with vertex set $V$ of $K$. Since faces of polytopes are
exposed, there is $g_0$ with $g_0\cdot a=g_0\cdot b=0<g_0\cdot w$ for $w\in V\setminus
\{a,b\}$; then every $g$ in a neighbourhood of $g_0$ in the subspace $H\coloneqq\{a,b\}^
\perp$ has $\min_Km\cdot g=0$ attained at $a$ and $b$. If the line through $a,b$ contains
no vertex of $\Delta$, then $e_l\notin\spn\{a,b\}$ for every $l$, so $H\not\subseteq
e_l^\perp$, and $H\cap e_l^\perp$ is a proper subspace of $H$. A nonempty open subset of
$H$ is not covered by finitely many proper subspaces, so some such $g$ has no zero
coordinate, contradicting (b).

(d) A facet of a full-dimensional polytope in $\aff\Delta$ has dimension $n-2\ge1$, so
contains two points; apply (b) to its normal.
\end{proof}

\subsection{The covering theorem}

\begin{definition}\label{def:chart}
Let $K$ be orthant-planar with chart $c$, and let $U$ be a domain. The \emph{chart ranking}
on $\M_U$ is the single-level integral ranking with $\psi_j(h)\coloneqq-c_j(h)$:
\[
  X\preceq_jY\iff\int_Xc_j\,d\mu\ \ge\ \int_Yc_j\,d\mu .
\]
\end{definition}

\begin{theorem}\label{thm:opcover}
Let $K\subseteq\relint\Delta$ be orthant-planar with chart $c$. Let $U$ be either the full
domain $\R^n$, or $U=D_{K'}$ with $K'\subseteq\relint K$ nonempty, compact and convex (this
requires $K$ to be full-dimensional). Then the chart ranking on $\M_U$ is complete and
transitive, covers $K$, and satisfies (W). Its undominated models are exactly the models
a.e.\ equal to
\[
  U\cap D_{\nu(\pi\odot K)}=
  \begin{cases}
    D_{\nu(\pi\odot K)}&\text{if }U=\R^n,\\
    D_{\conv(K'\cup\nu(\pi\odot K))}&\text{if }U=D_{K'},
  \end{cases}
  \qquad\pi\in\R^n_+\setminus\{0\},
\]
and a model tied at every world with an undominated model equals it a.e.
\end{theorem}

\begin{proof}
The integrands $c_j$ are integrable by (M), since $|c(h)|$ is bounded by a multiple of
$|h|$. (SC) holds because $c_j$ has the sign of $h_j$. For a single-level ranking we take
$W_0=\Omega$; an admissible triple is then $\beta\in\R^n_+\setminus\{0\}$, with
$g=-\beta\cdot c$ and $X^\ast=U\cap\{\beta\cdot c>0\}$. On each open orthant $\sigma$,
$\beta\cdot c(h)=(\beta\odot d_\sigma)\cdot h$ with $\beta\odot d_\sigma\ne0$, so $\{g=0\}$
lies in $Z$ together with finitely many hyperplanes, and is null: the ranking has null
ties. By Lemma~\ref{lem:pullback}, $X^\ast=U\cap D_{\nu(\beta\odot K)}$ a.e. On the full
domain this is coherent. If $U=D_{K'}$, it is $D_{\conv(K'\cup\nu(\beta\odot K))}$ by
Lemma~\ref{lem:D}(a), which is coherent. Theorem~\ref{thm:route} now gives all requirements
except (iv) and (vi) of Definition~\ref{def:cover}, together with (W), the description of
the undominated models, and the statement about ties.

\emph{Target.} With $\beta=\one$, $\nu(\one\odot K)=K$, so $X^\ast=U\cap D_K$. On the full
domain this is $D_K$; if $U=D_{K'}$ with $K'\subseteq K$, then $D_K\subseteq U$ and again
$X^\ast=D_K$. So $D_K$ is undominated.

\emph{Local freedom.} On the full domain (LF) holds with $N=\R^n$. If $U=D_{K'}$ with
$K'\subseteq\relint K$, Lemma~\ref{lem:LF}(a) gives $D_K\in\M_U$ and (LF).
\end{proof}

The undominated family depends on $K$ but not on the arbitrary choices in the chart, while
the comparisons the ranking makes do depend on them. Multiplying $c_j$ by a positive
constant does not change the order at world $j$.

\begin{lemma}\label{lem:charttarget}
Let $K_0\subseteq\relint\Delta$ be orthant-planar with chart $c$, let $U$ be any domain, and
let $K$ be a polytope such that $D_K$ is an undominated model of the chart ranking on
$\M_U$ and (LF) holds at $D_K$. Then $K=\nu(\pi\odot K_0)$ for some $\pi\in\R^n_+\setminus
\{0\}$. If moreover $K\subseteq\relint\Delta$, then $\pi$ has positive coordinates and $K$ is
orthant-planar.
\end{lemma}

\begin{proof}
The chart ranking satisfies (SC) and has null ties on every domain (proof of
Theorem~\ref{thm:opcover}). By Theorem~\ref{thm:route}(c) and Lemma~\ref{lem:pullback},
$D_K=U\cap D_E$ a.e.\ with $E\coloneqq\nu(\pi\odot K_0)$ for some $\pi\in\R^n_+\setminus
\{0\}$. So $D_K\subseteq D_E$ a.e., hence $D_K\subseteq D_E$ by Lemma~\ref{lem:open}, and
$E\subseteq K$ by Lemma~\ref{lem:D}(c). Suppose a vertex $x$ of $K$ lies outside $E$, and
choose $h$ with $\min_{m\in K}m\cdot h=0$ attained only at $x$. Then $h\in\partial D_K
\setminus\{0\}$, and $h\in D_E$ because $m\cdot h>0$ on $K\setminus\{x\}\supseteq E$. By
(LF), as in the proof of Lemma~\ref{lem:LF}(b), some ball around $h$ lies a.e.\ in $U$, and
since $D_E$ is open a smaller ball lies a.e.\ in $U\cap D_E=D_K$. But every ball around
$h$ contains a nonempty open set of gambles $h'$ with $x\cdot h'<0$, which lie outside
$\cl D_K$. So every vertex of $K$ lies in the convex set $E\subseteq K$, and $K=E$. If
$K\subseteq\relint\Delta$, then $\supp\pi=\Omega$, since $\nu(\pi\odot m)$ has the support of
$\pi$; for a mixed orthant $\sigma$ and $h\in\sigma$ we have $\pi\odot h\in\sigma$, so
$h\in D_K$ iff $\pi\odot h\in D_{K_0}$ iff $(d_\sigma\odot\pi)\cdot h>0$, and $K$ is
orthant-planar.
\end{proof}

\begin{corollary}\label{cor:precise}
On the full domain, the Brier ranking $X\preceq_jY\iff\int_Xh_j\,d\mu\ge\int_Yh_j\,d\mu$
covers every precise model $\{p\}$ with $p\in\relint\Delta$, and its undominated models
are exactly the precise models $D_{\{q\}}$, $q\in\Delta$.
\end{corollary}

\begin{proof}
$\{p\}$ is orthant-planar with $d_\sigma=p$ for every $\sigma$, and $c(h)=p\odot h$ gives at
world $j$ the order of $\int_Xp_jh_j\,d\mu$, which is the Brier order. By
Theorem~\ref{thm:opcover} the undominated models are the $D_{\nu(\pi\odot p)}$, $\pi\in
\R^n_+\setminus\{0\}$, and $\nu(\pi\odot p)$ ranges over all of $\Delta$.
\end{proof}

\subsection{Boxes and density-ratio classes}

\begin{proposition}\label{prop:box}
Let $0<l_j\le u_j$ for $j\in\Omega$, $Q\coloneqq\prod_j[l_j,u_j]$, and $K\coloneqq\nu(Q)$, a
\emph{normalised box}. Then $K\subseteq\relint\Delta$ is orthant-planar, with
$d_{\sigma,j}=l_j$ for $j\in P(\sigma)$ and $d_{\sigma,j}=u_j$ for $j\in N(\sigma)$. On $O$ a
chart is given by the concave, continuous, piecewise-linear map
\[
  c_j(h)=\min(l_jh_j,u_jh_j),
\]
and $D_K=\{h:\sum_j\min(l_jh_j,u_jh_j)>0\}$. In particular, for $p\in\relint\Delta$ and
$k\ge1$, the density-ratio class $\{m\in\Delta:m_i/p_i\le k\,m_j/p_j\text{ for all
}i,j\}=\nu(p\odot[1,k]^n)$ is orthant-planar, and $D_K$ is the set of $h$ with
$\sum_jp_j(h_j^+-kh_j^-)>0$. By Theorem~\ref{thm:opcover} all these targets are covered on
the full domain, and the full-dimensional ones on every $D_{K'}$ with $K'\subseteq\relint
K$ compact convex.
\end{proposition}

\begin{proof}
For $m=\nu(q)$ with $q\in Q$, $a\in P(\sigma)$ and $b\in N(\sigma)$ we have $m_b/m_a=q_b/q_a
\le u_b/l_a$, with equality at $d\coloneqq\nu(x)$, $x_j=l_j$ ($j\in P$), $x_j=u_j$ ($j\in
N$), which lies in $K$. Lemma~\ref{lem:ratio}(a) gives orthant-planarity with
$d_\sigma=x$, and $d_{\sigma}\odot h=(\min(l_jh_j,u_jh_j))_j$ on $\sigma$; on the pure
orthants we choose $d_\sigma=l$ (positive) and $u$ (negative), which agree with the same
formula. The formula for $D_K$ is the case $\pi=\one$ of Lemma~\ref{lem:pullback}, or
directly $\min_{q\in Q}q\cdot h=\sum_j\min(l_jh_j,u_jh_j)$. For the density-ratio class,
$m=\nu(p\odot q)$ with $q\in[1,k]^n$ iff the ratios $m_i/p_i$ lie in an interval
$[t,kt]$, which is the stated condition; and $\min(p_jh_j,kp_jh_j)=p_j(h_j^+-kh_j^-)$.
\end{proof}

The last criterion, that $h$ is desirable iff $E_p(h^+)>k\,E_p(h^-)$, is the desirability
criterion of the constant-odds-ratio model with reference $p$ \cite{Walley}. At a point
$\nu(x)$ of a box with $x_a=l_a$ and $x_b=u_b$, the point minimises $m_a/m_b$, so the
two-term vector $x_be_a-x_ae_b$ lies in its normal cone. By Lemma~\ref{lem:twopair} below,
a corner with one coordinate at its lower bound and two coordinates at their upper bounds
is visible from no world, provided $K$ is full-dimensional. Density-ratio classes are not
veto--Brier targets (Corollary~\ref{cor:out}(b)).

\subsection{Three worlds: cevian polygons}

Call a line in $\aff\Delta$ \emph{cevian} if it passes through a vertex of $\Delta$.

\begin{theorem}\label{thm:cevian}
Let $n=3$ and let $K\subseteq\relint\Delta$ be a full-dimensional polygon. Then $K$ is
orthant-planar if and only if every edge of $K$ lies on a cevian line.
\end{theorem}

\begin{proof}
Necessity is Lemma~\ref{lem:planarobs}(c). Conversely suppose every edge line is cevian.
An edge line through $e_l$ has an inward normal $g$ with $g_l=0$; since $g\cdot m=0$ at
points of $K\subseteq\relint\Delta$, the other two coordinates have opposite signs, so the
edge inequality reads $r\,m_i-m_j\ge0$ with $r>0$ and $\{i,j,l\}=\{1,2,3\}$. Since $K$ is
the intersection of its edge half-planes,
\[
  K=\{m\in\relint\Delta:r_fm_{i_f}\ge m_{j_f}\text{ for each edge }f\}.
\]
Let $H\coloneqq\{y\in\R^3:\one\cdot y=0\}$ and $\zeta(m)\coloneqq\log m-\frac13(\one\cdot
\log m)\one$, a homeomorphism of $\relint\Delta$ onto $H$ with inverse $y\mapsto\nu(\exp
y)$. Then $\zeta(K)=\{y\in H:(e_{j_f}-e_{i_f})\cdot y\le\log r_f\}$ is a compact convex
polygon whose outward edge normals are among the six vectors $e_a-e_b$ ($a\ne b$), which
lie in $H$ at angles of $60^\circ$ to their neighbours. Maximising $m_b/m_a$ over $K$ is
maximising the linear function $(e_b-e_a)\cdot y$ over $\zeta(K)$.

Let $\sigma$ be a mixed orthant. By Lemma~\ref{lem:ratio}(a) it suffices to find a point of
$K$ maximising both ratios $m_b/m_a$ with $a\in P(\sigma)$, $b\in N(\sigma)$. There are two
such pairs, and the corresponding vectors $\rho_1,\rho_2$ (namely $e_b-e_a$ and
$e_{b'}-e_a$ if $|P|=1$, or $e_b-e_a$ and $e_b-e_{a'}$ if $|N|=1$) make an angle of
$60^\circ$; so no vector $e_c-e_d$ lies strictly between them. For $f$ in the open cone
$\Gamma$ between $\rho_1$ and $\rho_2$, $f$ is not a positive multiple of an outward edge
normal of $\zeta(K)$, so $f$ is maximised over $\zeta(K)$ at a unique vertex $v(f)$
(otherwise it would be maximised on an edge, and $f$ would be an outward normal of that
edge). The map $f\mapsto v(f)$ is locally constant on $\Gamma$: if $f_k\to f$ and
$v(f_k)=v'$ for infinitely many $k$, then $f$ is maximised at $v'$ by continuity, so
$v'=v(f)$; and there are finitely many vertices. As $\Gamma$ is connected, $v(f)=v$ is
constant, and by continuity $\rho_1$ and $\rho_2$ are both maximised at $v$. The point
$\zeta^{-1}(v)\in K$ maximises both ratios.
\end{proof}

\begin{example}\label{ex:cyclic}
The \emph{cyclic triangle} $K\coloneqq\{m\in\Delta:m_3\le2m_1,\ m_1\le2m_2,\ m_2\le2m_3\}$
has vertices $w_1=\frac17(1,4,2)$, $w_2=\frac17(2,1,4)$ and $w_3=\frac17(4,2,1)$, and its
edge lines $m_3=2m_1$, $m_1=2m_2$, $m_2=2m_3$ pass through $e_2$, $e_3$ and $e_1$. So it is
orthant-planar. Explicitly, the six ratio maxima are $\max m_2/m_1=4$ (at $w_1$), $\max
m_3/m_1=2$ (at $w_1,w_2$), $\max m_1/m_2=2$ (at $w_2,w_3$), $\max m_3/m_2=4$ (at $w_2$),
$\max m_1/m_3=4$ (at $w_3$) and $\max m_2/m_3=2$ (at $w_1,w_3$), so the deciding vertex of
the orthant with $P=\{1\},\{2\},\{3\},\{1,2\},\{1,3\},\{2,3\}$ is $w_1,w_2,w_3,w_2,w_1,w_3$
respectively. The facet normals are $(-1,2,0)$, $(0,-1,2)$ and $(2,0,-1)$. By
Lemma~\ref{lem:poly}, the normal cone at $w_1$ is generated by $(0,-1,2)$ and $(2,0,-1)$:
their first coordinates are non-negative, their second are $-1$ and $0$, and their third
have both signs. So $w_1$ is anti-visible from world $1$, tangential from world $2$ and
crossing for world $3$, and is not visible. By symmetry no vertex of $K$ is
visible. The triangle is not a normalised box: for a box,
\[
  \max\tfrac{m_2}{m_1}\max\tfrac{m_3}{m_2}\max\tfrac{m_1}{m_3}
  =\frac{u_1u_2u_3}{l_1l_2l_3}=\max\tfrac{m_3}{m_1}\max\tfrac{m_2}{m_3}\max\tfrac{m_1}{m_2},
\]
whereas here the two products are $64$ and $8$. By Theorem~\ref{thm:opcover}, $K$ is
covered on the full domain and on every $D_{K'}$ with $K'\subseteq\relint K$; by
Corollary~\ref{cor:polyreach} below, no veto--Brier ranking on a domain $D_{K'}$ has it as
an undominated target with (LF).
\end{example}

\begin{example}\label{ex:hexagon}
The hexagon with vertices $\nu(1,2,2)$, $\nu(1,1,2)$, $\nu(2,1,2)$, $\nu(2,1,1)$,
$\nu(2,\frac{11}5,1)$ and $\nu(\frac{11}{10},\frac{11}5,1)$ has facet normals $(-1,0,2)$,
$(-1,2,0)$, $(0,-5,11)$, $(0,2,-1)$, $(2,-1,0)$ and $(2,0,-1)$, each with a zero
coordinate; so it is orthant-planar by Theorem~\ref{thm:cevian}. The two cyclic products of
maximal ratios are $8$ and $\frac{44}5$, so it is not a normalised box. Three of its
vertices are not visible. (Exact computation.)
\end{example}

\section{Veto--box rankings}\label{sec:vbox}

In this section the Brier level of a veto--Brier ranking is replaced by a box chart. The
single exempt point of Theorem~\ref{thm:exempt} then becomes an exempt normalised box,
which may be degenerate. Throughout, $K'\subseteq\Delta$ is nonempty, compact and convex,
$U=D_{K'}$, $B$ is a veto profile over $K'$ with vetoes $\varphi_i$ and normalised veto
sets $\hat A_i$ as in Definition~\ref{def:vb}, and $R\coloneqq\prod_j[l_j,u_j]$ with
$0<l_j\le u_j$.

\begin{definition}\label{def:vbox}
The \emph{box chart} of $R$ is the map with coordinates $c_j(h)\coloneqq\min(l_jh_j,
u_jh_j)$. The \emph{veto--box ranking} $\mathrm{VBox}(K',B,R)$ on $\M_U$ is the two-level
integral ranking with integrands $\varphi_i$ and $\psi_i\coloneqq-c_i$.
\end{definition}

With $R=\{\one\}$ this is $\VB(K',B)$. With $R=\{r\}$ a point, world $i$ ranks by
$(a_i,r_ib_i)$, which induces the same order as $\VB(K',B)$. With $B_i=\{0\}$ for all
$i$ it is the chart ranking of the normalised box $\nu(R)$ (Proposition~\ref{prop:box}).

\begin{lemma}\label{lem:boxchart}
\begin{enumerate}[label=(\alph*)]
  \item Each $c_j$ is continuous, concave and positively homogeneous, has the sign of
  $h_j$, and satisfies $|c_j(h)|\le u_j|h_j|$.
  \item $c_j(h+t\one)\ge c_j(h)+tl_j$ for $t\ge0$.
  \item For $\beta\in\R^n_+$, $\sum_j\beta_jc_j(h)=\min_{\rho\in R}(\beta\odot\rho)\cdot h$.
\end{enumerate}
\end{lemma}

\begin{proof}
(a) $c_j$ is the minimum of two linear functions of $h_j$ with positive slopes $l_j\le
u_j$. (b) A concave positively homogeneous function of one variable is superadditive, so
$c_j(h_j+t)\ge c_j(h_j)+c_j(t)=c_j(h_j)+l_jt$. (c) For each $j$, the linear function
$\rho_j\mapsto\beta_j\rho_jh_j$ attains its minimum over $[l_j,u_j]$ at $\rho_j=l_j$ if
$h_j\ge0$ and at $\rho_j=u_j$ if $h_j<0$, where it equals $\beta_jc_j(h)$; the choices for
different $j$ are independent.
\end{proof}

For an admissible triple $(W_0,\beta,\alpha)$ the function of Section~\ref{sec:route} is
$g=\sum_{\omega\notin W_0}\alpha_\omega\varphi_\omega-\sum_j\beta_jc_j$, because
$\supp\beta\subseteq W_0$. Put
\[
  W\coloneqq\Bigl\{\beta\odot\rho+\sum_{\omega\notin W_0}\alpha_\omega v_\omega:
  \rho\in R,\ v_\omega\in B_\omega\Bigr\},\qquad
  K(W_0,\beta,\alpha)\coloneqq\conv\Bigl(K'\cup\bigcup_{\omega\in W_0}\hat A_\omega\cup
  \nu(W)\Bigr).
\]

\begin{lemma}\label{lem:vboxcoh}
$\mathrm{VBox}(K',B,R)$ satisfies (SC) and has null ties, and for every admissible triple
$X^\ast(W_0,\beta,\alpha)=D_{K(W_0,\beta,\alpha)}$ a.e., a coherent model.
\end{lemma}

\begin{proof}
The integrands are integrable by Lemmas~\ref{lem:veto}(a), \ref{lem:boxchart}(a) and (M).
(SC) holds by Lemma~\ref{lem:veto}(b) and Lemma~\ref{lem:boxchart}(a).

\emph{Null ties.} $g$ is continuous, and by Lemma~\ref{lem:veto}(c) and
Lemma~\ref{lem:boxchart}(b), $g(h+t\one)\le g(h)-t\,\beta\cdot l<g(h)$ for $t>0$, because
$l$ has positive coordinates and $\beta\ge0$, $\beta\ne0$, so $\beta\cdot l>0$.
Lemma~\ref{lem:fubini} shows that $\{g=0\}$ is null. Only $\beta\cdot l>0$ is used, so
$\beta$ may have zero coordinates.

\emph{Coherence.} By Lemma~\ref{lem:boxchart}(c), $-\sum_j\beta_jc_j(h)=\max_{\rho\in R}
(-(\beta\odot\rho)\cdot h)$, and $\varphi_\omega(h)=\max_{v\in B_\omega}(-v\cdot h)$. Since
$\alpha\ge0$, a sum of maxima over independent choices is the maximum of the sums, so
\[
  g(h)=\max_{w\in W}(-w\cdot h).
\]
$W$ is compact, as a continuous image of a product of compact sets. Every $w\in W$ has
non-negative coordinates (as $\beta\odot\rho\ge0$ and $B_\omega\subseteq\R^n_+$) and
$\one\cdot w\ge\one\cdot(\beta\odot\rho)\ge\beta\cdot l>0$. So $\nu$ is continuous on $W$,
$\nu(W)\subseteq\Delta$ is compact, and, the maximum being attained, $\{g<0\}=D_{\nu(W)}$.
As in the proof of Lemma~\ref{lem:coh}, $\{\varphi_\omega=0\}=D_{\hat A_\omega}$ a.e.\ for
$\omega\in W_0$, so by Lemma~\ref{lem:D}(a)
\[
  X^\ast(W_0,\beta,\alpha)=D_{K'}\cap\bigcap_{\omega\in W_0}D_{\hat A_\omega}\cap
  D_{\nu(W)}=D_{K(W_0,\beta,\alpha)}\quad\text{a.e.},
\]
and $K(W_0,\beta,\alpha)$ is a credal set containing $K'$.
\end{proof}

The coherence step uses only that the second level, weighted and added to the weighted
vetoes, is a maximum of linear functionals $-w\cdot h$ over a compact set $W\subseteq\R^n_+$
with $\one\cdot w>0$ on $W$. Compactness makes the maximum attained and $\nu(W)$ compact.

\begin{theorem}\label{thm:BX}
Let $K'\subseteq\Delta$ be nonempty, compact and convex, let $A_1,\dots,A_n$ be compact
(possibly empty) with $A_i\subseteq\conv(K'\cup\{e_i\})$, put $B_i\coloneqq\conv(\{0\}
\cup A_i)$, let $R$ be a box as above, and let $\beta\in\R^n_+\setminus\{0\}$. Put
\[
  K\coloneqq\conv\bigl(K'\cup A_1\cup\dots\cup A_n\cup\nu(\beta\odot R)\bigr).
\]
If $K$ is full-dimensional and $K'\subseteq\relint K$, then $\mathrm{VBox}(K',B,R)$ on
$\M_{D_{K'}}$ is complete and transitive, covers $K$, and satisfies (W). Its undominated
models are exactly the models a.e.\ equal to $D_{K(W_0,\beta',\alpha)}$ for admissible
triples $(W_0,\beta',\alpha)$, and a model tied at every world with an undominated model
equals it a.e.
\end{theorem}

\begin{proof}
By Lemma~\ref{lem:vboxcoh} and Theorem~\ref{thm:route}, all requirements hold except (iv)
and (vi) of Definition~\ref{def:cover}, together with the remaining claims. The triple
$(\Omega,\beta,\emptyset)$ is admissible, with $W=\beta\odot R$; and $\hat A_\omega=\conv
A_\omega$ (Section~\ref{sec:vb}), so $K(\Omega,\beta,\emptyset)=K$, and $D_K$ is
undominated by Theorem~\ref{thm:route}(c). Lemma~\ref{lem:LF}(a) gives $D_K\in\M_U$ and
(LF).
\end{proof}

If $R=\{r\}$ is a point, then $\nu(\beta\odot R)$ is the single point $\nu(\beta\odot r)$,
which can be any point of $\Delta$; so Theorem~\ref{thm:BX} contains
Theorem~\ref{thm:exempt}. With $A=\emptyset$ and $\beta=\one$ it gives the box coverage of
Proposition~\ref{prop:box} on one-sided domains. The exempt set $\nu(\beta\odot R)$ is a
normalised box: if $\beta$ has positive coordinates, $\beta\odot R$ is again a box, and
$l_j=u_j$ is allowed. If $\beta$ has zero coordinates, $\nu(\beta\odot R)$ is a normalised
box of the face $\Delta_{\supp\beta}$, and the vertex $e_j$ when $\supp\beta=\{j\}$.

\begin{lemma}\label{lem:boxgeom}
Let $R'=\prod_j[l'_j,u'_j]$ with $0<l'_j\le u'_j$.
\begin{enumerate}[label=(\alph*)]
  \item $\nu(R')=\{m\in\Delta:l'_im_j\le u'_jm_i\text{ for all }i\ne j\}$.
  \item Every segment $[a,b]\subseteq\relint\Delta$ with $a\ne b$ lying on a cevian line
  is $\nu(R')$ for a box $R'$ with exactly one nondegenerate coordinate.
  \item Let $n=3$. If $\nu(R')$ has dimension $2$, each of its edges lies on a cevian line;
  if it has dimension $1$, it lies on a cevian line; and it is never a triangle.
\end{enumerate}
\end{lemma}

\begin{proof}
(a) If $m=\nu(y)$ with $y\in R'$, then $m_j/m_i=y_j/y_i\le u'_j/l'_i$. Conversely, if $m$
satisfies the inequalities, then $m$ has positive coordinates (if $m_i=0$ then $l'_im_j\le
0$ for all $j$, so $m=0$), and $\max_il'_i/m_i\le\min_ju'_j/m_j$; for $t$ between these
numbers, $tm\in R'$, so $m=\nu(tm)\in\nu(R')$.

(b) Let the line pass through $e_l$. After exchanging $a$ and $b$ if necessary,
$b=(1-\tau)a+\tau e_l$ with $\tau\in(0,1)$: indeed $b=a+s(e_l-a)$ with $s\ne0$, and $s<1$
because $b\in\relint\Delta$; if $s<0$ then $a=b+s'(e_l-b)$ with $s'=-s/(1-s)\in(0,1)$.
Put $R'_l\coloneqq[a_l,a_l+\tau/(1-\tau)]$ and $R'_i\coloneqq\{a_i\}$ for $i\ne l$. For
$t\ge0$, $\nu(a+te_l)=(1-\tau')a+\tau'e_l$ with $\tau'=t/(1+t)$, which runs over $[0,\tau]$
as $t$ runs over $[0,\tau/(1-\tau)]$. So $\nu(R')=[a,b]$.

(c) By (a), $\nu(R')$ is the intersection of $\relint\Delta$ with the six half-planes
$l'_im_j\le u'_jm_i$, each bounded by a cevian line (through the third vertex). If
$\nu(R')$ has dimension $2$, each of its edges lies on one of these lines. If exactly two
coordinates are nondegenerate, say $l'_3=u'_3$, then the ratios $m_1/m_3$ and $m_2/m_3$
range independently over intervals with nonempty interior, so $\nu(R')$ has dimension $2$;
if all three are nondegenerate, it contains such a set. So dimension $1$ means that
exactly one coordinate $l$ is nondegenerate, and then $\nu(R')$ lies on the line through
$e_l$ and $\nu(y)$ for any fixed $y\in R'$. Finally suppose $\nu(R')$ is a triangle. Apply
the homeomorphism $\zeta$ of the proof of Theorem~\ref{thm:cevian}; it maps cevian lines
to lines, so $\zeta(\nu(R'))=\{y\in H:(e_j-e_i)\cdot y\le a_{ij}\}$ with $a_{ij}\coloneqq
\log(u'_j/l'_i)$ is a triangle. Its three outward edge normals are among the vectors
$e_j-e_i$, contain no opposite pair, and positively span $H$; so they are $e_2-e_1$,
$e_3-e_2$, $e_1-e_3$ or their negatives. In the first case the vertex on the edges with
normals $e_2-e_1$ and $e_3-e_2$ has $y_3-y_1=a_{12}+a_{23}$, and it satisfies $y_3-y_1\le
a_{13}$; this gives $\log(u'_2/l'_2)\le0$, so $l'_2=u'_2$. The other two vertices give
$l'_3=u'_3$ and $l'_1=u'_1$ in the same way, so $\nu(R')$ is a point. The second case is
symmetric. So $\nu(R')$ is not a triangle.
\end{proof}

\begin{corollary}\label{cor:BX3}
Let $n=3$ and let $T\subseteq\relint\Delta$ be a triangle. Call a vertex \emph{bad} if it
is not visible. The following are equivalent.
\begin{enumerate}[label=(\roman*)]
  \item For some nonempty compact convex $K'\subseteq\Delta$, veto profile $B$ over $K'$
  and box $R$, $D_T$ is an undominated model of $\mathrm{VBox}(K',B,R)$ on
  $\M_{D_{K'}}$ and (LF) holds at $D_T$.
  \item $T$ is covered by some veto--box ranking on some $\M_{D_{K'}}$.
  \item Either (V) $T$ has at most one bad vertex, or (X) $T$ has exactly two bad vertices
  and the edge joining them lies on a cevian line.
\end{enumerate}
\end{corollary}

\begin{proof}
(ii)$\Rightarrow$(i) is immediate.

(iii)$\Rightarrow$(ii). In case (V), Corollary~\ref{cor:poly} gives a veto--Brier ranking
covering $T$, and veto--Brier rankings are veto--box rankings with $R=\{\one\}$. In case
(X), let $a,b$ be the bad vertices and $c$ the third, visible from some world $j$. Since
$c\ne e_j$, there is $\varepsilon>0$ with $k\coloneqq c+\varepsilon(c-e_j)\in\relint T$.
Put $K'\coloneqq\{k\}$, $A_j\coloneqq\{c\}$ (note $c=(k+\varepsilon e_j)/(1+\varepsilon)
\in\conv(K'\cup\{e_j\})$), $A_i\coloneqq\emptyset$ for $i\ne j$, $\beta\coloneqq\one$, and
let $R$ be the box of Lemma~\ref{lem:boxgeom}(b) with $\nu(R)=[a,b]$. Then $\conv(\{k,c\}
\cup[a,b])=T$, which is full-dimensional with $K'\subseteq\relint T$, and
Theorem~\ref{thm:BX} applies.

(i)$\Rightarrow$(iii). By Lemma~\ref{lem:LF}(b), $K'\subseteq\relint T$. By
Lemma~\ref{lem:vboxcoh}, Theorem~\ref{thm:route}(c) and Lemma~\ref{lem:open},
$T=K(W_0,\beta,\alpha)$ for an admissible triple. Taking $v_\omega=0$ in the definition of $W$, $\nu(\beta\odot\rho)\in\nu(W)
\subseteq T\subseteq\relint\Delta$ for $\rho\in R$, so $\supp\beta=\Omega$; hence
$W_0=\Omega$, there are no $\alpha$-weights, $W=\beta\odot R$, and $T=\conv(K'\cup
\bigcup_\omega\hat A_\omega\cup E)$ with $E\coloneqq\nu(\beta\odot R)$, a normalised box
contained in $T$. Every vertex of $T$ lies in $K'\cup\bigcup_\omega\hat A_\omega\cup E$.
It is not in $K'\subseteq\relint T$. A vertex $x\in\hat A_\omega\subseteq\conv(K'\cup
\{e_\omega\})$ is $(1-t)k+te_\omega$ with $k\in K'$; here $t<1$, because $x\in\relint\Delta$, and then
$x$ is visible from $\omega$ by
Lemma~\ref{lem:visbasic}(a). So every bad vertex lies in $E$. If all three vertices were
bad, then $T=E$, which is impossible by Lemma~\ref{lem:boxgeom}(c). Suppose exactly two
vertices $a,b$ are bad, and let $\ell$ be the line of the edge $[a,b]$, which supports $T$.
Since $E\subseteq T$ lies on one side of $\ell$ and contains $a$ and $b$, $E\cap\ell$ is a
face of $E$ containing two points. If $E$ has dimension $2$, this face is an edge of $E$
and $\ell$ is cevian by Lemma~\ref{lem:boxgeom}(c); if $E$ has dimension $1$, then
$E\subseteq\ell$ and $\ell$ is cevian by the same lemma. So $T$ satisfies (X).
\end{proof}

In (iii), both bad vertices lie in the exempt box; this is how a single veto--box ranking
exempts two corners at once.

\begin{example}\label{ex:T}
Let $n=3$, $p\coloneqq\frac15(1,2,2)$, $q\coloneqq\frac15(2,1,2)$, $s\coloneqq\frac15
(3,1,1)$ and $T\coloneqq\conv\{p,q,s\}$. Its facet normals are $(2,2,-3)$ for the edge
$[p,q]$ (the line $m_3=\frac25$), $(0,-1,1)$ for $[p,s]$ (the line $m_2=m_3$, through
$e_1$) and $(-1,4,-1)$ for $[q,s]$ (the line $m_2=\frac15$). By Lemma~\ref{lem:poly}:
\begin{itemize}
  \item at $p$ the normals are $(2,2,-3)$ and $(0,-1,1)$, so $p$ is anti-visible from
  world $1$ and worlds $2$ and $3$ are crossing; $p$ is not visible;
  \item at $s$ the normals are $(0,-1,1)$ and $(-1,4,-1)$, so $s$ is tangential from world
  $1$ and worlds $2$ and $3$ are crossing; $s$ is not visible;
  \item at $q$ the normals are $(2,2,-3)$ and $(-1,4,-1)$, whose third coordinates are
  negative, so $q$ is visible from world $3$.
\end{itemize}
So $T$ has exactly two bad vertices, joined by the cevian edge $[p,s]$, and $T$ satisfies
(X). It has a non-cevian edge, so it is not orthant-planar (Theorem~\ref{thm:cevian}), and
it is not a veto--Brier target (Corollary~\ref{cor:polyreach}).

An explicit cover is as follows. Let $k\coloneqq\frac1{50}(22,11,17)=\frac1{10}p+\frac35q+
\frac3{10}s\in\relint T$. Then $q=\frac{10}{11}k+\frac1{11}e_3$. Take $K'\coloneqq\{k\}$,
$A_3\coloneqq\{q\}$, $A_1=A_2\coloneqq\emptyset$, $R\coloneqq[1,6]\times\{2\}\times\{2\}$
and $\beta\coloneqq\one$. Then $\nu(R)=[\nu(1,2,2),\nu(6,2,2)]=[p,s]$, and
$\conv(\{k\}\cup\{q\}\cup[p,s])=T$, so Theorem~\ref{thm:BX} shows that the following
ranking covers $T$, on the domain $U=D_{\{k\}}=\{h:22h_1+11h_2+17h_3>0\}$:
\begin{center}
\begin{tabular}{lll}
\toprule
world & first level & second level\\
\midrule
$1$ & --- & $-\int_X\min(h_1,6h_1)\,d\mu=\int_X(6h_1^--h_1^+)\,d\mu$\\
$2$ & --- & $-\int_X2h_2\,d\mu$\\
$3$ & $\int_X\max(0,-q\cdot h)\,d\mu$ & $-\int_X2h_3\,d\mu$\\
\bottomrule
\end{tabular}
\end{center}
The factor $2$ at worlds $2$ and $3$ does not change their orders. With all weights equal,
the second levels add up to $-\min(5\,p\cdot h,10\,s\cdot h)$, since $(6,2,2)\cdot h-
(1,2,2)\cdot h=5h_1$; on the veto-free set $U\cap\{q\cdot h\ge0\}$ this selects exactly
$D_T$. The ranking is complete and transitive and satisfies (W). Finally, $T$ is not of the
form $L\star K_0$ with (L) at all: since $\nu(\pi^\ast\odot m)\in T\subseteq\relint\Delta$
has the support of $\pi^\ast$, such a factorisation would have $\pi^\ast\in\relint\Delta$,
and then $T$ would have at most one bad vertex (Theorem~\ref{thm:n3} below).
\end{example}

\section{Composition}\label{sec:comp}

\subsection{Composed rankings}

Throughout this section $K_0\subseteq\relint\Delta$ is an orthant-planar credal set, of any
dimension, with chart $c$, and the following holds.
\begin{itemize}
  \item[(L)] $L'\subseteq\Delta$ is nonempty, compact and convex; $A_1,\dots,A_n$ are
  compact with $A_i\subseteq\conv(L'\cup\{e_i\})$; $\pi^\ast\in\Delta$; and
  $L\coloneqq\conv(L'\cup A_1\cup\dots\cup A_n\cup\{\pi^\ast\})$ is full-dimensional with
  $L'\subseteq\relint L$.
\end{itemize}
Taking $\pi^\ast\in L'$ recovers the case without an exempt point. We write
$K(W_0,\beta,\alpha)$ for the credal sets of Lemma~\ref{lem:coh} for $\VB(L',A)$, and put
\[
  K\coloneqq L\star K_0,\qquad K'\coloneqq L'\star K_0,\qquad U\coloneqq D_{K'}.
\]

\begin{definition}\label{def:composed}
The \emph{composed ranking} on $\M_U$ is the two-level integral ranking with
\[
  \varphi_i(h)\coloneqq\max\Bigl(0,\max_{v\in A_i}(-v\cdot c(h))\Bigr),\qquad
  \psi_i(h)\coloneqq-c_i(h).
\]
\end{definition}

Taking $K_0=\{\one/n\}$ and $c=$ identity recovers $\VB(L',A)$. The exempt point does not
enter the ranking.

\begin{lemma}\label{lem:conj}
Let $\mu_c\coloneqq c_\sharp\mu$, that is, $\mu_c(E)\coloneqq\mu(c^{-1}(E))$.
\begin{enumerate}[label=(\alph*)]
  \item $\mu_c$ satisfies (M).
  \item $U\cap O=c^{-1}(D_{L'})\cap O$, and $\Phi(X)\coloneqq c(X\cap O)$ is a bijection
  from $\M_U$ (under $\mu$) onto $\M_{D_{L'}}$ (under $\mu_c$), modulo null sets, which
  commutes with unions, intersections and differences.
  \item For each $i$ and $X,Y\in\M_U$: $X\preceq_iY$ in the composed ranking iff
  $\Phi(X)\preceq_i\Phi(Y)$ in $\VB(L',A)$ under $\mu_c$. Moreover $X\subseteq\{h_i>0\}$
  a.e.\ iff $\Phi(X)\subseteq\{x_i>0\}$ a.e., and likewise for $\{h_i<0\}$.
\end{enumerate}
\end{lemma}

\begin{proof}
(a) $\mu_c$ is a finite Borel measure, and $\int|x|\,d\mu_c=\int|c(h)|\,d\mu<\infty$. Let
$E$ be Borel. Since $c$ maps each open orthant $\sigma$ onto itself by the invertible
linear map $\operatorname{diag}(d_\sigma)$, and maps $Z$ into $Z$ (as the identity),
$c^{-1}(E)$ is the union of the sets $\operatorname{diag}(d_\sigma)^{-1}(E\cap\sigma)$ and
of $E\cap Z$. If $E$ is Lebesgue-null, so is each of these, and $\mu_c(E)=0$. If
$\mu_c(E)=0$, each set $\operatorname{diag}(d_\sigma)^{-1}(E\cap\sigma)$ is null, so
$E\cap\sigma$ is null, and $E\cap Z$ is null; hence $E$ is null.

(b) The first claim is Lemma~\ref{lem:pullback} with $Q=L'$. Since $c$ is a bijection of $O$
onto itself with measurable inverse, $\Phi$ maps measurable subsets of $U\cap O$
bijectively onto measurable subsets of $D_{L'}\cap O$ and commutes with set operations;
$Z$ is null for both measures; and $\mu_c(\Phi(X))=\mu(X\cap O)=\mu(X)$, so null sets
correspond.

(c) For integrable $f$, $\int_Xf(c(h))\,d\mu(h)=\int_{\Phi(X)}f\,d\mu_c$, by the definition
of the image measure and $1_X=1_{\Phi(X)}\circ c$ on $O$. With
$f(x)=\max(0,\max_{v\in A_i}(-v\cdot x))$ and $f(x)=-x_i$ this identifies the scores of $X$
in the composed ranking with the scores of $\Phi(X)$ in $\VB(L',A)$. The last claim holds
because $c_i(h)$ has the sign of $h_i$.
\end{proof}

\begin{theorem}\label{thm:comp}
Assume (L). Then $K=L\star K_0$ is a full-dimensional credal set, $K'=L'\star K_0$ is a
compact convex subset of $\relint K$, and the composed ranking on $\M_{D_{K'}}$ is complete
and transitive, covers $K$, and satisfies (W). Its undominated models are exactly the
models a.e.\ equal to $D_{K(W_0,\beta,\alpha)\star K_0}$ for admissible triples
$(W_0,\beta,\alpha)$, and a model tied at every world with an undominated model equals it
a.e.
\end{theorem}

\begin{proof}
\emph{Geometry.} For $m\in K_0$ the map $T_m:q\mapsto\nu(q\odot m)$ is a homeomorphism of
$\relint\Delta$ onto itself. Since $\relint L$ is open in $\aff\Delta$ and contained in
$\relint\Delta$, $T_m(\relint L)$ is open in $\aff\Delta$ and contained in $K$; so $K$ is
full-dimensional and $T_m(\relint L)\subseteq\relint K$. Hence every generator $\nu(q\odot
m)$ of $K'$ ($q\in L'\subseteq\relint L$) lies in $\relint K$, and so does their convex
hull $K'$, which is compact.

\emph{Ranking.} By Lemma~\ref{lem:conj}, $\Phi$ is an isomorphism between the composed
ranking on $\M_U$ under $\mu$ and $\VB(L',A)$ on $\M_{D_{L'}}$ under $\mu_c$, which
preserves set operations, null sets and the correctness of acceptances at every world.
So (G), (IB), completeness, transitivity, the dominance relations, undominatedness and ties
correspond. Theorem~\ref{thm:family} applies to $\VB(L',A)$ under $\mu_c$. Its undominated
models are the $D_{K(W_0,\beta,\alpha)}$, and every model is weakly dominated by one of
them. Since $K(W_0,\beta,\alpha)\supseteq L'$, Lemma~\ref{lem:pullback} gives
$\Phi^{-1}(D_{K(W_0,\beta,\alpha)})=D_{K(W_0,\beta,\alpha)\star K_0}$ a.e., a coherent
model. Theorem~\ref{thm:route}(e), applied through $\Phi$, gives (GC), (C), (ID) and (W).

\emph{Target.} By Theorem~\ref{thm:exempt}, $D_L$ is undominated in $\VB(L',A)$ under
$\mu_c$, and $L=K(\Omega,\pi^\ast,\emptyset)$. So $\Phi^{-1}(D_L)=D_{L\star K_0}=D_K$ is
undominated.

\emph{Local freedom.} $K'\subseteq\relint K$ and Lemma~\ref{lem:LF}(a) give $D_K\in\M_U$
and (LF).
\end{proof}

\begin{lemma}\label{lem:reach}
Let $Q\subseteq\Delta$ be compact and convex. Every extreme point of $Q\star K_0$ has the
form $\nu(q\odot m)$ with $q\in\ext Q$ and $m\in\ext K_0$.
\end{lemma}

\begin{proof}
An extreme point $x$ of the convex hull of the compact generating set is a generator,
$x=\nu(q\odot m)$. Write $m=\sum_t\lambda_tm_t$ as a convex combination of extreme points
of $K_0$ with positive weights. By~\eqref{eq:proj}, $x$ is a convex combination with
positive weights of the points $\nu(q\odot m_t)\in Q\star K_0$, so $x=\nu(q\odot m_t)$ for
each $t$. Repeating the argument in $q$ with $m_t$ fixed gives the claim.
\end{proof}

\begin{lemma}\label{lem:U}
Assume (L) and let $\delta>0$ be such that $K'\subseteq K_\delta$. Every extreme point of
$K$ that does not lie in $\nu(\pi^\ast\odot K_0)$ lies in $\bigcup_j\conv(K_\delta\cup
\{e_j\})$. The extreme points of $K$ that lie in $\nu(\pi^\ast\odot K_0)$ are among the
finitely many points $\nu(\pi^\ast\odot w)$, $w\in\ext K_0$.
\end{lemma}

\begin{proof}
Such $\delta$ exists by Lemma~\ref{lem:visbasic}(b), since $K'$ is a compact subset of
$\relint K$. Let $x$ be extreme in $K$. By Lemma~\ref{lem:reach}, $x=\nu(q\odot m)$ with
$q\in\ext L\subseteq L'\cup\bigcup_jA_j\cup\{\pi^\ast\}$. If $q\in L'$, then $x\in
K'\subseteq\relint K$, which is impossible. If $q\in A_j$, write $q=(1-t)k+te_j$ with $k\in
L'$, $t\in[0,1]$; then $t>0$, and by~\eqref{eq:proj} (applied in $q$ with $m$ fixed), $x$
lies on the segment from $\nu(k\odot m)\in K'\subseteq K_\delta$ to $\nu(e_j\odot m)=e_j$.
Otherwise $q=\pi^\ast$. For the second claim, if $x=\nu(\pi^\ast\odot m)$ is extreme in $K$,
the argument of Lemma~\ref{lem:reach} shows $x=\nu(\pi^\ast\odot w)$ for an extreme point
$w$ of $K_0$, and $K_0$ has finitely many by Lemma~\ref{lem:planarobs}(a).
\end{proof}

\subsection{Three worlds: no new interior targets}

For a credal set $S$, a point $x\in S$ and $i\ne j$, say that $x$ \emph{minimises $m_i/m_j$
over $S$} if $m_ix_j-m_jx_i\ge0$ for all $m\in S$, and put
\[
  M_i(S)\coloneqq\{x\in S\cap\relint\Delta:x\text{ minimises }m_i/m_j\text{ over }S\text{
  for every }j\ne i\}.
\]
Equivalently, the two-term vector $g^{ij}_x\coloneqq x_je_i-x_ie_j$ lies in $F^S_x$ for
every $j\ne i$: indeed $g^{ij}_x\cdot x=0$ always, and $g^{ij}_x\in C_S$ is the stated
inequality.

\begin{lemma}\label{lem:N3}
Let $n=3$, let $S$ be a full-dimensional credal set and $x\in S\cap\relint\Delta$. Then $x$
is visible from no world if and only if $x\in M_i(S)$ for some $i$. Each $M_i(S)$ has at
most one point, which is an extreme point of $S$, and $M_i(S)\cap M_j(S)=\emptyset$ for
$i\ne j$.
\end{lemma}

\begin{proof}
Let $\Pi\coloneqq x^\perp$, a plane, and $F\coloneqq F^S_x\subseteq\Pi$, a closed convex
cone with $F\cap(-F)=\{0\}$ (Lemma~\ref{lem:fulldim}(c)). By Lemma~\ref{lem:V1}, $x$ is
visible from no world iff $F$ meets $H_l\coloneqq\{h_l\ge0\}\setminus\{0\}$ for every
$l$. Since $x$ has positive coordinates, no nonzero $h\in\Pi$ has $h\ge0$ or $h\le0$. For
each $l$, $\Pi\cap e_l^\perp=\R g^{ij}_x$ with $\{i,j,l\}=\{1,2,3\}$, and $g^{ij}_x$ is
positive at $i$, negative at $j$ and zero at $l$. So the nonzero points of $\Pi$ with a
zero coordinate form six open rays, $\pm\R_+g^{ij}_x$, which divide $\Pi\setminus\{0\}$ into
six open sectors, each of angle less than $\pi$. On each sector every coordinate has a
constant nonzero sign, not all equal; crossing a ray on which $h_l=0$ changes the sign of
$h_l$ only. So the number of positive coordinates is $1$ or $2$ on each sector, and the
two kinds alternate. The closure of a sector with one positive coordinate $i$ is
$\cone\{g^{ij}_x,g^{ik}_x\}$.

($\Leftarrow$) If $g^{ij}_x,g^{ik}_x\in F$, then $g^{ij}_x$ meets $H_i$ and $H_k$, and
$g^{ik}_x$ meets $H_j$.

($\Rightarrow$) $F\neq\{0\}$, since otherwise $F$ meets no $H_l$. If $F$ contained none of
the six rays, $F\setminus\{0\}$ would be a connected set on which no coordinate vanishes,
so each coordinate would have constant sign there, some coordinate $l$ negative, and $F$
would miss $H_l$. If $F$ contained exactly one ray, on which $h_l=0$, $h_i>0>h_j$, then
$F\setminus\{0\}$ would lie in that ray and the two adjacent sectors, on all of which
$h_j<0$; so $F$ would miss $H_j$. Hence $F$ contains two rays, and, being convex with
angle less than $\pi$, the closed sector between two consecutive ones. If this sector has
one positive coordinate $i$, it is $\cone\{g^{ij}_x,g^{ik}_x\}\subseteq F$, and we are
done. Otherwise it has two positive coordinates and one negative coordinate $l$, which is
negative on both bounding rays; if $F$ contained no further ray, $F\setminus\{0\}$ would
lie in this closed sector and the two adjacent open sectors, on all of which $h_l<0$, and
$F$ would miss $H_l$. So $F$ contains a third ray adjacent to this sector, and hence the
closed sector beyond it, which has one positive coordinate.

The points of $S$ minimising $m_i/m_j$ form an exposed face of $S$ lying on a line through
$e_k$, and those minimising $m_i/m_k$ an exposed face on a line through $e_j$; the lines
are distinct, since a line through $e_j$ and $e_k$ is the edge $\{m_i=0\}$, which misses
$x$. So $M_i(S)$ is contained in a face of $S$ with at most one point, which is then an
extreme point. If $x\in M_i(S)\cap M_j(S)$, then $x$ minimises both $m_i/m_j$ and $m_j/m_i$,
so $m_i/m_j$ is constant on $S$ and $S$ lies on a line, contradicting
full-dimensionality.
\end{proof}

\begin{lemma}\label{lem:decouple}
Let $S\subseteq\Delta$ be compact and convex, $q\in S\cap\relint\Delta$, $w\in K_0$, and
$x\coloneqq\nu(q\odot w)$. For $i\ne j$, $x$ minimises $m_i/m_j$ over $S\star K_0$ if and
only if $q$ minimises $m_i/m_j$ over $S$ and $w$ minimises $m_i/m_j$ over $K_0$.
Consequently $x\in M_i(S\star K_0)$ iff $q\in M_i(S)$ and $w\in M_i(K_0)$.
\end{lemma}

\begin{proof}
The inequality $m_ix_j-m_jx_i\ge0$ is linear in $m$, so it holds on $S\star K_0$ iff it
holds at the generators $\nu(q'\odot m')$, that is, iff $q'_im'_iq_jw_j-q'_jm'_jq_iw_i\ge0$
for all $q'\in S$, $m'\in K_0$. If $q'_iq_j\ge q'_jq_i$ and $m'_iw_j\ge m'_jw_i$ for all
$q',m'$, the products of these non-negative quantities satisfy the inequality. Conversely,
$m'=w$ gives $(q'_iq_j-q'_jq_i)w_iw_j\ge0$, and $q'=q$ gives $(m'_iw_j-m'_jw_i)q_iq_j\ge0$;
and $q,w$ have positive coordinates.
\end{proof}

For a full-dimensional credal set $S$ let $\mathsf N(S)$ be the number of extreme points of
$S$ in $\relint\Delta$ that are not visible.

\begin{theorem}\label{thm:n3}
Let $n=3$, assume (L) with $\pi^\ast\in\relint\Delta$, and let $K_0\subseteq\relint\Delta$
be orthant-planar. Then $\mathsf N(L\star K_0)=\mathsf N(L)\le1$, and every extreme point of
$L\star K_0$, with at most one exception, is uniformly visible. Consequently $L\star K_0$ is
covered by a veto--Brier ranking (Corollary~\ref{cor:exemptvis}): in three worlds,
composition with an interior exempt point produces no target outside the class of
Section~\ref{sec:vb}.
\end{theorem}

\begin{proof}
Write $K=L\star K_0$. \emph{The factor $L$.} Every extreme point of $L$ other than
$\pi^\ast$ lies in $L'\cup\bigcup_jA_j$; points of $L'$ are not extreme, points of $A_j$
other than $e_j$ have the form $(1-t)k+te_j$ with $k\in L'\subseteq\relint L$, $t<1$, so are
visible by Lemma~\ref{lem:visbasic}(a), and $e_j\notin\relint\Delta$. So $\mathsf N(L)\le
1$, with equality iff $\pi^\ast$ is not visible, iff (Lemma~\ref{lem:N3}) $\pi^\ast\in
M_i(L)$ for some $i$, which is then unique.

\emph{The factor $K_0$.} For each $i$, $M_i(K_0)$ is a single point $w_i$: apply
Lemma~\ref{lem:ratio}(b) to the mixed orthant with $P=\{i\}$, $N=\Omega\setminus\{i\}$,
which gives a point of $K_0$ maximising $m_j/m_i$ for both $j\ne i$; and $|M_i(K_0)|\le1$
by the argument in the proof of Lemma~\ref{lem:N3} (two distinct lines), which does not use
full-dimensionality of $K_0$.

\emph{The composite.} By Lemma~\ref{lem:U}, extreme points of $K$ outside
$\nu(\pi^\ast\odot K_0)$ are uniformly visible, hence visible. Let $x$ be a non-visible
extreme point of $K$ in $\relint\Delta$; then $x=\nu(\pi^\ast\odot w)$ with $w\in K_0$, and
by Lemma~\ref{lem:N3}, $x\in M_i(K)$ for some $i$. By Lemma~\ref{lem:decouple},
$\pi^\ast\in M_i(L)$ and $w=w_i$. So $\mathsf N(K)\le\mathsf N(L)$, and $x$ is determined.
Conversely, if $\pi^\ast\in M_i(L)$, Lemma~\ref{lem:decouple} gives $\nu(\pi^\ast\odot w_i)
\in M_i(K)$, which by Lemma~\ref{lem:N3} is an extreme point of $K$ in $\relint\Delta$ that is
not visible. So $\mathsf N(K)=\mathsf N(L)$.

\emph{Uniformity.} The remaining extreme points are those of Lemma~\ref{lem:U} outside
$\nu(\pi^\ast\odot K_0)$, which lie in $\bigcup_j\conv(K_\delta\cup\{e_j\})$, and finitely
many visible points $\nu(\pi^\ast\odot w)$, each of which is $(1-t)k+te_j$ with
$k\in\relint K$, $t<1$, hence lies in $\conv(K_{d(k)}\cup\{e_j\})$. Taking the least depth,
all extreme points except the non-visible one (if any) are uniformly visible.
\end{proof}

\begin{corollary}\label{cor:trireach}
Let $n=3$ and let $T\subseteq\relint\Delta$ be a triangle. Then $T$ is covered by one of
the constructions of this paper, namely a veto--Brier ranking (Section~\ref{sec:vb}), a
chart ranking of an orthant-planar credal set on any domain (Definition~\ref{def:chart}), a
veto--box
ranking (Theorem~\ref{thm:BX}), or a composed ranking with a target of the form
$L\star K_0$ of Theorem~\ref{thm:comp}, if and only if $T$ satisfies (V) or (X) of
Corollary~\ref{cor:BX3}, or all three edges of $T$ lie on cevian lines.
\end{corollary}

\begin{proof}
The constructions give the three classes by Corollary~\ref{cor:BX3} and
Theorem~\ref{thm:cevian}. Conversely, a veto--Brier ranking is a veto--box ranking, so
Corollary~\ref{cor:BX3} applies to both. If $D_T$ is an undominated model of a chart
ranking of an orthant-planar credal set on any domain, with (LF) at $D_T$, then $T$ is
orthant-planar by Lemma~\ref{lem:charttarget}, and all its edges are cevian by
Theorem~\ref{thm:cevian}. Finally, if $T=L\star K_0$ with (L),
then $\pi^\ast\in\relint\Delta$, since $\nu(\pi^\ast\odot m)\in T$ has the support of
$\pi^\ast$; by Theorem~\ref{thm:n3}, $T$ has at most one bad vertex.
\end{proof}

\begin{remark}\label{rem:n3}
The mechanism is specific to three worlds: there the normal cone of a non-visible interior
point must contain a whole sector bounded by two two-term vectors (Lemma~\ref{lem:N3}), and
these two-term conditions decouple across the factors (Lemma~\ref{lem:decouple}). In four
or more worlds the two-term condition is still sufficient for non-visibility
(Lemma~\ref{lem:twopair} below); we do not know whether it is necessary there, and
composition does produce new targets (Theorem~\ref{thm:F}).
\end{remark}

\subsection{Four or more worlds}

\begin{lemma}\label{lem:twopair}
Let $S$ be a full-dimensional credal set and $x\in S\cap\relint\Delta$. If there are pairs
$(a,b)$ and $(a',b')$ with $b\ne b'$ such that $x$ minimises $m_a/m_b$ and $m_{a'}/m_{b'}$
over $S$, then $x$ is visible from no world.
\end{lemma}

\begin{proof}
$g\coloneqq x_be_a-x_ae_b\in F^S_x$ is nonzero and has $g_l\ge0$ for every $l\ne b$, and
similarly $g'\coloneqq x_{b'}e_{a'}-x_{a'}e_{b'}$ for $l\ne b'$. Every world $l$ differs
from $b$ or from $b'$, so by Lemma~\ref{lem:V1} $x$ is not visible from $l$; and
$x\ne e_l$.
\end{proof}

\begin{theorem}\label{thm:F}
Let $n\ge4$, $0<\alpha_j<\beta_j$ for $j\in\Omega$, $K_0\coloneqq\nu(\prod_j[\alpha_j,
\beta_j])$, and $\ub\coloneqq\one/n$. Assume (L) with $\pi^\ast=\ub$ and $L\subseteq\relint
\Delta$, and suppose that for some world $a$,
\begin{itemize}
  \item[(F)] $q_a>q_b$ for every $q\in L\setminus\{\ub\}$ and every $b\ne a$.
\end{itemize}
Put $K\coloneqq L\star K_0$. Then:
\begin{enumerate}[label=(\alph*)]
  \item $K$ is covered by the composed ranking on $\M_{D_{L'\star K_0}}$, whose chart is
  $c_j(h)=\min(\alpha_jh_j,\beta_jh_j)$ (Proposition~\ref{prop:box});
  \item for every corner $x$ of $\prod_j[\alpha_j,\beta_j]$ with $x_a=\alpha_a$ and at least
  two coordinates $x_j=\beta_j$, the point $\nu(x)$ is an extreme point of $K$ that is
  visible from no world; there are $2^{n-1}-n\ge4$ such points;
  \item if some facet normal of $L$ has no zero coordinate, then so does some facet normal
  of $K$, and $K$ is not orthant-planar.
\end{enumerate}
By (b) and Theorem~\ref{thm:reach} below, no veto--Brier ranking on any domain $D_{K''}$
has $D_K$ as an undominated model with (LF).
\end{theorem}

\begin{proof}
(a) is Theorem~\ref{thm:comp}.

(b) Let $x$ be such a corner, $H\coloneqq\{j:x_j=\beta_j\}\subseteq\Omega\setminus\{a\}$,
and $w\coloneqq\nu(x)$; $w\in K_0\subseteq K$ because $\ub\in L$. For $j\in H$ and a
generator $\nu(q\odot m)$ of $K$, $m=\nu(y)$ with $y$ in the box: $y_j/y_a\le\beta_j/\alpha_a
=w_j/w_a$, so $w_jm_a\ge w_am_j$; and $q_a\ge q_j$ by (F). Hence
$w_jq_am_a\ge w_jq_jm_a\ge q_jw_am_j$, so $w$ minimises $m_a/m_j$ over $K$. As $|H|\ge2$,
Lemma~\ref{lem:twopair} shows that $w$ is visible from no world. For extremality, let
$E\coloneqq\{m\in K:m_aw_j=m_jw_a\text{ for all }j\in H\}$, an intersection of exposed
faces, hence a face of $K$, equal to the convex hull of the generators it contains.
Equality throughout the chain above forces $q_a=q_j$, hence $q=\ub$ by (F), and
$y_j/y_a=\beta_j/\alpha_a$, that is $y_a=\alpha_a$ and $y_j=\beta_j$ for $j\in H$. So
$E=\nu(G_x)$ with $G_x\coloneqq\{y\in\prod_j[\alpha_j,\beta_j]:y_a=\alpha_a,\ y_j=\beta_j\
(j\in H)\}$, a face of the box having $x$ as a vertex. On the affine set $\{y_a=\alpha_a\}$
the map $\nu$ is injective, and by~\eqref{eq:proj} it maps segments to segments, so it maps
the extreme point $x$ of $G_x$ to an extreme point of $\nu(G_x)=E$. Hence $w$ is extreme in
$K$. Distinct corners with $x_a=\alpha_a$ have distinct images. The number of subsets
$H\subseteq\Omega\setminus\{a\}$ with $|H|\ge2$ is $2^{n-1}-1-(n-1)$.

(c) Let $\eta$ be a facet normal of $L$: $\eta\cdot q\ge0$ on $L$, with equality on a facet
$F_L$ of dimension $n-2$. Put $x^\ast_j\coloneqq\alpha_j$ if $\eta_j\ge0$ and
$x^\ast_j\coloneqq\beta_j$ if $\eta_j<0$, and $\eta'\coloneqq\eta\oslash x^\ast$. For
$q\in L$ and $y$ in the box, $\eta'\cdot(q\odot y)=\sum_j\eta_jq_jy_j/x^\ast_j\ge\sum_j
\eta_jq_j\ge0$, since $y_j/x^\ast_j\ge1$ where $\eta_j\ge0$ and $\le1$ where $\eta_j<0$.
So $\eta'\cdot m\ge0$ on $K$, with equality on $T_{\nu(x^\ast)}(F_L)$, which has dimension
$n-2$ because $T_{\nu(x^\ast)}$ is a homeomorphism of $\relint\Delta$ mapping segments to
segments. So $\eta'$ is a facet normal of $K$, with the zero pattern of $\eta$.
Lemma~\ref{lem:planarobs}(d) completes the proof.
\end{proof}

\begin{corollary}\label{cor:Ffamily}
Let $n\ge4$, $p_0\in\relint\Delta$ with $p_{0,a}>p_{0,b}$ for all $b\ne a$, and
\[
  0<\varepsilon<\min_{b\ne a}\frac{p_{0,a}-p_{0,b}}{1+p_{0,a}-p_{0,b}}.
\]
Put $v_j\coloneqq(1-\varepsilon)p_0+\varepsilon e_j$, $L'\coloneqq\{p_0\}$, $A_j\coloneqq
\{v_j\}$ and $L\coloneqq\conv(\{\ub,v_1,\dots,v_n\})$. Then every conclusion of
Theorem~\ref{thm:F}, including (c), holds for $L\star K_0$, for every box
$K_0=\nu(\prod_j[\alpha_j,\beta_j])$ with $\alpha\ll\beta$.
\end{corollary}

\begin{proof}
$\Sigma\coloneqq\conv\{v_j\}=(1-\varepsilon)p_0+\varepsilon\Delta$ is a full-dimensional
simplex in $\relint\Delta$ containing $p_0$ in its relative interior, so $L$ is
full-dimensional, $L\subseteq\relint\Delta$, $p_0\in\relint L$, and $A_j\subseteq[p_0,
e_j]$: (L) holds with $\pi^\ast=\ub$. For each $j$, $v_{j,a}-v_{j,b}\ge(1-\varepsilon)
(p_{0,a}-p_{0,b})-\varepsilon>0$ by the choice of $\varepsilon$, so $q_a>q_b$ on $\Sigma$;
a point of $L$ is $(1-\lambda)s+\lambda\ub$ with $s\in\Sigma$, and $\ub$ has equal
coordinates, so (F) holds. Finally, $\sum_k(1-\varepsilon)p_{0,k}<1=\sum_k\ub_k$, so
$\ub_k>(1-\varepsilon)p_{0,k}$ for some $k$. The facet $\Sigma\cap\{m_k=(1-\varepsilon)
p_{0,k}\}$ of $\Sigma$ is then a facet of $L$, since both $\Sigma$ and $\ub$ satisfy
$m_k\ge(1-\varepsilon)p_{0,k}$. Its normal $e_k-(1-\varepsilon)p_{0,k}\one$ has no zero
coordinate.
\end{proof}

\begin{example}\label{ex:n4}
Let $n=4$, $K_0\coloneqq\nu([1,2]^4)=\{m:m_j\le2m_i\text{ for all }i,j\}$, $p_0\coloneqq
\frac1{10}(4,3,2,1)$, $a=1$ and $\varepsilon=\frac1{20}<\frac1{11}$. So $L=\conv\{\ub,
v_1,\dots,v_4\}$ with
\[
  v_1=\bigl(\tfrac{43}{100},\tfrac{57}{200},\tfrac{19}{100},\tfrac{19}{200}\bigr),\
  v_2=\bigl(\tfrac{19}{50},\tfrac{67}{200},\tfrac{19}{100},\tfrac{19}{200}\bigr),\
  v_3=\bigl(\tfrac{19}{50},\tfrac{57}{200},\tfrac{6}{25},\tfrac{19}{200}\bigr),\
  v_4=\bigl(\tfrac{19}{50},\tfrac{57}{200},\tfrac{19}{100},\tfrac{29}{200}\bigr).
\]
At world $i$ the composed ranking compares, lexicographically, on measurable
$X\subseteq D_{K'}$ with $K'=\nu(p_0\odot K_0)=\{m:m_j/m_i\le2p_{0,j}/p_{0,i}\}$,
\[
  \Bigl(\int_X\max\bigl(0,-v_i\cdot c(h)\bigr)\,d\mu(h),\ -\int_Xc_i(h)\,d\mu(h)\Bigr),
  \qquad c(h)=h^+-2h^- .
\]
It covers $K=L\star K_0$. An exact computation of the hull of the $70$ generators
$\nu(q\odot m)$, $q\in\ext L$, $m\in\ext K_0$, shows that $K$ has $32$ vertices, $63$
edges and $33$ facets. Exactly five vertices are visible from no world: the normalised box
corners
\[
  \nu(1,1,2,2),\ \nu(1,2,1,2),\ \nu(1,2,2,1),\ \nu(1,2,2,2),\ \nu(2,1,2,2);
\]
the first four are those of Theorem~\ref{thm:F}(b), and all five are tangential from two
or three worlds. Six facet normals have no zero coordinate, among them
$(-19,-19,-19,362)$ and $(134,-57,-57,94)$; so $K$ is not orthant-planar. Thus $K$ is
covered, lies outside the visibility class of Section~\ref{sec:vb}, and lies outside the
orthant-planar class. Its non-visible vertices are tangential; we do not claim that $K$
lies outside every class of rankings with a veto level followed by a Brier-type level.
\end{example}

\subsection{Exempt points on the boundary}

\begin{proposition}\label{prop:facecount}
Let $K''\subseteq\relint\Delta$ be nonempty, compact and convex, $B$ a veto profile over
$K''$, $(W_0,\beta,\alpha)$ an admissible triple, and $S\subsetneq\Omega$ with $|S|\ge2$.
Then $K(W_0,\beta,\alpha)\cap\Delta_S$ has at most one extreme point whose support is
exactly $S$.
\end{proposition}

\begin{proof}
Put $K^\ast\coloneqq K(W_0,\beta,\alpha)=\conv T$ with $T\coloneqq K''\cup\bigcup_{\omega
\in W_0}\hat A_\omega\cup\hat W$, a compact set. Since $m\mapsto\sum_{k\notin S}m_k$ is
non-negative on $K^\ast$, $G\coloneqq K^\ast\cap\Delta_S$ is a face, $G=\conv(T\cap\Delta
_S)$, and $\ext G\subseteq T\cap\Delta_S$. Points of $K''$ have full support. By
Lemma~\ref{lem:veto}(d), points of $\hat A_\omega$ have full support or equal $e_\omega$.
Every $v\in B_\omega\subseteq\cone(K''\cup\{e_\omega\})$ is $0$, has full support, or is a
positive multiple of $e_\omega$. A point $w=\beta+\sum_{\omega\notin W_0}\alpha_\omega
v_\omega$ of $W$ has $\supp w\subseteq S$ iff $\supp\beta\subseteq S$ and, whenever
$\alpha_\omega>0$, either $v_\omega=0$ or $v_\omega=te_\omega$ with $\omega\in S$. With
$T_\omega\coloneqq\max\{t\ge0:te_\omega\in B_\omega\}$, the set of such $w$ is the box
$\{\beta+\sum_{\omega\in S'}t_\omega\alpha_\omega e_\omega:t_\omega\in[0,T_\omega]\}$,
$S'\coloneqq(S\setminus W_0)\cap\{\alpha_\omega>0\}$. The image under $\nu$ of a convex
combination of corners is a convex combination of the images of the corners, so an
extreme point of $G$ with support exactly $S$ is $\nu$ of a corner $\beta+\sum_{\omega\in
M}\alpha_\omega T_\omega e_\omega$, $M\subseteq S'$, with support exactly $S$. As
$\supp\beta\subseteq W_0$ is disjoint from $S'$, this support is $\supp\beta\cup\{\omega\in
M:T_\omega>0\}$, so the corner is determined: it is the one with $\{\omega\in M:T_\omega>0
\}=S\setminus\supp\beta$.
\end{proof}

\begin{theorem}\label{thm:BE}
Let $n\ge3$. Assume (L) with $\pi^\ast$ in the relative interior of a face $\Delta_S$ of
$\Delta$, where $S\subsetneq\Omega$ and $|S|\ge2$, and suppose $e_i\notin L$ for every
$i\in S$. Let $K_0\subseteq\relint\Delta$ be a full-dimensional orthant-planar credal set
and $K\coloneqq L\star K_0$. Then:
\begin{enumerate}[label=(\alph*)]
  \item $K$ is covered by the composed ranking on $\M_{D_{L'\star K_0}}$;
  \item $K\cap\Delta_S=\nu(\pi^\ast\odot K_0)$, which has at least two extreme points, all
  with support exactly $S$;
  \item no veto--Brier ranking $\VB(K'',B)$ with $K''\subseteq\relint\Delta$ on
  $\M_{D_{K''}}$ has $D_K$ as an undominated model;
  \item if $n=3$ and $K$ is a polygon (which holds whenever $L$ is a polygon, since $K_0$
  is then a polytope by Lemma~\ref{lem:planarobs}(a) and Lemma~\ref{lem:reach} applies),
  some edge line of $K$ passes through no vertex of $\Delta$.
\end{enumerate}
\end{theorem}

\begin{proof}
(a) is Theorem~\ref{thm:comp}. (b) $L\cap\Delta_S$ is a face of $L$, the convex hull of the
generators of $L$ in $\Delta_S$. Points of $L'\subseteq\relint L\subseteq\relint\Delta$ have
full support; points of $A_i$ are $e_i$ or have full support; $e_i\in\Delta_S$ only for
$i\in S$, and those are excluded. So $L\cap\Delta_S=\{\pi^\ast\}$. Likewise
$K\cap\Delta_S$ is the convex hull of the generators $\nu(q\odot m)$ in $\Delta_S$; as $m$
has full support, these are those with $q\in L\cap\Delta_S$, so $K\cap\Delta_S=\nu(\pi^\ast
\odot K_0)$, a compact convex set whose points have support exactly $S$. It is not a point:
for $i\ne j$ in $S$ the ratio $m_i/m_j$ is not constant on the full-dimensional $K_0$, and
$\nu(\pi^\ast\odot m)$ has ratio $(\pi^\ast_i/\pi^\ast_j)(m_i/m_j)$. So it has at least two
extreme points, which are extreme in $K$ since $K\cap\Delta_S$ is a face.
(c) If $D_K$ were undominated, then $D_K=D_{K(W_0,\beta,\alpha)}$ a.e.\ by
Theorem~\ref{thm:family}, so $K=K(W_0,\beta,\alpha)$ by Lemma~\ref{lem:open}, contradicting
(b) and Proposition~\ref{prop:facecount}.
(d) Let $S=\{a,b\}$ and let $c$ be the third world. By (b), $K\cap\Delta_S=\nu(\pi^\ast
\odot K_0)$ is a segment $[x^-,x^+]$ on the line $\{m_c=0\}$; it is an edge of $K$, and
$x^\pm$ are vertices of $K$. Let $\theta^-<\theta^+$ be the least and greatest values of
$m_a/m_b$ on $K_0$, attained at $m^\pm$, so that the ratio $m_a/m_b$ takes the values
$\rho^\pm\coloneqq(\pi^\ast_a/\pi^\ast_b)\theta^\pm$ at $x^\pm$. Suppose every edge line of
$K$ is cevian. The other edge of $K$ at $x^+$ lies on a line through $x^+$ and some $e_l$.
For $l\in S$ this line is $\{m_c=0\}$, which carries only the edge $[x^-,x^+]$; so $l=c$,
and the line is $\{m_a=\rho^+m_b\}$. It supports $K$, and $x^-$ lies on the side
$\{m_a<\rho^+m_b\}$, so $m_a\le\rho^+m_b$ on $K$. Likewise $m_a\ge\rho^-m_b$ on $K$. Since
$L$ is full-dimensional, some $q\in\relint L$ has $q_a/q_b\ne\pi^\ast_a/\pi^\ast_b$. If
$q_a/q_b>\pi^\ast_a/\pi^\ast_b$, the point $\nu(q\odot m^+)\in K$ has ratio
$(q_a/q_b)\theta^+>\rho^+$; otherwise $\nu(q\odot m^-)\in K$ has ratio below $\rho^-$. This
is a contradiction.
\end{proof}

\begin{example}\label{ex:pentagon}
Let $n=3$, let $K_0$ be the cyclic triangle of Example~\ref{ex:cyclic}, and let
\[
  L\coloneqq\conv\Bigl\{\tfrac15(1,2,2),\ \bigl(\tfrac12,\tfrac12,0\bigr),\
  \tfrac15(1,3,1)\Bigr\},\qquad \pi^\ast\coloneqq\bigl(\tfrac12,\tfrac12,0\bigr).
\]
The facet normals of $L$ are $(-2,2,-1)$, $(1,-1,2)$ and $(4,-1,-1)$. By
Lemma~\ref{lem:poly}, $\frac15(1,2,2)$ is visible from world $3$ and $\frac15(1,3,1)$ from
world $2$. So, by Corollary~\ref{cor:exemptvis} and its proof, (L) holds with
$L'=L_\delta$ for small $\delta$ and $A_j=L\cap\conv(L_\delta\cup\{e_j\})$. Neither $e_1$
nor $e_2$ lies in $L$. An exact computation gives $K=L\star K_0$ as the pentagon with
vertices
\[
  \tfrac1{15}(1,12,2),\quad\tfrac1{13}(1,8,4),\quad\tfrac16(1,1,4),\quad
  \tfrac13(2,1,0),\quad\tfrac15(1,4,0),
\]
and facet normals $(-4,8,-1)$, $(0,0,1)$, $(4,-1,4)$, $(4,0,-1)$ and $(16,-1,-2)$. Its
edge joining $\frac13(2,1,0)$ and $\frac15(1,4,0)$ is $\nu(\pi^\ast\odot K_0)$ and lies on
the boundary of $\Delta$. By Theorem~\ref{thm:BE}, $K$ is covered and is not the
undominated target of any veto--Brier ranking on a domain $D_{K''}$ with
$K''\subseteq\relint\Delta$. Three of its edge lines pass through no vertex of $\Delta$, in
accordance with Theorem~\ref{thm:BE}(d).
\end{example}

\section{Scoped exclusions}\label{sec:excl}

Each result in this section concerns a stated class of rankings, or a stated route to
coverage, and nothing more. Every class considered consists of complete and transitive
rankings, or (in Section~\ref{sec:excl-other}, items~(2) and~\ref{it:unanimity}) of
transitive unanimity rankings. None of these results shows that the targets concerned cannot be
covered in the sense of Definition~\ref{def:cover} by rankings outside the class; that
question is open (Section~\ref{sec:open}).

\subsection{Veto--Brier rankings}\label{sec:excl-vb}

\begin{lemma}\label{lem:sublinear}
Let $K''\subseteq\Delta$ be nonempty, compact and convex, and let $\varphi:\R^n\to\R$ be
non-negative, convex and positively homogeneous, with $\varphi(h)=0$ whenever $h\in
D_{K''}$ and $h_i\ge0$. Then $\varphi(h)=\max_{v\in B}(-v\cdot h)$ for a compact convex $B$
with $0\in B\subseteq\cone(K''\cup\{e_i\})$.
\end{lemma}

\begin{proof}
A finite convex positively homogeneous function on $\R^n$ is continuous, and is the
support function of a nonempty closed convex set $\tilde B$ \cite[Thm.~13.2]{Rockafellar}, % citation number to be verified
which is bounded because the function is finite: $\varphi(h)=\max_{y\in\tilde B}y\cdot h$.
Put $B\coloneqq-\tilde B$, compact and convex. If $0\notin B$, separation
gives $h$ with $-v\cdot h<0$ for all $v\in B$, so $\varphi(h)<0$; hence $0\in B$. Let
$V\coloneqq D_{K''}\cap\{h_i\ge0\}$. The hypothesis says $v\cdot h\ge0$ for $v\in B$,
$h\in V$, that is, $B\subseteq V^\ast$, the dual cone. Now $\cl V=\{h:k\cdot h\ge0\text{
for }k\in K'',\ h_i\ge0\}$ (if $h$ is in the right side, $h+t\one\in V$ for $t>0$), which
is the dual cone of $Q\coloneqq\cone(K''\cup\{e_i\})$. $Q=\R_+\conv(K''\cup\{e_i\})$ is
closed, as $\conv(K''\cup\{e_i\})$ is compact and avoids $0$. By the bipolar theorem
$V^\ast=(\cl V)^\ast=Q^{\ast\ast}=Q$.
\end{proof}

So the rankings $\VB(K'',B)$ are exactly the two-level rankings on $\M_{D_{K''}}$ whose
first level is a non-negative sublinear integral veto that vanishes on correct
acceptances, and whose second level is the error $b_i$. By Remark~\ref{rem:convexU}, the
domains $D_{K''}$ are exactly the open convex one-sided cones containing the positive
orthant. We call a credal set a \emph{veto--Brier target} if it is covered by some
$\VB(K'',B)$ on some $\M_{D_{K''}}$.

\begin{remark}\label{rem:convexU}
Let $U\subseteq\R^n$ be an open convex cone containing every gamble with positive
coordinates, with $U\cap(-U)=\emptyset$, and let $C\coloneqq\cl U$. Then $C\neq\R^n$, so
its dual cone $C^\ast$ is nonzero and contained in $\R^n_+$; put $K''\coloneqq C^\ast\cap
\Delta$, a nonempty compact convex set. For a closed convex cone with nonempty interior,
$\inter C=\{h:y\cdot h>0\text{ for all }y\in C^\ast\setminus\{0\}\}$: if $h\in\inter C$ and
$y\in C^\ast\setminus\{0\}$, then $h-\epsilon y\in C$ for small $\epsilon>0$, so $y\cdot
h\ge\epsilon|y|^2>0$; conversely, if the minimum $c$ of $y\cdot h$ over unit vectors $y\in
C^\ast$ is positive, the ball of radius $c$ about $h$ lies in $C^{\ast\ast}=C$. Since $U$
is open and convex, $U=\inter C=D_{K''}$.
\end{remark}

\begin{theorem}\label{thm:reach}
Let $K\subseteq\relint\Delta$ be a full-dimensional credal set. The following are
equivalent.
\begin{enumerate}[label=(\alph*)]
  \item For some nonempty compact convex $K''\subseteq\Delta$ and veto profile $B$ over
  $K''$, $D_K$ is an undominated model of $\VB(K'',B)$ on $\M_{D_{K''}}$ and (LF) holds at
  $D_K$.
  \item $K$ is a veto--Brier target.
  \item For some $p\in K$, the set $\ext K\setminus\{p\}$ is uniformly visible.
\end{enumerate}
\end{theorem}

\begin{proof}
(b)$\Rightarrow$(a) is immediate, and (c)$\Rightarrow$(b) is
Corollary~\ref{cor:exemptvis}.

(a)$\Rightarrow$(c). Let $U=D_{K''}$. Since $D_K\in\M_U$, Lemma~\ref{lem:LF}(b) gives
$K''\subseteq\relint K$. By Theorem~\ref{thm:family}, $D_K=D_{K(W_0,\beta,\alpha)}$ a.e.\
for an admissible triple, so $K=K(W_0,\beta,\alpha)$ by Lemma~\ref{lem:open}. Taking every
$v_\omega=0$ in the definition of $W$ shows $\nu(\beta)\in\hat W\subseteq K\subseteq
\relint\Delta$. So $\supp\beta=\Omega$, hence $W_0=\Omega$, there are no
$\alpha$-weights, and $\hat W=\{\nu(\beta)\}$. Thus
\[
  K=\conv\Bigl(K''\cup\bigcup_{\omega\in\Omega}\hat A_\omega\cup\{\nu(\beta)\}\Bigr).
\]
Then $\ext K\subseteq K''\cup\bigcup_\omega\hat A_\omega\cup\{\nu(\beta)\}$, and
$K''\cap\ext K=\emptyset$. By Lemma~\ref{lem:visbasic}(b), $K''\subseteq K_\delta$ for some
$\delta>0$, and $\hat A_\omega\subseteq\conv(K''\cup\{e_\omega\})\subseteq\conv(K_\delta
\cup\{e_\omega\})$ by Lemma~\ref{lem:veto}(d). So $\ext K\setminus\{\nu(\beta)\}$ is
uniformly visible.
\end{proof}

\begin{corollary}\label{cor:polyreach}
Let $K\subseteq\relint\Delta$ be a full-dimensional polytope. $K$ is a veto--Brier target
if and only if at most one vertex of $K$ is not visible. Tangential vertices count as not
visible.
\end{corollary}

\begin{proof}
For a finite set of points of $K$, uniform visibility is equivalent to visibility of each
point (Lemma~\ref{lem:visbasic}(a), and the proof of Corollary~\ref{cor:poly}); and no
vertex of $\Delta$ lies in $K$. Apply Theorem~\ref{thm:reach}.
\end{proof}

\begin{corollary}\label{cor:out}
Let $n\ge3$. None of the following credal sets is a veto--Brier target.
\begin{enumerate}[label=(\alph*)]
  \item \emph{Pari-mutuel models} $K=\{m\in\Delta:m_i\le(1+\delta)p_i\}$ with
  $p\in\relint\Delta$, $\delta>0$ and $K\subseteq\relint\Delta$; in particular the reversed
  triangle $R$. No vertex is visible.
  \item \emph{Density-ratio classes} $K=\{m\in\Delta:m_i/p_i\le k\,m_j/p_j\text{ for all
  }i,j\}$ with $p\in\relint\Delta$ and $k>1$. At least three vertices are not visible.
  \item Full-dimensional probability-interval models $K\subseteq\relint\Delta$ with at
  least two vertices having two or more coordinates at their upper bounds.
\end{enumerate}
\end{corollary}

\begin{proof}
In each case $K$ is a full-dimensional polytope in $\relint\Delta$; we apply
Corollary~\ref{cor:polyreach} and Lemma~\ref{lem:poly}. We use that a point $v$ of
$\{m:\one\cdot m=1,\ g_f\cdot m\ge0\}$ is a vertex iff $\one$ and the active $g_f$ span
$\R^n$, equivalently iff $v$ is the only point of $\aff\Delta$ at which all active
constraints are equalities \cite[Sec.~8.5]{Schrijver}. % citation number to be verified

(a) Put $u_i\coloneqq(1+\delta)p_i$. $p$ satisfies every constraint strictly, so $K$ is
full-dimensional. At a vertex $v$ the constraints $m_k\ge0$ are inactive, and with $H$ the
set of $i$ with $v_i=u_i$, the span of $\one$ and the normals $u_i\one-e_i$ ($i\in H$)
is the span of $\one$ and $e_i$ ($i\in H$); so $|H|\ge n-1\ge2$. Each $u_i\one-e_i$ with
$i\in H$ is negative exactly at $i$ (as $u_i=v_i<1$), so by Lemma~\ref{lem:poly} $v$ is
not visible. A full-dimensional polytope in $\aff\Delta$ has at least $n\ge3$ vertices.

(b) If $m_j=0$ for some $m\in K$ then $m=0$; so $K\subseteq\relint\Delta$, and $p$
satisfies all constraints strictly. Write the constraints as $m\cdot g^{ij}\ge0$ with
$g^{ij}\coloneqq(k/p_j)e_j-(1/p_i)e_i$, negative exactly at $i$. For $\emptyset\neq H
\subsetneq\Omega$ let $w^H\coloneqq\nu\bigl(\sum_{i\in H}kp_ie_i+\sum_{i\notin H}p_ie_i
\bigr)\in K$. The active constraints at $w^H$ are those with $i\in H$, $j\notin H$. If all
of them hold with equality at $m\in\aff\Delta$, the ratios $m_i/p_i$ take a common value
$c$ on $H$ and $c/k$ off $H$, so $m=w^H$; hence $w^H$ is a vertex. If $|H|\ge2$, two
active normals are negative at different coordinates, so $w^H$ is not visible. Distinct
$H$ give distinct vertices (the coordinates of maximal ratio form $H$), and there are
$2^n-n-2\ge3$ sets $H$ with $2\le|H|\le n-1$.

(c) This is Proposition~\ref{prop:pri}.
\end{proof}

The density-ratio classes of (b) are covered by orthant-diagonal rankings
(Proposition~\ref{prop:box}). Coverage of the classes (a) and (c) is not settled.

\begin{proposition}\label{prop:contamzero}
Let $n\ge3$, $0<\varepsilon<1$, and $p\in\Delta\setminus\relint\Delta$. Then
$K\coloneqq(1-\varepsilon)p+\varepsilon\Delta$ is not a veto--Brier target. More precisely,
there are no $K''$ and $B$ as in Theorem~\ref{thm:reach}(a).
\end{proposition}

\begin{proof}
$K$ is a full-dimensional simplex with vertices $v_i=(1-\varepsilon)p+\varepsilon e_i$.
Suppose $K''$ and $B$ are as in Theorem~\ref{thm:reach}(a). As in that proof, $K''\subseteq
\relint K\subseteq\relint\Delta$ and $K=K(W_0,\beta,\alpha)$ for an admissible triple. Let
$S\coloneqq\supp p$, a proper subset of $\Omega$.

\emph{Case $|S|\ge2$.} The face $K\cap\Delta_S=\conv\{v_i:i\in S\}$ has $|S|\ge2$ extreme
points, each with support exactly $S$. This contradicts Proposition~\ref{prop:facecount}.

\emph{Case $|S|=1$}, say $p=e_1$. Write $T\coloneqq K''\cup\bigcup_{\omega\in W_0}\hat
A_\omega\cup\hat W$, so $K=\conv T$. Points of $K''$ have full support; points of $\hat
A_\omega$ have full support or equal $e_\omega$. For $i\ge2$, $v_i=(1-\varepsilon)e_1+
\varepsilon e_i$ is extreme with support $\{1,i\}$, so it lies in $\hat W$. Write
$v_i=\nu(w^i)$. Since $\supp w^i=\{1,i\}$ for $i=2,3$, $\supp\beta\subseteq\{1\}$, so
$\beta=\beta_1e_1$ with $\beta_1>0$ and $1\in W_0$. Each $w^i$ is then $\beta_1e_1+\alpha_i
t_ie_i$ with $i\notin W_0$, $t_ie_i\in B_i$, all other terms zero; and $\nu(w^i)=v_i$
forces $\alpha_it_i=a\beta_1$ with $a\coloneqq\varepsilon/(1-\varepsilon)$. Choosing
$v_2=t_2e_2$, $v_3=t_3e_3$ and all other $v_\omega=0$ gives $w\coloneqq\beta_1(e_1+ae_2+
ae_3)\in W$. But the first coordinate of $\nu(w)$ is $1/(1+2a)=(1-\varepsilon)/
(1+\varepsilon)<1-\varepsilon\le\min_{m\in K}m_1$, so $\nu(w)\notin K$, contradicting
$\hat W\subseteq K$.
\end{proof}

\subsection{Vetoes followed by a Brier-type error}\label{sec:excl-veto}

The following class contains every veto--Brier ranking. Fix a one-sided domain $U$ and
$K\subseteq\relint\Delta$ with $D_K\subseteq U$ and (LF) at $D_K$.

\begin{definition}\label{def:VB}
A ranking on $\M_U$ is in the class $\VBB$ if at each world $k$ it compares models
lexicographically by finitely many integrated levels
\[
  \int_X\varphi_{k,1}\,d\mu,\ \dots,\ \int_X\varphi_{k,r_k}\,d\mu,\ \int_X\ell_k\,d\mu,\
  \text{followed by any finite number of further integrated levels},
\]
where the \emph{vetoes} $\varphi_{k,t}\ge0$ are integrable ($r_k=0$ allowed), and with
$Z_k\coloneqq U\cap\bigcap_t\{\varphi_{k,t}=0\}$ and $N_{\mathrm v}\coloneqq\bigcap_kZ_k$,
the \emph{error level} $\ell_k$ is integrable, nonzero a.e.\ on $Z_k$, and equal on
$N_{\mathrm v}$ to $-w(h)h_k$ for a measurable $w>0$ that does not depend on $k$. The
ranking is assumed to satisfy (G).
\end{definition}

Rankings in $\VBB$ are complete and transitive and satisfy (IB). The vetoes need not be
convex, homogeneous or continuous, and the target need not be exempt from them in
advance. Say that $x\in K$ is \emph{invisible from $j$} if $F_x$ meets $\{h_j>0\}$, and
\emph{invisible} if it is invisible from every world. Invisible points are not visible;
tangential points are not invisible.

\begin{lemma}\label{lem:V1a}
Let a ranking in $\VBB$ be given. For every $k$: every $\varphi_{k,t}=0$ a.e.\ on $U\cap
\{h_k>0\}$, so $Z_k\supseteq U\cap\{h_k>0\}$ a.e.; and on $Z_k$, $\ell_k$ has the sign of
$-h_k$ a.e.
\end{lemma}

\begin{proof}
Suppose some veto at world $k$ is not a.e.\ zero on $U\cap\{h_k>0\}$, and let $t$ be the
first such level. Then $S\coloneqq U\cap\{h_k>0\}\cap\{\varphi_{k,t}>0\}$ is not null, and
the earlier vetoes vanish a.e.\ on $S$. Comparing $\emptyset$ with $S$, the first $t-1$
levels tie at $0$ and level $t$ favours $\emptyset$; so $\emptyset\prec_kS$. But (G) gives
$S\prec_k\emptyset$. Next suppose $S\subseteq Z_k\cap\{h_k<0\}$ is not null with
$\ell_k\le0$, hence $\ell_k<0$ a.e., on $S$. Then all vetoes of $S$ at $k$ vanish and
$\int_S\ell_k<0$, so $S\prec_k\emptyset$, while (G) gives $\emptyset\prec_kS$. The case
$\{h_k>0\}$ is symmetric.
\end{proof}

\begin{lemma}\label{lem:V2}
Let a ranking in $\VBB$ be given with $D_K$ undominated. Then $D_K\subseteq Z_k$ a.e.\ for
every $k$.
\end{lemma}

\begin{proof}
Let $W_0$ be the set of worlds at which $D_K$ pays no veto. For $k\notin W_0$ let
$\kappa_k$ be the first veto level that $D_K$ pays. Let $\mathcal Z$ be the class of
models contained a.e.\ in the set where all vetoes of the worlds in $W_0$ vanish and the
vetoes at levels below $\kappa_k$ vanish for $k\notin W_0$. It contains $D_K$ and
$\emptyset$, and by Lyapunov's theorem the range of the vector
$\bigl((\int_Y\ell_k)_{k\in W_0},(\int_Y\varphi_{k,\kappa_k})_{k\notin W_0}\bigr)$ over
$Y\in\mathcal Z$ is convex. A model in $\mathcal Z$ strictly smaller than $D_K$ in every
coordinate of this vector would strictly dominate $D_K$. So by separation there are
$\beta\ge0$ on $W_0$ and $\alpha\ge0$ off $W_0$, not both zero, such that $D_K$ minimises
$\int_Yg$ over $Y\in\mathcal Z$, where $g\coloneqq\sum_{W_0}\beta_k\ell_k+\sum_{k\notin W_0}
\alpha_k\varphi_{k,\kappa_k}$. If $\beta=0$, then $\int_\emptyset g=0<\int_{D_K}g$, as
$D_K$ pays each $\varphi_{k,\kappa_k}$, a contradiction; so $\beta\neq0$. Minimality gives
$g\le0$ a.e.\ on $D_K$, hence $\sum_{W_0}\beta_k\ell_k\le0$ a.e.\ on $D_K$. If
$W_0\ne\Omega$, pick $j_0\notin W_0$. The open set $O'\coloneqq D_K\cap\{h_k<0\text{ for
all }k\in W_0\}$ is nonempty: it contains $-\sum_{k\in W_0}e_k+Me_{j_0}$ for large $M$,
because $\min_{m\in K}m_{j_0}>0$. On $O'\subseteq D_K\subseteq Z_k$ ($k\in W_0$),
Lemma~\ref{lem:V1a} gives $\ell_k>0$ a.e., so $\sum_{W_0}\beta_k\ell_k>0$ a.e.\ on $O'$,
a contradiction. Hence $W_0=\Omega$.
\end{proof}

\begin{lemma}\label{lem:V3}
Let a ranking in $\VBB$ be given. For $j\ne k$, the model $X_{kj}\coloneqq Z_k\cap
\{h_j>0\}$ is undominated.
\end{lemma}

\begin{proof}
By Lemma~\ref{lem:V1a}, $X_{kj}\subseteq Z_j\cap Z_k$ a.e. Let $Y$ weakly dominate
$X_{kj}$. Since $X_{kj}$ pays no veto at worlds $j$ and $k$, the lexicographic order forces
every veto of $Y$ at $j$ and $k$ to vanish, so $Y\subseteq Z_j\cap Z_k$ a.e., and
$\int_Y\ell_j\le\int_{X_{kj}}\ell_j$. On $Z_j$, $\ell_j<0$ a.e.\ where $h_j>0$ and
$\ell_j>0$ a.e.\ where $h_j<0$, so
$\int_Y\ell_j-\int_{X_{kj}}\ell_j=\int_{Y\setminus X_{kj}}\ell_j+\int_{X_{kj}\setminus Y}
|\ell_j|$ is positive unless $Y=X_{kj}$ a.e.
\end{proof}

\begin{lemma}\label{lem:V4}
Let a ranking in $\VBB$ be given with $D_K$ undominated and the models $X_{kj}$
coherent. Let $v\in\ext K$ be invisible. Then for every $r\in\relint F_v$ there is
$\rho>0$ with $B(r,\rho)\subseteq N_{\mathrm v}$ a.e.
\end{lemma}

\begin{proof}
Let $N_{\mathrm{LF}}$ be the open set of (LF); as in the proof of Lemma~\ref{lem:LF}(b),
$N_{\mathrm{LF}}\subseteq U$ a.e. For each $j$ let $P_j\coloneqq\relint F_v\cap
\{h_j>0\}$, a relatively open convex cone; $P_j\neq\emptyset$ because $v$ is invisible from
$j$ and $\relint F_v$ is dense in $F_v$. Every nonzero point of $F_v$ has a positive
coordinate (Lemma~\ref{lem:fulldim}(c) and $K\subseteq\relint\Delta$), so $\relint
F_v=\bigcup_jP_j$, and $\relint F_v\subseteq\partial D_K\setminus\{0\}\subseteq
N_{\mathrm{LF}}$.

\emph{Spreading.} Let $D_Q$ be an open convex cone containing $D_K\cap\{h_j>0\}$ and a
ball $B(p,\rho)$ with $p\in P_j$. Then $D_Q$ contains a ball around every $r\in P_j$. If
$r$ is a positive multiple of $p$, use homogeneity. Otherwise $q\coloneqq r+t(r-p)\in P_j$
for small $t>0$, and $r=\lambda p+(1-\lambda)q$ with $\lambda=t/(1+t)$. Any $h\in
B(r,\lambda\rho/2)$ can be written $\lambda(p+d)+(1-\lambda)(q+\zeta\one)$ with
$|d|<\rho$ and small $\zeta>0$. Since $q+\zeta\one\in D_K\cap\{h_j>0\}$ and $p+d\in
B(p,\rho)$, convexity gives $h\in D_Q$.

\emph{Induction.} Fix $k$. Join $j$ and $l$ when $P_j\cap P_l\ne\emptyset$. This graph is
connected: along a segment from a point of $P_j$ to a point of $P_l$ inside $\relint F_v$,
the relatively open sets $\{t:\gamma(t)_i>0\}$ cover $[0,1]$, and an open cover of a
connected set has a connected overlap graph. For $r\in P_k$ there is a ball around $r$
inside $N_{\mathrm{LF}}\cap\{h_k>0\}$, which lies a.e.\ in $U\cap\{h_k>0\}\subseteq Z_k$
(Lemma~\ref{lem:V1a}). Suppose balls around all points of $P_j$ lie a.e.\ in $Z_k$, and
let $l\ne k$ be joined to $j$. Pick $p\in P_j\cap P_l$ and a ball $B(p,\rho)$ lying a.e.\ in
$Z_k\cap\{h_l>0\}=X_{kl}$. By hypothesis $X_{kl}=D_Q$ a.e.\ for a credal set $Q$, and by
Lemma~\ref{lem:open} $B(p,\rho)\subseteq D_Q$. Also $D_K\subseteq Z_k$ a.e.\
(Lemma~\ref{lem:V2}), so $D_K\cap\{h_l>0\}\subseteq X_{kl}=D_Q$ a.e., hence everywhere by
Lemma~\ref{lem:open}. Spreading puts a ball around every point of $P_l$ inside $D_Q$; after
shrinking, inside $D_Q\cap N_{\mathrm{LF}}\cap\{h_l>0\}$, which lies a.e.\ in $X_{kl}
\subseteq Z_k$. By connectedness, every point of $\relint F_v$ has a ball lying a.e.\ in
$Z_k$. Intersecting over the finitely many $k$ gives the claim.
\end{proof}

\begin{theorem}\label{thm:A}
Let $U$ be a one-sided domain and $K\subseteq\relint\Delta$ with $D_K\subseteq U$ and (LF)
at $D_K$. Suppose a ranking in $\VBB$ on $\M_U$ has $D_K$ undominated and the models
$X_{kj}$ of Lemma~\ref{lem:V3} coherent (as (GC) requires). Then some $\pi\in K$ has
$\pi\cdot h=0$ for every $h\in F_v$ and every invisible extreme point $v$ of $K$. In
particular a polytope $K$ has at most one invisible vertex.
\end{theorem}

\begin{proof}
By Lemma~\ref{lem:V2}, $D_K\subseteq N_{\mathrm v}$ a.e. Models $Y\subseteq N_{\mathrm v}$
pay no veto anywhere, and their error vectors $(\int_Y-wh_k\,d\mu)_k$ form a convex set.
$D_K$ is not strictly dominated by such a model, so separation gives $\pi\in\Delta$ such
that $D_K$ minimises $\int_Y-w\,\pi\cdot h\,d\mu$ over $Y\subseteq N_{\mathrm v}$. Hence
$N_{\mathrm v}\cap\{\pi\cdot h>0\}\subseteq D_K$ a.e., and $D_K\cap\{\pi\cdot h<0\}$ is null,
hence empty (it is open). So $D_K\subseteq\{\pi\cdot h\ge0\}$, and as $D_K$ is open,
$D_K\subseteq\{\pi\cdot h>0\}=D_{\{\pi\}}$; Lemma~\ref{lem:D}(c) gives $\pi\in K$. Let $v$ be invisible and $r\in\relint F_v$, and take the ball of
Lemma~\ref{lem:V4}. If $\pi\cdot r>0$, shrink the ball so that $\pi\cdot h>0$ on it; then it
lies a.e.\ in $D_K$. But it contains the nonempty open set of points $h$ near $r-\epsilon v$
with $v\cdot h<0$, which lies outside $\cl D_K$. So $\pi\cdot r=0$ on $\relint F_v$, hence on
$F_v$. If $K$ is a polytope, it is full-dimensional by Proposition~\ref{prop:lowdim}, since $U$ is
one-sided and (LF) holds; for a vertex $v$ of a full-dimensional polytope, the only point
of $K$ orthogonal to all of $F_v$ is $v$, so $\pi=v$ for every invisible vertex $v$.
\end{proof}

\begin{corollary}\label{cor:Apmm}
Let $n\ge3$ and let $K\subseteq\relint\Delta$ be a pari-mutuel model as in
Corollary~\ref{cor:out}(a), for instance the reversed triangle. Every vertex of $K$ is
invisible, so no ranking in $\VBB$ on any one-sided domain with (LF) at $D_K$ has $D_K$
undominated and all models $X_{kj}$ coherent.
\end{corollary}

\begin{proof}
At a vertex $v$ with upper set $H$, $|H|\ge n-1\ge2$, the normal cone is generated by the
vectors $u_i\one-e_i$, $i\in H$ (Lemma~\ref{lem:poly}). For $j\in H$ choose $i\in H$,
$i\ne j$: the $j$-th coordinate of $u_i\one-e_i$ is $u_i>0$. For $j\notin H$, every
generator has $j$-th coordinate $u_i>0$. So $v$ is invisible, and $K$ has at least $n\ge3$
vertices. Apply Theorem~\ref{thm:A}.
\end{proof}

\subsection{Further exclusions, recorded without proof}\label{sec:excl-other}

The following statements have been established, for the stated classes, in unpublished
notes; we record them here without proof, with their hypotheses. In items
(3)--(5) a \emph{chart} is a positively homogeneous integrable map $c:U\to\R^3$, and the
ranking at world $i$ is the order of $\int_Xc_i\,d\mu$; for $\pi\in\R^3_+\setminus\{0\}$
put $B_\pi\coloneqq U\cap\{\pi\cdot c>0\}$. The \emph{four-condition certificate} for $K$
is: (C1) $c_i$ has the sign of $h_i$ a.e.\ on $U$; (C2) $\{\pi\cdot c=0\}$ is null for
every $\pi\in\R^3_+\setminus\{0\}$; (C3) every $B_\pi$, $\pi\in\R^3_+\setminus\{0\}$, is
coherent; (C4) $B_{\pi^\ast}=D_K$ a.e.\ for some $\pi^\ast$ with positive coordinates. By
Theorem~\ref{thm:route}, the certificate together with (LF) suffices for coverage; items
(3)--(5) concern the certificate, or stated parts of it, not coverage by such charts in
general. All items concern three worlds and the reversed triangle $R$, except where
stated; $C_R=\cl D_R$.
\begin{enumerate}[label=(\arabic*)]
  \item \emph{Convex error levels.} Consider complete lexicographic rankings that have, at
  each world, finitely many non-negative measurable veto levels, then an error level that
  is nonzero a.e.\ on its veto-free set and convex on an open convex set $G$, then any
  further levels, and that satisfy (G), on a one-sided domain $U\subseteq G$ with (LF) at
  $D_K$. No such ranking has $D_K$ undominated and all undominated models coherent when $K$
  is $R$ or any pari-mutuel triangle in $\relint\Delta_3$: for these targets the collar
  mechanism of Lemma~\ref{lem:V4} (whose proof uses only Lemmas~\ref{lem:V1a}--\ref{lem:V3},
  and these hold for every ranking of this kind with (G)) reaches all three faces of
  $C_K$, and the resulting linear system on
  the three extreme rays of $C_K$ is nonsingular. The same holds for every polytope
  $K\subseteq\relint\Delta_3$ with two or more vertices at which every world is crossing.
  The conclusion does not extend to all targets with two non-visible vertices: the cover
  in Example~\ref{ex:T} has a veto followed by a convex error level.
  \item \emph{Concave charts.} Consider rankings whose worlds are unanimities of finitely
  many views, each view consisting of a veto $\varphi_v\ge0$ vanishing on correct
  acceptances (paid by the target or not), then a level $\int_X\sigma_v$ whose integrand is
  strictly sign-correct, then any further levels, where $\varphi_v$ is convex and
  $-\sigma_v$ is concave on an open convex set containing $C_R\setminus\{0\}$. On any
  one-sided domain with (LF) at $D_R$, no such ranking covers $R$.
  \item \emph{Regular single charts.} No chart ranking with strictly sign-correct chart,
  a.e.\ unique coherent Bayes models $B_\pi$ at all $\pi$ in the open simplex, $D_R=
  B_{\pi^\ast}$ with $\pi^\ast$ interior, (LF), and the first-order regularity hypotheses
  that $\pi\mapsto\min_{m\in Q_\pi}m\cdot h$ (where $B_\pi=D_{Q_\pi}$) is Gâteaux
  differentiable at $\pi^\ast$ for $h$ in the relative interiors of the faces of $C_R$,
  with derivative not identically zero on any face, exists. The resulting corner equations
  reduce to the map $x\mapsto\frac32-\frac1x$, which has no real periodic point. On the
  full domain, in any number $n\ge3$ of worlds, a chart that is locally Lipschitz on each
  closed orthant, with $c_j/h_j$ bounded above and below by positive constants on compact
  sets away from $0$, and that meets the certificate (with null ties at every prior),
  forces $D_K$ to be orthant-planar. The argument uses only gambles that are not sure
  losses. Since $R$ is not orthant-planar (Lemma~\ref{lem:planarobs}(d)), with
  Theorem~\ref{thm:dichotomy} no such chart meets the certificate for $R$ on any domain
  that is not essentially one-sided.
  \item \emph{Continuous piecewise-linear charts.} No chart that is continuous and linear
  on each cone of a finite conic partition of a neighbourhood of $C_R\setminus\{0\}$
  satisfies (C1)--(C4) for $R$, on any open convex cone $U\supseteq C_R\setminus\{0\}$ that
  contains no nonzero $h\le0$.
  \item \emph{Direction-continuous piecewise-linear charts.} The same holds, on the same
  cones $U$, for charts whose direction $c/|c|$ agrees on each cone of a finite conic
  partition with the direction of a linear map, the scale being an arbitrary positive
  measurable function, and whose direction is continuous, provided (C1), (C4), (C3) for
  all $\pi$ in a neighbourhood of $\pi^\ast$, and (C2) at $\pi=\pi^\ast$. The corner
  equations reduce to $2z^2-3z+2=0$, which has no real root.
  \item\label{it:unanimity} \emph{Finite unanimity.} Let $U\supseteq C_R\setminus\{0\}$ be
  an open convex one-sided cone containing no nonzero $h\le0$. Let each world $j$ be the
  unanimity of finitely many views, each either a single level whose gain is concave,
  positively homogeneous and weakly sign-correct on $U$, or a convex positively
  homogeneous veto vanishing on $\{h_j>0\}$ (paid by the target or not) followed by such a
  level. Suppose the views are jointly strictly guiding: for each $j$, a.e.\ on $\{h_j\ne
  0\}$ some view at $j$ has a nonzero integrand. If every undominated model is coherent,
  then $D_R$ is dominated.
\end{enumerate}
Items (1)--(6) use the collar mechanism of Lemma~\ref{lem:V4} or a variant of it, or a
corner system at the extreme rays of $C_R$. Where a variant of Lemma~\ref{lem:V4} is used
for unanimity rankings, transitivity of each world's relation enters through the step: if
$Y\preceq_jX$ and $X$ is pointwise at least as good at $j$ as $M$, then $Y\preceq_jM$.
This holds for unanimities of lexicographic views, but it is not implied by (TC).

\subsection{Composition, veto--box and chart rankings}\label{sec:excl-comp}

\begin{proposition}\label{prop:compexcl}
Let $K\subseteq\relint\Delta$ be a full-dimensional credal set that is not orthant-planar
and none of whose extreme points is visible. Then for no orthant-planar $K_0$ and data
$L',A$ as in (L) is $D_K$ an undominated model of the composed ranking on
$\M_{D_{L'\star K_0}}$ with (LF) at $D_K$.
\end{proposition}

\begin{proof}
Suppose it is. By Theorem~\ref{thm:comp} and Lemma~\ref{lem:open}, $K=Q\star K_0$ with
$Q=K(W_0,\beta,\alpha)$. Since $\nu(\beta)\in\hat W$, $\nu(\nu(\beta)\odot m)\in K\subseteq
\relint\Delta$ for $m\in K_0$, so $\supp\beta=\Omega$; hence $W_0=\Omega$, $\hat W=\{\hat
\beta\}$ with $\hat\beta\coloneqq\nu(\beta)$, and $Q=\conv(L'\cup\bigcup_\omega\hat
A_\omega\cup\{\hat\beta\})$. By Lemma~\ref{lem:LF}(b), $L'\star K_0\subseteq\relint K$. Let
$x$ be extreme in $K$; by Lemma~\ref{lem:reach}, $x=\nu(q\odot m)$ with $q\in\ext Q
\subseteq L'\cup\bigcup_\omega\hat A_\omega\cup\{\hat\beta\}$. If $q\in L'$, then
$x\in\relint K$, which is impossible. If $q\in\hat A_\omega\subseteq\conv(L'\cup\{e_\omega
\})$, $q=(1-t)k+te_\omega$ with $k\in L'$: $t=0$ is excluded as before, $t=1$ gives $x=
e_\omega\notin\relint\Delta$, and $0<t<1$ puts $x$ on the open segment from $\nu(k\odot m)\in
\relint K$ to $e_\omega$, so $x$ is visible from $\omega$ (Lemma~\ref{lem:visbasic}(a)),
which is excluded. So $\ext K\subseteq\nu(\hat\beta\odot K_0)\subseteq K$, and since
$\nu(\hat\beta\odot K_0)$ is convex, $K=\nu(\hat\beta\odot K_0)$. For each mixed orthant
$\sigma$ and $h\in\sigma$, $h\in D_K$ iff $\hat\beta\odot h\in D_{K_0}$ iff $(d_\sigma\odot
\hat\beta)\cdot h>0$, because $\hat\beta\odot h\in\sigma$. So $K$ is orthant-planar, a
contradiction.
\end{proof}

\begin{corollary}\label{cor:comppmm}
For $n\ge3$, no pari-mutuel model $K\subseteq\relint\Delta$, in particular not the
reversed triangle, is an undominated target with (LF) of any composed ranking.
\end{corollary}

\begin{proof}
No vertex is visible (Corollary~\ref{cor:out}(a)). Each facet normal has the form
$u_i\one-e_i$ with $0<u_i<1$, so its coordinates $u_i$ and $u_i-1$ are nonzero, so $K$ is not orthant-planar by
Lemma~\ref{lem:planarobs}(d). Apply Proposition~\ref{prop:compexcl}.
\end{proof}

\begin{corollary}\label{cor:pmmvbox}
Let $n\ge3$. No pari-mutuel model $K\subseteq\relint\Delta$ is an undominated model with
(LF) of a veto--box ranking on some $\M_{D_{K'}}$, or of a chart ranking of an
orthant-planar credal set on any domain.
\end{corollary}

\begin{proof}
By Lemma~\ref{lem:planarobs}(d), as in the proof of Corollary~\ref{cor:comppmm}, $K$ is not
orthant-planar. \emph{Veto--box.} Suppose $D_K$ is an undominated model of
$\mathrm{VBox}(K',B,R)$ on $\M_{D_{K'}}$ with (LF). The first steps of the proof of
Corollary~\ref{cor:BX3}, (i)$\Rightarrow$(iii), do not use $n=3$: they give $K'\subseteq
\relint K$, $K=\conv(K'\cup\bigcup_\omega\hat A_\omega\cup E)$ with $E=\nu(\beta\odot R)$
and $\beta$ with positive coordinates, and every vertex of $K$ that is not visible lies in
$E$. Every vertex of $K$ is not visible (Corollary~\ref{cor:out}(a)), so $K=E$, a
normalised box, which is orthant-planar by Proposition~\ref{prop:box}; a contradiction.
\emph{Chart.} By Lemma~\ref{lem:charttarget}, $K$ would be orthant-planar; a contradiction.
\end{proof}

\begin{corollary}\label{cor:twocevian}
Let $n=3$ and let $K\subseteq\relint\Delta$ be a triangle exactly two of whose edges lie on
cevian lines. Then $K$ is neither a veto--Brier target nor orthant-planar, and it is not
$L\star K_0$ for any $L$ as in (L) with $\pi^\ast\in\relint\Delta$ and orthant-planar
$K_0$.
\end{corollary}

\begin{proof}
A vertex $w$ of a triangle with other vertices $w',w''$ is visible from $j$ iff
$e_j=w+s(w-w')+t(w-w'')$ for some $s,t>0$, by Lemma~\ref{lem:visbasic}(a) and the
description of $\relint K$ near $w$; that is, iff $e_j$ lies in the open angle opposite to
$K$ at $w$. These three open angles are disjoint, and none contains a point of an edge line.
The two cevian edges pass through two different vertices of $\Delta$ (two edges through
the same $e_l$ would share the line through $e_l$ and their common vertex). So at most one
vertex of $\Delta$ lies in an opposite angle, and at least two vertices of $K$ are not
visible. By Corollary~\ref{cor:polyreach}, $K$ is not a veto--Brier target; by
Theorem~\ref{thm:cevian} it is not orthant-planar; by Theorem~\ref{thm:n3} it is not such a
composed target.
\end{proof}

Such a triangle may still be covered: if exactly two of its vertices are not visible and
the edge joining them is one of the two cevian edges, it is a veto--box target
(Corollary~\ref{cor:BX3}).

\begin{proposition}\label{prop:ZC}
Let $n\ge3$, $0<\varepsilon<1$, let $p\in\Delta$ have exactly one zero coordinate, and put
$K\coloneqq(1-\varepsilon)p+\varepsilon\Delta$. For no orthant-planar $K_0\subseteq
\relint\Delta$, of any dimension, and no data $L',A$ as in (L) is $D_K$ an undominated model
of the composed ranking on $\M_{D_{L'\star K_0}}$. In particular $K$ is not of the form
$L\star K_0$ with $L$ as in (L).
\end{proposition}

\begin{proof}
Suppose $D_K$ is undominated. By Theorem~\ref{thm:comp} and Lemma~\ref{lem:open},
$K=Q\star K_0$ with $Q\coloneqq K(W_0,\beta,\alpha)$, a credal set, for an admissible
triple of $\VB(L',A)$. (If $K=L\star K_0$, then $D_K$ is undominated by
Theorem~\ref{thm:comp}.) Here $L'\subseteq\relint L\subseteq\relint\Delta$. Let $k$ be the
zero coordinate of $p$ and $S\coloneqq\Omega\setminus\{k\}$, so $2\le|S|<n$. The vertices
of $K$ are $v_i=(1-\varepsilon)p+\varepsilon e_i$, and no $e_i$ lies in $K$, since $p$ has
at least two nonzero coordinates. The face $K\cap\Delta_S=\conv\{v_i:i\in S\}$ has
$n-1\ge2$ extreme points, each with support exactly $S$. By
Proposition~\ref{prop:facecount}, applied with $K''=L'$, the face $Q\cap\Delta_S$ has at
most one extreme point with support exactly $S$.

\emph{$\dim K_0=0$.} Then $K_0=\{p_0\}$ and $K=T_{p_0}(Q)$. The bijection $T_{p_0}$ of
$\Delta$ preserves supports and maps extreme points of convex sets to extreme points, so
$Q\cap\Delta_S=T_{p_0}^{-1}(K\cap\Delta_S)$ has $n-1\ge2$ extreme points with support
exactly $S$, a contradiction.

\emph{$\dim K_0\ge1$.} The vertex $v_k$ lies in $\relint\Delta$ and the tangent cone of the
simplex $K$ at $v_k$ is $\mathcal T\coloneqq\cone\{e_i-e_k:i\ne k\}=\{y:\one\cdot y=0,\ y_i\ge0
\text{ for }i\ne k\}$. Of the directions $\pm(e_j-v_k)$, only $v_k-e_k$ lies in $\mathcal T$: for
$j\ne k$, $e_j-v_k$ is negative at each $i\notin\{j,k\}$ (as $p_i>0$) and $v_k-e_j$ is
negative at $j$; and $e_k-v_k$ is negative off $k$. By Lemma~\ref{lem:reach},
$v_k=\nu(q\odot m)$ with $q\in\ext Q$, $m\in\ext K_0$, and $q\in\relint\Delta$ since $v_k$
is. The copy $T_q(K_0)\subseteq K$ is a polytope of dimension $\dim K_0$ with vertex
$v_k=T_q(m)$, and its edges lie on cevian lines, because those of $K_0$ do
(Lemma~\ref{lem:planarobs}(c)) and $T_q$ maps lines through $e_j$ into lines through
$e_j$. The direction of each edge at $v_k$ lies in $\mathcal T$ and is $\pm(e_j-v_k)$ for some $j$,
hence a positive multiple of $v_k-e_k$. If $\dim K_0\ge2$, the vertex $v_k$ of the copy
lies on at least two edges with independent directions \cite{Ziegler}, which is impossible.
If $\dim K_0=1$, the copy is a segment on the line through $e_k$ and $v_k$, so $K_0$ lies
on a line through $e_k$, and the ratios $m_i/m_j$ with $i,j\in S$ are constant on $K_0$.
The face $K\cap\Delta_S$ is the convex hull of the generators $\nu(q'\odot m')$ with
$q'\in Q\cap\Delta_S$, and these do not depend on $m'\in K_0$. So, for any fixed
$m'\in K_0$, $K\cap\Delta_S=T_{m'}(Q\cap\Delta_S)$, and $Q\cap\Delta_S=T_{m'}^{-1}(K\cap
\Delta_S)$ has $n-1\ge2$ extreme points with support exactly $S$, a contradiction.
\end{proof}

Contaminations whose centre has two or more zero coordinates are not treated by
Proposition~\ref{prop:ZC}; all their vertices lie on $\rbd\Delta$. Every contamination
whose centre lies on $\rbd\Delta$ is excluded from the veto--Brier targets by
Proposition~\ref{prop:contamzero}.

\subsection{What the exclusions do and do not show}

Each exclusion above says that a particular construction, or a particular sufficient
route, does not reach a target. The reversed triangle is reached by none of the
constructions of this paper and lies in the scope of every exclusion above. For the other
pari-mutuel models, the exclusions that apply are Corollary~\ref{cor:out}(a),
Theorem~\ref{thm:A} with Corollary~\ref{cor:Apmm}, Corollaries~\ref{cor:comppmm}
and~\ref{cor:pmmvbox}, and, for triangles in $\relint\Delta_3$, item~(1) of
Section~\ref{sec:excl-other}. None of this shows that these targets are uncovered. A ranking
covering the reversed triangle would have to lie outside every class above; one covering
another pari-mutuel model would have to lie outside the classes whose exclusions apply to
it. It could do so, for instance, by using incomparability in an essential way, by having
levels that are not integrals of measurable functions, or by being additive with a chart
outside the regularity classes listed. In particular, nothing here says that
pari-mutuel models are impermissible as coherent models.

\section{Open problems}\label{sec:open}

\begin{enumerate}[label=(\arabic*)]
  \item \emph{Universal coverage.} Is every full-dimensional credal set covered, in the
  sense of Definition~\ref{def:cover}, by some ranking on some domain with (LF)? The
  simplest open case is the reversed triangle. An impossibility proof would have to work
  from the requirements of Section~\ref{sec:framework} alone: it must identify the rival
  models, the comparisons the requirements force (including disjunctions over worlds), the
  places where coherent undominated dominators are needed, and how (IB) transfers
  comparisons, without assuming a numerical representation, transitivity or completeness.
  \item \emph{Necessary relational lemmas.} Lemma~\ref{lem:worldbest} and
  Theorem~\ref{thm:dichotomy} use only strict guidance. Which further restrictions on
  rankings, beyond these, follow from the requirements of Definition~\ref{def:cover}
  under (TC) alone?
  \item \emph{Incomparability.} Can a ranking that leaves some comparisons undecided, while
  satisfying (G), (IB) and (TC), protect a target such as the reversed triangle against
  the collar rivals of Lemma~\ref{lem:V4} and still have only coherent undominated models?
  \item \emph{Three worlds.} The covered targets in $\relint\Delta_3$ are not just the
  veto--Brier targets together with the cevian polygons: the triangle of
  Example~\ref{ex:T} is covered and lies in neither class. Among interior triangles, the
  ones covered by the constructions of this paper, in the sense of
  Corollary~\ref{cor:trireach} (in particular, with composition restricted to targets of
  the form $L\star K_0$), are exactly those in which every edge joining two non-visible
  vertices lies on a cevian line.
  We conjecture that an interior triangle is covered, in the sense of
  Definition~\ref{def:cover}, by some ranking on some domain with (LF), if and only if
  every edge joining two of its non-visible vertices lies on a cevian line. The ``if''
  half holds by Corollary~\ref{cor:trireach}, so the content of the conjecture is its
  ``only if'' half. That half implies that the reversed triangle is not covered, so proving
  the conjecture would settle the simplest open case of problem (1), negatively. Test cases are
  triangles with exactly two non-visible vertices joined by an edge on no cevian line,
  such as $\conv\{\frac18(1,6,1),\frac1{10}(6,1,3),\frac18(6,1,1)\}$, with one cevian edge,
  and $\conv\{\frac1{11}(6,4,1),\frac1{10}(6,1,3),\frac18(6,1,1)\}$, with two.
  \item \emph{Composition in four or more worlds.} Characterise the extreme points of
  $L\star K_0$ that are not visible, and decide which targets factor as $L\star K_0$. The
  two-term criterion of Lemma~\ref{lem:twopair} is sufficient; decide whether it is
  necessary. Decide
  whether the targets of Theorem~\ref{thm:F}, whose bad vertices are tangential, are
  reached by rankings with a veto level followed by a Brier-type level.
  \item \emph{Boundary targets.} Decide coverage of $\varepsilon$-contaminations whose centre
  lies on $\rbd\Delta$, including centres with two or more zero coordinates, and decide
  whether the targets of Theorem~\ref{thm:BE} for $n\ge4$ have an edge on no cevian line.
  \item \emph{Two-sided domains.} By Theorem~\ref{thm:dichotomy}, for interior targets in
  $n\ge3$ worlds the only domains that are not essentially one-sided contain every gamble
  that is not a sure loss. Which full-dimensional interior targets
  that are not orthant-planar are covered on the full domain?
  \item \emph{Probability intervals.} Decide coverage of probability-interval models with
  two or more vertices having at least two coordinates at their upper bounds.
\end{enumerate}

\begin{thebibliography}{99}

\bibitem{Ichiishi}
T.~Ichiishi.
\newblock Super-modularity: applications to convex games and to the greedy algorithm
for {LP}.
\newblock \emph{Journal of Economic Theory}, 25:283--286, 1981.

\bibitem{Lindenstrauss}
J.~Lindenstrauss.
\newblock A short proof of {L}iapounoff's convexity theorem.
\newblock \emph{Journal of Mathematics and Mechanics}, 15:971--972, 1966.

\bibitem{Lyapunov}
A.~Lyapunov.
\newblock Sur les fonctions-vecteurs compl\`etement additives.
\newblock \emph{Izvestiya Akademii Nauk SSSR, Seriya Matematicheskaya}, 4:465--478,
1940.

\bibitem{Rockafellar}
R.~T. Rockafellar.
\newblock \emph{Convex Analysis}.
\newblock Princeton University Press, 1970.

\bibitem{Schrijver}
A.~Schrijver.
\newblock \emph{Theory of Linear and Integer Programming}.
\newblock Wiley, 1986.

\bibitem{Shafer}
G.~Shafer.
\newblock \emph{A Mathematical Theory of Evidence}.
\newblock Princeton University Press, 1976.

\bibitem{Shapley}
L.~S. Shapley.
\newblock Cores of convex games.
\newblock \emph{International Journal of Game Theory}, 1:11--26, 1971.

\bibitem{Walley}
P.~Walley.
\newblock \emph{Statistical Reasoning with Imprecise Probabilities}.
\newblock Chapman and Hall, 1991.

\bibitem{Ziegler}
G.~M. Ziegler.
\newblock \emph{Lectures on Polytopes}.
\newblock Springer, 1995.

\end{thebibliography}

\end{document}
